\chapter{Deep feedforward architectures, backpropagation, and convolutional networks}
\label{sec:deep_networks_and_cnns}

\section{The perceptron, linear separability, and the dawn of neural computing}
\label{sec:perceptron_and_linear_separability}

The quest for computational learning models began with the ambition to emulate biological neural processing. Having investigated Bayesian decision theory, parametric estimation, and linear classifiers in previous chapters, we now examine the foundational connectionist architecture: the \textbf{Perceptron}.

\subsection{The biological neuron and Rosenblatt's artificial threshold unit}
\label{sec:biological_and_artificial_neuron}

In biological nervous systems, an individual neuron consists of four primary morphological structures:
\begin{itemize}[noitemsep]
  \item \textbf{Dendrites}: Branching nerve fibers that receive incoming chemical neurotransmitter signals from adjacent upstream neurons across synaptic clefts.
  \item \textbf{Soma (Cell Body)}: The metabolic hub containing the nucleus, which integrates and accumulates incoming post-synaptic electrical potentials over space and time.
  \item \textbf{Axon}: A long, slender nerve transmission cable (often insulated by fatty myelin sheaths) that propagates an all-or-nothing action potential (voltage spike) down its length once the soma's membrane potential surpasses an excitation threshold.
  \item \textbf{Synapses}: Chemical junctions connecting axon terminal boutons to recipient dendrites, where biological synaptic plasticity modulates the efficacy of electrical transmission.
\end{itemize}

\begin{figure}[H]
\centering
\begin{tikzpicture}[scale=0.92, every node/.style={transform shape}]
  % Left Panel: Biological Neuron
  \begin{scope}[shift={(0,0)}]
    \node[font=\bfseries\sffamily\color{politoNavy}] at (3.0, 2.7) {(a) Biological Neuron Anatomy};
    
    % Dendrites branching
    \draw[thick, deepblue] (0.2, 1.8) -- (1.5, 1.2);
    \draw[thick, deepblue] (0.0, 1.0) -- (1.5, 0.8);
    \draw[thick, deepblue] (0.3, 0.0) -- (1.5, 0.4);
    \draw[thick, deepblue] (-0.2, 1.5) -- (0.8, 1.2);
    \draw[thick, deepblue] (-0.2, 0.4) -- (0.8, 0.6);
    \node[font=\scriptsize\color{deepblue}, align=center] at (0.3, -0.4) {Dendrites\\(Inputs $\mathbf{x}$)};

    % Soma (cell body)
    \draw[fill=deepblue!15, draw=deepblue, very thick] (2.2, 0.8) circle (0.85);
    \draw[fill=deepblue!40, draw=deepblue] (2.2, 0.8) circle (0.35);
    \node[font=\scriptsize\bfseries\color{politoNavy}, align=center] at (2.2, 0.8) {Soma\\($\sum$)};
    \node[font=\scriptsize\color{deepblue}, align=center] at (2.2, -0.4) {Cell Body\\(Accumulator)};

    % Axon with myelin sheath
    \draw[ultra thick, deepblue] (3.05, 0.8) -- (5.2, 0.8);
    \foreach \x in {3.3, 3.9, 4.5} {
      \draw[fill=skyBlue!40, draw=politoNavy, rounded corners=2pt] (\x, 0.60) rectangle (\x+0.45, 1.00);
    }
    \node[font=\scriptsize\color{deepblue}, align=center] at (4.1, -0.4) {Axon Cable\\(Transmission)};

    % Axon Terminals / Synapses
    \draw[thick, deepblue] (5.2, 0.8) -- (5.9, 1.3);
    \draw[thick, deepblue] (5.2, 0.8) -- (6.0, 0.8);
    \draw[thick, deepblue] (5.2, 0.8) -- (5.9, 0.3);
    \foreach \pos in {(5.9, 1.3), (6.0, 0.8), (5.9, 0.3)} {
      \fill[examred] \pos circle (2.2pt);
    }
    \node[font=\scriptsize\color{examred}, align=center] at (6.0, -0.4) {Synaptic Boutons\\(Output $\hat{y}$)};
  \end{scope}

  % Vertical Separator
  \draw[dashed, gray!40, thick] (6.9, -0.8) -- (6.9, 3.1);

  % Right Panel: Artificial Perceptron
  \begin{scope}[shift={(7.8,0)}]
    \node[font=\bfseries\sffamily\color{politoNavy}] at (3.5, 2.7) {(b) Rosenblatt's Artificial Perceptron};

    % Inputs
    \node[circle, draw=deepblue, fill=deepblue!10, thick, inner sep=2pt, font=\footnotesize] (x1) at (0, 1.8) {$x_1$};
    \node[circle, draw=deepblue, fill=deepblue!10, thick, inner sep=2pt, font=\footnotesize] (x2) at (0, 0.9) {$x_2$};
    \node[font=\footnotesize] (xdots) at (0, 0.2) {$\vdots$};
    \node[circle, draw=deepblue, fill=deepblue!10, thick, inner sep=2pt, font=\footnotesize] (xd) at (0, -0.6) {$x_d$};
    \node[circle, draw=gray!70, fill=gray!10, thick, inner sep=2pt, font=\footnotesize] (xb) at (0, -1.4) {$+1$};

    % Summation Node
    \node[circle, draw=politoNavy, fill=politoNavy!15, very thick, inner sep=5pt, font=\normalsize] (sum) at (2.8, 0.6) {$\sum$};
    \node[font=\tiny\color{politoNavy}, below=2pt of sum] {$z = \mathbf{w}^\top\mathbf{x} + b$};

    % Connections with weights
    \draw[-{Stealth}, thick] (x1) -- node[above, font=\scriptsize] {$w_1$} (sum);
    \draw[-{Stealth}, thick] (x2) -- node[above, font=\scriptsize] {$w_2$} (sum);
    \draw[-{Stealth}, thick] (xd) -- node[below, font=\scriptsize] {$w_d$} (sum);
    \draw[-{Stealth}, thick, dashed] (xb) -- node[below, font=\scriptsize] {$b$} (sum);

    % Activation Function Block
    \node[rectangle, draw=intuitiongreen, fill=intuitiongreen!15, thick, inner sep=5pt, font=\footnotesize, right=0.9cm of sum] (act) {$\operatorname{sgn}(z)$};

    \draw[-{Stealth}, thick] (sum) -- node[above, font=\scriptsize] {$z$} (act);

    % Output
    \node[right=0.9cm of act, font=\footnotesize\bfseries] (out) {$\hat{y} \in \{-1, +1\}$};
    \draw[-{Stealth}, very thick, examred] (act) -- (out);
  \end{scope}
\end{tikzpicture}
\caption{Morphological and mathematical correspondence between biological and artificial neural processing. (a) A biological neuron integrates post-synaptic potentials across dendrites into the metabolic soma, firing an action potential down the myelinated axon toward synaptic terminals. (b) Rosenblatt's artificial perceptron models this mechanism: numerical inputs $x_j$ are modulated by synaptic parameters $w_j$, accumulated by affine summation $z = \mathbf{w}^\top \mathbf{x} + b$, and thresholded by a Heaviside sign function to produce binary classification $\hat{y} \in \{-1, +1\}$.}
\label{fig:biological_vs_artificial_neuron}
\end{figure}

In 1958, Frank Rosenblatt~\cite{rosenblatt1958perceptron} formalized this biological mechanism into the \textbf{Perceptron}. Given an input vector of features $\mathbf{x} = [x_1, \dots, x_d]^\top \in \R^d$, the artificial neuron associates each input channel with a tunable real-valued weight $w_j \in \R$, incorporates a scalar bias threshold $b \in \R$, and computes a linear activation score:
\begin{equation}
z \triangleq \mathbf{w}^\top \mathbf{x} + b = \sum_{j=1}^d w_j x_j + b.
\label{eq:perceptron_preactivation}
\end{equation}
The continuous scalar score $z$ is then passed through a discontinuous Heaviside step activation function:
\begin{equation}
\hat{y} = f(z) = \operatorname{sgn}(z) = \begin{cases} +1 & \text{if } \mathbf{w}^\top \mathbf{x} + b \ge 0, \\ -1 & \text{if } \mathbf{w}^\top \mathbf{x} + b < 0. \end{cases}
\label{eq:perceptron_activation}
\end{equation}

\begin{figure}[H]
\centering
\begin{tikzpicture}[scale=1.15]
  % Axes
  \draw[-{Stealth}, gray!70, thick] (-1.0, 0) -- (5.0, 0) node[right, font=\footnotesize] {$x_1$};
  \draw[-{Stealth}, gray!70, thick] (0, -1.0) -- (0, 4.2) node[above, font=\footnotesize] {$x_2$};

  % Half-space shading
  \fill[deepblue!8, opacity=0.8] (-0.8, 3.8) -- (4.2, 0.46) -- (4.5, 3.8) -- (-0.8, 3.8) -- cycle;
  \fill[examred!8, opacity=0.8] (-0.8, 3.8) -- (4.2, 0.46) -- (4.2, -0.8) -- (-0.8, -0.8) -- cycle;

  % Hyperplane: w1*x1 + w2*x2 + b = 0
  \draw[very thick, politoNavy] (-0.6, 2.5) -- (4.4, -0.8) node[right, font=\footnotesize, align=left] {Decision Hyperplane:\\$\mathbf{w}^\top \mathbf{x} + b = 0$};

  % Normal vector w from a point on the line (1.5, 1.0)
  \coordinate (P) at (1.5, 1.0);
  \draw[-{Stealth[length=6pt]}, ultra thick, intuitiongreen] (P) -- ++(1.0, 1.5) node[above right, font=\footnotesize\bfseries] {$\mathbf{w}$ (Orthogonal Normal Vector)};

  % Perpendicular projection from Origin (0,0) to line
  \coordinate (Proj) at (0.923, 1.385);
  \draw[dashed, thick, gray!80] (0, 0) -- (Proj);
  \draw[thick, gray!80] (0.8, 1.2) -- (0.95, 1.1) -- (1.07, 1.28);
  \node[left, font=\scriptsize\bfseries, text=politoNavy] at (0.45, 0.75) {$d = \frac{|b|}{\|\mathbf{w}\|_2}$};

  % Data points: Positive Class (+1)
  \node[star, star points=5, star point ratio=2.2, draw=deepblue, fill=deepblue!30, thick, inner sep=2.5pt] at (2.2, 2.6) {};
  \node[star, star points=5, star point ratio=2.2, draw=deepblue, fill=deepblue!30, thick, inner sep=2.5pt] at (3.5, 2.1) {};
  \node[star, star points=5, star point ratio=2.2, draw=deepblue, fill=deepblue!30, thick, inner sep=2.5pt] at (1.2, 3.1) {};
  \node[star, star points=5, star point ratio=2.2, draw=deepblue, fill=deepblue!30, thick, inner sep=2.5pt] at (3.8, 3.3) {};
  \node[font=\footnotesize\bfseries, text=deepblue] at (3.2, 3.8) {Positive Half-Space: $\mathbf{w}^\top\mathbf{x} + b > 0$};

  % Data points: Negative Class (-1)
  \node[circle, draw=examred, fill=examred!20, thick, inner sep=2.5pt] at (0.5, 0.5) {};
  \node[circle, draw=examred, fill=examred!20, thick, inner sep=2.5pt] at (1.5, -0.2) {};
  \node[circle, draw=examred, fill=examred!20, thick, inner sep=2.5pt] at (2.8, -0.4) {};
  \node[circle, draw=examred, fill=examred!20, thick, inner sep=2.5pt] at (0.2, 1.2) {};
  \node[font=\footnotesize\bfseries, text=examred] at (1.5, -0.7) {Negative Half-Space: $\mathbf{w}^\top\mathbf{x} + b < 0$};

  % Origin
  \fill (0,0) circle (2pt) node[below left, font=\footnotesize] {$\mathbf{0}$};
\end{tikzpicture}
\caption{The geometry of the linear decision boundary defined by the perceptron. The parameter vector $\mathbf{w} \in \mathbb{R}^d$ is the orthogonal normal vector to the hyperplane pointing directly toward the positive decision half-space. The scalar bias $b \in \mathbb{R}$ offsets the hyperplane from the origin, establishing a signed perpendicular Euclidean distance of $\frac{|b|}{\|\mathbf{w}\|_2}$.}
\label{fig:hyperplane_geometry}
\end{figure}

As visualized in Figure~\ref{fig:hyperplane_geometry}, the equation $\mathbf{w}^\top \mathbf{x} + b = 0$ defines a $(d-1)$-dimensional \textbf{decision hyperplane} that partitions the $d$-dimensional input space $\mathcal{X} \subseteq \R^d$ into two open half-spaces. The weight vector $\mathbf{w}$ is the normal vector to the hyperplane, pointing toward the positive decision region, while $b / \|\mathbf{w}\|_2$ quantifies the signed perpendicular distance from the coordinate origin to the hyperplane.

\subsection{The perceptron learning algorithm and Novikoff's convergence theorem}
\label{sec:perceptron_learning_algorithm}

Given a supervised training dataset $\mathcal{D} = \{(\mathbf{x}_i, y_i)\}_{i=1}^N$ with labels $y_i \in \{-1, +1\}$, the perceptron learning algorithm operates in an online, mistake-driven fashion. Parameters are initialized arbitrarily (typically $\mathbf{w}^{(0)} = \mathbf{0}, b^{(0)} = 0$).

At step $t$, the algorithm inspects a training instance $(\mathbf{x}_i, y_i)$. If the current hypothesis correctly classifies the instance ($y_i (\mathbf{w}^{(t)\top} \mathbf{x}_i + b^{(t)}) > 0$), no modification is made. If the instance is misclassified ($y_i (\mathbf{w}^{(t)\top} \mathbf{x}_i + b^{(t)}) \le 0$), the parameters are updated via the additive correction:
\begin{align}
\mathbf{w}^{(t+1)} &= \mathbf{w}^{(t)} + \eta y_i \mathbf{x}_i, \label{eq:perceptron_update_w} \\
b^{(t+1)} &= b^{(t)} + \eta y_i R^2, \label{eq:perceptron_update_b}
\end{align}
where $\eta > 0$ is the learning rate (often set to $\eta = 1$ without loss of generality, as scaling both sides leaves the sign unchanged), and $R = \max_i \|\mathbf{x}_i\|_2$.

\begin{figure}[H]
\centering
\begin{minipage}{0.48\linewidth}
  \centering
  \begin{tikzpicture}[scale=1.05]
    \node[font=\footnotesize\bfseries\color{politoNavy}] at (1.5, 3.2) {(a) Parameter Space: Vector Correction};
    \draw[-{Stealth}, gray!50] (-0.5, 0) -- (3.2, 0) node[right, font=\tiny] {$w_1$};
    \draw[-{Stealth}, gray!50] (0, -0.5) -- (0, 3.0) node[above, font=\tiny] {$w_2$};
    
    % w(t)
    \coordinate (W0) at (0.8, 1.8);
    \draw[-{Stealth[length=6pt]}, very thick, politoNavy] (0, 0) -- (W0) node[above left, font=\footnotesize] {$\mathbf{w}^{(t)}$};
    
    % Misclassified sample vector y_i * x_i (positive class: +x_i)
    \coordinate (X) at (1.6, -0.4);
    \draw[-{Stealth[length=6pt]}, very thick, examred] (W0) -- ++(X) node[midway, above right, font=\footnotesize] {$+y_i \mathbf{x}_i$};
    
    % w(t+1)
    \coordinate (W1) at ($(W0) + (X)$);
    \draw[-{Stealth[length=6pt]}, very thick, intuitiongreen] (0, 0) -- (W1) node[right, font=\footnotesize\bfseries] {$\mathbf{w}^{(t+1)}$};
    
    \node[draw=politoNavy, fill=politoNavy!5, rounded corners, font=\scriptsize, align=center] at (1.5, -0.8) {$\mathbf{w}^{(t+1)} = \mathbf{w}^{(t)} + \eta y_i \mathbf{x}_i$};
  \end{tikzpicture}
\end{minipage}\hfill
\begin{minipage}{0.48\linewidth}
  \centering
  \begin{tikzpicture}[scale=1.05]
    \node[font=\footnotesize\bfseries\color{politoNavy}] at (1.5, 3.2) {(b) Input Space: Boundary Rotation};
    \draw[-{Stealth}, gray!50] (-1.0, 0) -- (3.2, 0) node[right, font=\tiny] {$x_1$};
    \draw[-{Stealth}, gray!50] (0, -1.0) -- (0, 3.0) node[above, font=\tiny] {$x_2$};

    % Old boundary w(t)^T x = 0
    \draw[thick, politoNavy, dashed] (-0.9, 0.4) -- (2.5, -1.1) node[below right, font=\tiny, align=left] {Old boundary:\\$\mathbf{w}^{(t)\top}\mathbf{x}=0$};
    
    % Misclassified point x_i (positive class y_i = +1)
    \node[star, star points=5, star point ratio=2.2, draw=examred, fill=examred!30, thick, inner sep=2.5pt, label={right:\scriptsize $\mathbf{x}_i \ (y_i = +1)$}] (xi) at (1.2, -0.8) {};

    % New boundary w(t+1)^T x = 0
    \draw[very thick, intuitiongreen] (-0.8, -0.6) -- (2.0, 2.2) node[above right, font=\tiny, align=left] {New boundary:\\$\mathbf{w}^{(t+1)\top}\mathbf{x}=0$};

    % Rotation arrow
    \draw[-{Stealth}, thick, deepdiveTeal] (1.8, -0.7) to[bend right=25] (1.3, 1.4) node[midway, right, font=\tiny] {Boundary rotates toward $\mathbf{x}_i$};
  \end{tikzpicture}
\end{minipage}
\caption{Geometric dynamics of the Perceptron Learning Algorithm update. (a) In parameter space, whenever a training sample $(\mathbf{x}_i, y_i)$ is misclassified, the algorithm performs vector addition $\mathbf{w}^{(t+1)} = \mathbf{w}^{(t)} + y_i \mathbf{x}_i$. (b) In input space, adding $y_i \mathbf{x}_i$ rotates the orthogonal decision hyperplane directly toward the misclassified sample, pulling it into the correct decision region.}
\label{fig:perceptron_learning_dynamics}
\end{figure}

The historical significance of Rosenblatt's perceptron stems from the fact that it was the first machine learning model accompanied by a rigorous mathematical convergence proof, formalized by Albert Novikoff in 1962~\cite{novikoff1962convergence}.

\begin{theorem}[Novikoff's Perceptron Convergence Theorem]
\label{thm:novikoff_convergence}
Let $\mathcal{D} = \{(\mathbf{x}_i, y_i)\}_{i=1}^N$ be a training dataset where $\|\mathbf{x}_i\|_2 \le R$ for all $i$. Assume the dataset is strictly \textbf{linearly separable}: that is, there exists an optimal unit-norm parameter vector $\mathbf{w}^*$ ($\|\mathbf{w}^*\|_2 = 1$) and an optimal margin $\gamma > 0$ such that:
\begin{equation}
y_i (\mathbf{w}^{*\top} \mathbf{x}_i) \ge \gamma, \quad \forall i \in \{1, \dots, N\}.
\label{eq:linear_separability_condition}
\end{equation}
Then, starting from $\mathbf{w}^{(0)} = \mathbf{0}$, the perceptron learning algorithm makes at most $k$ mistake updates before converging to a perfectly separating hyperplane, where:
\begin{equation}
k \le \left( \frac{R}{\gamma} \right)^2.
\label{eq:novikoff_bound}
\end{equation}
\end{theorem}

\begin{proof}
Let $\mathbf{w}^{(k)}$ denote the weight vector after $k$ mistaken updates. Assume the $k$-th update was triggered by misclassified sample $(\mathbf{x}_i, y_i)$, meaning $y_i (\mathbf{w}^{(k-1)\top} \mathbf{x}_i) \le 0$. The update rule is $\mathbf{w}^{(k)} = \mathbf{w}^{(k-1)} + y_i \mathbf{x}_i$.

We bound the inner product $\mathbf{w}^{(k)\top} \mathbf{w}^*$ from below. Expanding recursively:
\begin{equation}
\mathbf{w}^{(k)\top} \mathbf{w}^* = (\mathbf{w}^{(k-1)} + y_i \mathbf{x}_i)^\top \mathbf{w}^* = \mathbf{w}^{(k-1)\top} \mathbf{w}^* + y_i \mathbf{x}_i^\top \mathbf{w}^* \ge \mathbf{w}^{(k-1)\top} \mathbf{w}^* + \gamma.
\end{equation}
Inductively from $\mathbf{w}^{(0)} = \mathbf{0}$, after $k$ mistakes:
\begin{equation}
\mathbf{w}^{(k)\top} \mathbf{w}^* \ge k \gamma.
\label{eq:novikoff_inner_bound}
\end{equation}

Next, we bound the squared Euclidean norm $\|\mathbf{w}^{(k)}\|_2^2$ from above:
\begin{align}
\|\mathbf{w}^{(k)}\|_2^2 &= \|\mathbf{w}^{(k-1)} + y_i \mathbf{x}_i\|_2^2 = \|\mathbf{w}^{(k-1)}\|_2^2 + 2 y_i \mathbf{w}^{(k-1)\top} \mathbf{x}_i + \|\mathbf{x}_i\|_2^2 \nonumber \\
&\le \|\mathbf{w}^{(k-1)}\|_2^2 + 0 + R^2 = \|\mathbf{w}^{(k-1)}\|_2^2 + R^2,
\end{align}
where we used the mistake condition $y_i \mathbf{w}^{(k-1)\top} \mathbf{x}_i \le 0$ and the radius bound $\|\mathbf{x}_i\|_2 \le R$. Inductively from $\mathbf{w}^{(0)} = \mathbf{0}$:
\begin{equation}
\|\mathbf{w}^{(k)}\|_2^2 \le k R^2 \implies \|\mathbf{w}^{(k)}\|_2 \le \sqrt{k} R.
\label{eq:novikoff_norm_bound}
\end{equation}

By the Cauchy-Schwarz inequality, for any unit vector $\mathbf{w}^*$:
\begin{equation}
\mathbf{w}^{(k)\top} \mathbf{w}^* \le \|\mathbf{w}^{(k)}\|_2 \|\mathbf{w}^*\|_2 = \|\mathbf{w}^{(k)}\|_2.
\end{equation}
Combining Equation~\eqref{eq:novikoff_inner_bound} and Equation~\eqref{eq:novikoff_norm_bound}:
\begin{equation}
k \gamma \le \mathbf{w}^{(k)\top} \mathbf{w}^* \le \|\mathbf{w}^{(k)}\|_2 \le \sqrt{k} R \implies k \gamma \le \sqrt{k} R \implies \sqrt{k} \le \frac{R}{\gamma} \implies k \le \left( \frac{R}{\gamma} \right)^2.
\end{equation}
This completes the proof.
\end{proof}

\subsection{The fatal flaws: The XOR impasse and Minsky \& Papert's critique}
\label{sec:xor_and_minsky_critique}

Despite the mathematical elegance of Novikoff's theorem, the single-layer perceptron suffers from fundamental theoretical limitations that ultimately triggered the historical ``First AI Winter''.

\paragraph{1. The Non-Separability Catastrophe: Infinite Cycling.}
Novikoff's Perceptron Convergence Theorem (Theorem~\ref{thm:novikoff_convergence}) holds strictly under the premise that the training dataset $\mathcal{D}$ is linearly separable, meaning that a strictly positive margin exists: $\gamma > 0$. 

However, in realistic real-world datasets, sensory noise, label ambiguity, and overlapping class probability distributions inevitably render data \textbf{linearly non-separable} ($\gamma \le 0$). When applied to non-separable data, the perceptron learning algorithm fails catastrophically:
\begin{itemize}[leftmargin=*]
  \item \textbf{Destruction of the Convergence Guarantee}: If no hyperplane can simultaneously satisfy all training constraints $y_i(\mathbf{w}^\top \mathbf{x}_i) > 0$, the upper bound $k \le (R/\gamma)^2$ becomes undefined ($\gamma \to 0$ implies $k \to \infty$).
  \item \textbf{The Inseparable Outlier and Boundary Thrashing}: Consider a dataset where an inseparable outlier (e.g., a negative point located within a cluster of positive samples) is encountered. Because the point is misclassified, the algorithm updates the weight vector:
  \begin{equation}
  \mathbf{w}^{(t+1)} = \mathbf{w}^{(t)} + y_i \mathbf{x}_i.
  \end{equation}
  As illustrated in Figure~\ref{fig:perceptron_flaw_nonseparable}, this update forcibly rotates the decision boundary from $H^{(t)}$ to $H^{(t+1)}$ to satisfy the outlier.
  \item \textbf{Perpetual Ping-Pong Oscillation}: Crucially, rotating the hyperplane to correctly classify the outlier inevitably moves the boundary across previously well-classified points, immediately causing them to become misclassified. In subsequent iterations, those newly misclassified points trigger counter-updates that rotate the boundary back toward $H^{(t)}$.
  \item \textbf{Infinite Looping}: The algorithm becomes trapped in an infinite ping-pong oscillation between conflicting hyperplanes, violently thrashing back and forth and \textbf{never terminating}.
\end{itemize}

To mitigate this catastrophic behavior in practical applications, Stephen Gallant introduced the heuristic \textbf{Pocket Algorithm}~\cite{gallant1990perceptron}, which runs the perceptron algorithm for a predetermined number of iterations while storing (``in its pocket'') the parameter vector that achieved the lowest classification error seen so far. Nevertheless, the fundamental theoretical guarantee of exact convergence is irrevocably lost on non-separable data.

\begin{figure}[H]
\centering
\begin{tikzpicture}[scale=1.15]
  \draw[-{Stealth}, gray!50] (-0.5, 0) -- (4.4, 0) node[right, font=\footnotesize] {$x_1$};
  \draw[-{Stealth}, gray!50] (0, -0.5) -- (0, 3.6) node[above, font=\footnotesize] {$x_2$};

  % Positive cluster (+1 stars)
  \node[star, star points=5, star point ratio=2.2, draw=politoNavy, fill=politoNavy!30, thick, inner sep=2pt] at (0.8, 2.5) {};
  \node[star, star points=5, star point ratio=2.2, draw=politoNavy, fill=politoNavy!30, thick, inner sep=2pt] at (1.3, 3.1) {};
  \node[star, star points=5, star point ratio=2.2, draw=politoNavy, fill=politoNavy!30, thick, inner sep=2pt] at (2.2, 3.0) {};
  \node[star, star points=5, star point ratio=2.2, draw=politoNavy, fill=politoNavy!30, thick, inner sep=2pt] at (2.5, 2.4) {};

  % Negative cluster (-1 circles)
  \node[circle, draw=examred, fill=examred!20, thick, inner sep=2pt] at (2.7, 0.7) {};
  \node[circle, draw=examred, fill=examred!20, thick, inner sep=2pt] at (3.4, 1.1) {};
  \node[circle, draw=examred, fill=examred!20, thick, inner sep=2pt] at (3.2, 0.4) {};

  % The Inseparable Outlier Negative point inside Positive cluster
  \node[circle, draw=examred, fill=examred!70, very thick, inner sep=2.5pt] (outlier) at (1.5, 2.0) {};
  \node[font=\scriptsize\bfseries, text=examred, above left] at (1.5, 2.0) {Inseparable Outlier ($-1$)};

  % Two alternating oscillating hyperplanes
  \draw[thick, deepblue, dashed] (0.2, 1.4) -- (3.8, 2.6) node[right, font=\scriptsize, text=deepblue] {Hyperplane $H^{(t)}$};
  \draw[thick, examred, dashdotted] (0.8, 3.4) -- (2.6, 0.3) node[below right, font=\scriptsize, text=examred] {Hyperplane $H^{(t+1)}$};

  % Oscillation cycle arrows in clear zone
  \draw[-{Stealth}, thick, deepdiveTeal] (3.3, 2.2) to[bend left=35] (2.9, 1.2);
  \draw[-{Stealth}, thick, deepdiveTeal] (2.7, 1.2) to[bend left=35] (3.1, 2.2);
  \node[font=\scriptsize\bfseries, text=deepdiveTeal, right, align=left] at (3.4, 1.7) {Perpetual\\Oscillation};

  \node[draw=examred, fill=examback, rounded corners, font=\scriptsize, align=center] at (2.0, -0.6) {Non-separable data: Algorithm oscillates indefinitely ($k \to \infty$).};
\end{tikzpicture}
\caption{Flaw 1 of the Perceptron: Non-separability and infinite cycling. When classes overlap or contain an inseparable outlier, no separating hyperplane exists ($\gamma \le 0$). The perceptron learning algorithm continually rotates to satisfy the misclassified outlier, which promptly causes previously classified points to fail, trapping the algorithm in an infinite ping-pong oscillation between conflicting hyperplanes $H^{(t)}$ and $H^{(t+1)}$.}
\label{fig:perceptron_flaw_nonseparable}
\end{figure}

The non-separability catastrophe demonstrates that the perceptron fails in the presence of noise and outliers. Yet a far deeper question remains: what if the data contains zero noise, zero outliers, and is completely deterministic? Can a single-layer perceptron learn even the simplest logical functions? As Marvin Minsky and Seymour Papert demonstrated, the answer is an emphatic no.

\paragraph{2. Inability to Compute Exclusive-OR (XOR).}
In their landmark 1969 monograph \textit{Perceptrons}, Marvin Minsky and Seymour Papert~\cite{minsky1969perceptrons} established that a single-layer perceptron is geometrically incapable of computing basic non-linear Boolean predicates, most famously the \textbf{Exclusive-OR (XOR)} function.

Consider the four canonical vertices of the unit square $[0, 1]^2$ defining the XOR logic gate:
\begin{equation}
\begin{aligned}
\mathbf{x}_1 = [0, 0]^\top &\implies y = -1, \\
\mathbf{x}_2 = [0, 1]^\top &\implies y = +1,
\end{aligned}
\qquad
\begin{aligned}
\mathbf{x}_3 = [1, 0]^\top &\implies y = +1, \\
\mathbf{x}_4 = [1, 1]^\top &\implies y = -1.
\end{aligned}
\label{eq:xor_dataset}
\end{equation}

To classify these four points correctly with zero training error, a perceptron parameterized by weights $(w_1, w_2)$ and bias $b$ must satisfy the four simultaneous strict linear inequalities:
\begin{align}
\text{For } \mathbf{x}_1 = [0, 0]^\top: \quad & w_1 \cdot 0 + w_2 \cdot 0 + b < 0 \implies b < 0, \label{eq:xor_ineq1} \\
\text{For } \mathbf{x}_2 = [0, 1]^\top: \quad & w_1 \cdot 0 + w_2 \cdot 1 + b \ge 0 \implies w_2 + b \ge 0, \label{eq:xor_ineq2} \\
\text{For } \mathbf{x}_3 = [1, 0]^\top: \quad & w_1 \cdot 1 + w_2 \cdot 0 + b \ge 0 \implies w_1 + b \ge 0, \label{eq:xor_ineq3} \\
\text{For } \mathbf{x}_4 = [1, 1]^\top: \quad & w_1 \cdot 1 + w_2 \cdot 1 + b < 0 \implies w_1 + w_2 + b < 0. \label{eq:xor_ineq4}
\end{align}

We now derive the formal algebraic impossibility proof:
\begin{proof}
Summing the two positive-class inequalities~\eqref{eq:xor_ineq2} and~\eqref{eq:xor_ineq3}:
\begin{equation}
(w_2 + b) + (w_1 + b) \ge 0 \implies w_1 + w_2 + 2b \ge 0 \implies w_1 + w_2 + b \ge -b.
\end{equation}
Because Equation~\eqref{eq:xor_ineq1} requires $b < 0$, multiplying by $-1$ yields $-b > 0$. Therefore:
\begin{equation}
w_1 + w_2 + b \ge -b > 0 \implies w_1 + w_2 + b > 0.
\label{eq:xor_contradiction_positive}
\end{equation}
However, Equation~\eqref{eq:xor_ineq4} strictly requires:
\begin{equation}
w_1 + w_2 + b < 0.
\label{eq:xor_contradiction_negative}
\end{equation}
Equations~\eqref{eq:xor_contradiction_positive} and~\eqref{eq:xor_contradiction_negative} are directly contradictory: a real quantity cannot be strictly positive and strictly negative simultaneously ($0 < w_1 + w_2 + b < 0$, a manifest contradiction in $\mathbb{R}$). 
\end{proof}

Hence, \textbf{no parameter vector $\mathbf{w} \in \mathbb{R}^2$ and bias $b \in \mathbb{R}$ can ever classify the XOR problem}. Geometrically, as illustrated in Figure~\ref{fig:perceptron_flaw_xor}, the positive points $(0, 1)$ and $(1, 0)$ and negative points $(0, 0)$ and $(1, 1)$ occupy opposite diagonals of the unit square. Any straight line drawn through the plane can at best isolate three of the four points, always misclassifying at least one vertex.

\begin{figure}[H]
\centering
\begin{tikzpicture}[scale=1.3]
  % Axes
  \draw[-{Stealth}, gray!70, thick] (-0.4, 0) -- (2.2, 0) node[right, font=\footnotesize] {$x_1$};
  \draw[-{Stealth}, gray!70, thick] (0, -0.4) -- (0, 2.2) node[above, font=\footnotesize] {$x_2$};
  \draw[help lines, color=gray!20, step=1] (-0.2,-0.2) grid (1.5,1.5);

  % XOR points with non-conflicting placement
  \node[circle, draw=examred, fill=examred!30, very thick, inner sep=3pt, label={below left:\footnotesize $(0,0): y=-1$}] at (0, 0) {};
  \node[circle, draw=examred, fill=examred!30, very thick, inner sep=3pt, label={above right:\footnotesize $(1,1): y=-1$}] at (1, 1) {};
  \node[star, star points=5, star point ratio=2.2, draw=politoNavy, fill=politoNavy!40, very thick, inner sep=3pt, label={above left:\footnotesize $(0,1): y=+1$}] at (0, 1) {};
  \node[star, star points=5, star point ratio=2.2, draw=politoNavy, fill=politoNavy!40, very thick, inner sep=3pt, label={below right:\footnotesize $(1,0): y=+1$}] at (1, 0) {};

  % Candidate hyperplanes failing (placed in open spaces)
  \draw[thick, gray!60, dashed] (-0.3, 0.5) -- (1.7, 0.5) node[right, font=\tiny, text=gray!80] {Candidate 1 ($x_2 = 0.5$)};
  \draw[thick, gray!60, dashed] (0.5, -0.3) -- (0.5, 1.7) node[above, font=\tiny, text=gray!80] {Candidate 2 ($x_1 = 0.5$)};
  \draw[thick, examred, dashdotted] (-0.3, 0.8) node[left, font=\tiny, text=examred] {Diagonal partition} -- (1.1, -0.2);

  \node[draw=examred, fill=examback, rounded corners, font=\scriptsize, align=center] at (0.9, -0.9) {Minsky \& Papert (1969): No single linear boundary can isolate opposite diagonal pairs.};
\end{tikzpicture}
\caption{Flaw 2 of the Perceptron: The XOR non-linear separability impasse. The Exclusive-OR (XOR) function assigns positive labels ($y = +1$) to coordinates $(0, 1)$ and $(1, 0)$, and negative labels ($y = -1$) to $(0, 0)$ and $(1, 1)$. Because positive and negative pairs occupy alternate diagonals of the unit square $[0, 1]^2$, no single linear decision hyperplane can separate the two classes.}
\label{fig:perceptron_flaw_xor}
\end{figure}

The publication of Minsky and Papert's critique had devastating historical ramifications. Because early AI researchers lacked an algorithm to train multi-layer neural architectures with hidden units, the field broadly concluded that artificial neural networks were fundamentally flawed and unscalable, causing government and academic funding to evaporate and initiating the decade-long First AI Winter.

\paragraph{3. Zero Derivative and Optimization Incompatibility.}
The third fatal flaw of Rosenblatt's perceptron lies in its activation function. The discontinuous step activation:
\begin{equation}
f(z) = \operatorname{sgn}(z) = \begin{cases} +1 & \text{if } z \ge 0, \\ -1 & \text{if } z < 0, \end{cases}
\end{equation}
is piecewise constant. Its derivative is zero everywhere that it is defined:
\begin{equation}
f'(z) = \frac{d}{dz}\operatorname{sgn}(z) = 0, \quad \forall z \ne 0,
\end{equation}
and is completely undefined (or infinite) at the discontinuity $z = 0$. 

Consequently, the parameter gradient of the pre-activation with respect to weights:
\begin{equation}
\nabla_\mathbf{w} f(\mathbf{w}^\top \mathbf{x} + b) = f'(z) \cdot \mathbf{x} = \mathbf{0}, \quad \forall z \ne 0.
\end{equation}
Because the gradient is identically zero almost everywhere, \textbf{gradient-based optimization algorithms cannot receive any directional learning signal}. Perturbing the weights by an infinitesimal step $\Delta \mathbf{w}$ either produces zero change in output or an abrupt, unpredictable jump across the step discontinuity. The perceptron cannot be trained using gradient descent, precluding integration into modern continuous optimization and backpropagation frameworks.

---

\section{Multilayer perceptrons, representation learning, and universal approximation}
\label{sec:mlp_and_representation_learning}

To overcome the perceptron's limitations, the architecture must undergo two simultaneous mathematical transformations: (1) transitioning from discrete step functions to continuous differentiable activations, and (2) stacking neurons into cascading \textbf{hidden layers} to construct \textbf{Multilayer Perceptrons (MLPs)}.

\subsection{Why we need activation functions and task-dependent probability squashing}
\label{sec:continuous_probabilities_softmax}

Before introducing specific activation functions, we must answer a foundational theoretical question: \textit{why can we not simply stack multiple layers of linear perceptrons without activation functions?}

\paragraph{The Compositional Collapse Theorem of Linear Networks.}
Suppose we cascade $L$ layers where each layer computes a pure affine transformation without non-linear activations:
\begin{align}
\mathbf{h}^{(1)} &= \mathbf{W}_1 \mathbf{x} + \mathbf{b}_1, \\
\mathbf{h}^{(2)} &= \mathbf{W}_2 \mathbf{h}^{(1)} + \mathbf{b}_2 = \mathbf{W}_2 (\mathbf{W}_1 \mathbf{x} + \mathbf{b}_1) + \mathbf{b}_2 = (\mathbf{W}_2 \mathbf{W}_1)\mathbf{x} + (\mathbf{W}_2 \mathbf{b}_1 + \mathbf{b}_2), \\
&\;\,\vdots \nonumber \\
\hat{\mathbf{y}} &= \mathbf{W}_L \mathbf{h}^{(L-1)} + \mathbf{b}_L = \left( \prod_{l=L}^1 \mathbf{W}_l \right) \mathbf{x} + \mathbf{b}_{\text{eff}} = \mathbf{W}_{\text{eff}} \mathbf{x} + \mathbf{b}_{\text{eff}}.
\end{align}
Because matrix multiplication is associative and affine transformations are closed under composition, \textbf{any multi-layer linear network collapses mathematically into a single linear transformation}! Stacking 100 linear layers provides zero additional representational capacity over a simple single-layer linear model. 

Non-linear activation functions are the sole mathematical mechanism that breaks affine subspace closure, enabling the network to bend, warp, and fold coordinate space to carve complex, non-linear decision boundaries.

\paragraph{Task-Dependent Activation Selection: Sigmoid vs. Softmax.}
Once non-linearity is introduced, the mathematical form of the output activation function must be chosen strictly according to the probabilistic structure of the target task:
\begin{enumerate}
  \item \textbf{Binary Classification ($y \in \{0, 1\}$ or $\{-1, +1\}$)}:
    The output is a single scalar probability modeling a Bernoulli likelihood $P(y = 1 \mid \mathbf{x})$. We apply the smooth \textbf{Logistic Sigmoid activation}:
    \begin{equation}
    \sigma(z) \triangleq \frac{1}{1 + e^{-z}} \in (0, 1), \quad \text{with } P(y = 1 \mid \mathbf{x}) = \sigma(\mathbf{w}^\top \mathbf{x} + b).
    \label{eq:sigmoid_formula}
    \end{equation}
    Sigmoid maps unconstrained real-valued logits $z \in (-\infty, +\infty)$ monotonically into the valid probability interval $(0, 1)$, with derivative $\sigma'(z) = \sigma(z)(1 - \sigma(z))$.
  \item \textbf{Multiclass Categorical Classification ($y \in \{1, \dots, K\}$ with $K > 2$)}:
    When classes are mutually exclusive, the output must be a probability vector residing on the $(K-1)$-dimensional probability simplex $\Delta^{K-1} = \{\mathbf{p} \in \R^K : p_k \ge 0, \, \sum_k p_k = 1\}$. We compute a vector of $K$ unconstrained logits:
    \begin{equation}
    \mathbf{z} = \mathbf{W}\mathbf{x} + \mathbf{b} \in \R^K,
    \end{equation}
    and apply the \textbf{Softmax activation function}:
    \begin{equation}
    \hat{y}_k \triangleq P(y = k \mid \mathbf{x}) = \operatorname{Softmax}(\mathbf{z})_k = \frac{\exp(z_k)}{\sum_{j=1}^K \exp(z_j)}, \quad \forall k \in \{1, \dots, K\}.
    \label{eq:softmax_definition}
    \end{equation}
\end{enumerate}

Softmax satisfies three essential analytical properties:
\begin{itemize}[noitemsep]
  \item \textbf{Strict Positivity}: $\hat{y}_k > 0$ for all finite logits $z_k$.
  \item \textbf{Unit Normalization}: $\sum_{k=1}^K \hat{y}_k = 1$.
  \item \textbf{Shift Invariance}: For any constant $c \in \R$, $\operatorname{Softmax}(\mathbf{z} + c \mathbf{1}) = \operatorname{Softmax}(\mathbf{z})$. Subtracting $\max_k z_k$ prevents numerical floating-point overflow.
\end{itemize}

Representing the ground truth label via a one-hot indicator vector $\mathbf{y} \in \{0, 1\}^K$, the corresponding loss is the \textbf{Categorical Cross-Entropy}:
\begin{equation}
\mathcal{L}_{\text{CE}}(\mathbf{y}, \hat{\mathbf{y}}) \triangleq -\sum_{k=1}^K y_k \ln \hat{y}_k = -\ln \hat{y}_{\text{true class}}.
\label{eq:categorical_cross_entropy}
\end{equation}
Differentiating with respect to logit $z_i$ reveals that the pre-activation gradient is simply the residual prediction error:
\begin{equation}
\frac{\partial \mathcal{L}_{\text{CE}}}{\partial z_i} = \hat{y}_i - y_i.
\label{eq:softmax_cross_entropy_gradient}
\end{equation}

\subsection{Resolving the XOR dilemma via hidden representation warping}
\label{sec:resolving_xor_hidden_representations}

Having established differentiable activations, we now resolve Minsky and Papert's XOR impasse by introducing an intermediate \textbf{hidden layer}.

\begin{figure}[H]
\centering
\begin{tikzpicture}[scale=1.0, every node/.style={transform shape}]
  % Layer 0: Input Layer
  \node[circle, draw=deepblue, fill=deepblue!10, thick, inner sep=3pt] (x1) at (0, 1.2) {$x_1$};
  \node[circle, draw=deepblue, fill=deepblue!10, thick, inner sep=3pt] (x2) at (0, -1.2) {$x_2$};
  \node[font=\footnotesize\bfseries\color{politoNavy}] at (0, 2.2) {Input $\mathbf{x} \in \R^2$};

  % Layer 1: Hidden Layer
  \node[circle, draw=intuitiongreen, fill=intuitiongreen!15, very thick, inner sep=3pt] (h1) at (3.2, 1.2) {$h_1$};
  \node[circle, draw=intuitiongreen, fill=intuitiongreen!15, very thick, inner sep=3pt] (h2) at (3.2, -1.2) {$h_2$};
  \node[font=\footnotesize\bfseries\color{intuitiongreen}] at (3.2, 2.2) {Hidden $\mathbf{h} = \sigma(\mathbf{W}_1\mathbf{x}+\mathbf{b}_1)$};

  % Layer 2: Output Layer
  \node[circle, draw=examred, fill=examred!15, very thick, inner sep=3pt] (y) at (6.4, 0.0) {$\hat{y}$};
  \node[font=\footnotesize\bfseries\color{examred}] at (6.4, 2.2) {Output $\hat{y} = \mathbf{w}_2^\top\mathbf{h}+b_2$};

  % Connections: Input to Hidden (W1)
  \draw[-{Stealth}, thick] (x1) -- node[above, font=\scriptsize] {$W_{11}=1$} (h1);
  \draw[-{Stealth}, thick] (x1) -- node[near start, above right, font=\scriptsize] {$W_{21}=1$} (h2);
  \draw[-{Stealth}, thick] (x2) -- node[near start, below right, font=\scriptsize] {$W_{12}=1$} (h1);
  \draw[-{Stealth}, thick] (x2) -- node[below, font=\scriptsize] {$W_{22}=1$} (h2);

  % Biases for Hidden
  \node[gray!70, font=\scriptsize] (b1_1) at (3.2, 0.4) {$b_1=0$};
  \node[gray!70, font=\scriptsize] (b1_2) at (3.2, -2.0) {$b_2=-1$};

  % Connections: Hidden to Output (w2)
  \draw[-{Stealth}, thick] (h1) -- node[above right, font=\scriptsize] {$w_{2,1}=1$} (y);
  \draw[-{Stealth}, thick] (h2) -- node[below right, font=\scriptsize] {$w_{2,2}=-2$} (y);

  % Output Arrow
  \draw[-{Stealth}, very thick, examred] (y) -- ++(1.2, 0) node[right, font=\footnotesize\bfseries] {$\hat{y} \in \{0, 1\}$};
\end{tikzpicture}
\caption{The Multi-Layer Perceptron (MLP) architecture resolving XOR. Input coordinates $(x_1, x_2)$ are projected into a 2D hidden layer $\mathbf{h} = \operatorname{ReLU}(\mathbf{W}_1 \mathbf{x} + \mathbf{b}_1)$ that extracts non-linear representations, followed by a linear output unit that classifies the transformed points cleanly.}
\label{fig:mlp_architecture}
\end{figure}

Consider the two-layer feedforward network depicted in Figure~\ref{fig:mlp_architecture}:
\begin{align}
\mathbf{h} &= \operatorname{ReLU}(\mathbf{W}_1 \mathbf{x} + \mathbf{b}_1), \quad \text{with } \mathbf{W}_1 = \begin{bmatrix} 1 & 1 \\ 1 & 1 \end{bmatrix}, \, \mathbf{b}_1 = \begin{bmatrix} 0 \\ -1 \end{bmatrix}, \label{eq:xor_hidden_forward} \\
\hat{y} &= \mathbf{w}_2^\top \mathbf{h} + b_2, \quad \text{with } \mathbf{w}_2 = \begin{bmatrix} 1 \\ -2 \end{bmatrix}, \, b_2 = 0, \label{eq:xor_output_forward}
\end{align}
where $\operatorname{ReLU}(z) \triangleq \max(0, z)$.

Evaluating the forward pass across the four XOR input configurations:
\begin{enumerate}
  \item \textbf{Input $\mathbf{x} = [0, 0]^\top$ (Negative Class)}:
    \begin{equation}
    \mathbf{W}_1 \mathbf{x} + \mathbf{b}_1 = \begin{bmatrix} 0 \\ -1 \end{bmatrix} \implies \mathbf{h} = \operatorname{ReLU}\left(\begin{bmatrix} 0 \\ -1 \end{bmatrix}\right) = \begin{bmatrix} 0 \\ 0 \end{bmatrix} \implies \hat{y} = [1, -2] \begin{bmatrix} 0 \\ 0 \end{bmatrix} = 0.
    \end{equation}
  \item \textbf{Input $\mathbf{x} = [1, 0]^\top$ (Positive Class)}:
    \begin{equation}
    \mathbf{W}_1 \mathbf{x} + \mathbf{b}_1 = \begin{bmatrix} 1 \\ 0 \end{bmatrix} \implies \mathbf{h} = \operatorname{ReLU}\left(\begin{bmatrix} 1 \\ 0 \end{bmatrix}\right) = \begin{bmatrix} 1 \\ 0 \end{bmatrix} \implies \hat{y} = [1, -2] \begin{bmatrix} 1 \\ 0 \end{bmatrix} = 1.
    \end{equation}
  \item \textbf{Input $\mathbf{x} = [0, 1]^\top$ (Positive Class)}:
    \begin{equation}
    \mathbf{W}_1 \mathbf{x} + \mathbf{b}_1 = \begin{bmatrix} 1 \\ 0 \end{bmatrix} \implies \mathbf{h} = \operatorname{ReLU}\left(\begin{bmatrix} 1 \\ 0 \end{bmatrix}\right) = \begin{bmatrix} 1 \\ 0 \end{bmatrix} \implies \hat{y} = [1, -2] \begin{bmatrix} 1 \\ 0 \end{bmatrix} = 1.
    \end{equation}
  \item \textbf{Input $\mathbf{x} = [1, 1]^\top$ (Negative Class)}:
    \begin{equation}
    \mathbf{W}_1 \mathbf{x} + \mathbf{b}_1 = \begin{bmatrix} 2 \\ 1 \end{bmatrix} \implies \mathbf{h} = \operatorname{ReLU}\left(\begin{bmatrix} 2 \\ 1 \end{bmatrix}\right) = \begin{bmatrix} 2 \\ 1 \end{bmatrix} \implies \hat{y} = [1, -2] \begin{bmatrix} 2 \\ 1 \end{bmatrix} = 0.
    \end{equation}
\end{enumerate}

\begin{figure}[H]
\centering
\begin{minipage}{0.48\linewidth}
  \centering
  \begin{tikzpicture}[scale=1.1]
    \draw[-{Stealth}, gray!70, thick] (-0.3, 0) -- (1.6, 0) node[right, font=\footnotesize] {$x_1$};
    \draw[-{Stealth}, gray!70, thick] (0, -0.3) -- (0, 1.6) node[above, font=\footnotesize] {$x_2$};
    \draw[help lines, color=gray!20, step=1] (-0.2,-0.2) grid (1.4,1.4);
    
    % Points
    \node[circle, draw=examred, fill=examred!20, thick, inner sep=3pt, label={below left:\scriptsize $(0,0): y=0$}] at (0, 0) {};
    \node[circle, draw=examred, fill=examred!20, thick, inner sep=3pt, label={above right:\scriptsize $(1,1): y=0$}] at (1, 1) {};
    \node[star, star points=5, star point ratio=2.2, draw=politoNavy, fill=politoNavy!30, thick, inner sep=2.5pt, label={above left:\scriptsize $(0,1): y=1$}] at (0, 1) {};
    \node[star, star points=5, star point ratio=2.2, draw=politoNavy, fill=politoNavy!30, thick, inner sep=2.5pt, label={below right:\scriptsize $(1,0): y=1$}] at (1, 0) {};

    \node[font=\footnotesize\bfseries, align=center, text=examred] at (0.5, -0.7) {Input Space $\mathcal{X}$:\\Non-linearly Separable};
  \end{tikzpicture}
\end{minipage}\hfill
\begin{minipage}{0.48\linewidth}
  \centering
  \begin{tikzpicture}[scale=1.1]
    \draw[-{Stealth}, gray!70, thick] (-0.3, 0) -- (2.4, 0) node[right, font=\footnotesize] {$h_1$};
    \draw[-{Stealth}, gray!70, thick] (0, -0.3) -- (0, 1.6) node[above, font=\footnotesize] {$h_2$};
    \draw[help lines, color=gray!20, step=1] (-0.2,-0.2) grid (2.2,1.4);

    % Transformed Points
    \node[circle, draw=examred, fill=examred!20, thick, inner sep=3pt, label={below left:\scriptsize $(0,0)$}] at (0, 0) {};
    \node[circle, draw=examred, fill=examred!20, thick, inner sep=3pt, label={above right:\scriptsize $(2,1)$}] at (2, 1) {};
    \node[star, star points=5, star point ratio=2.2, draw=politoNavy, fill=politoNavy!30, thick, inner sep=2.5pt, label={above:\scriptsize Coincident: $(1,0)$}] at (1, 0) {};

    % Separating hyperplane
    \draw[very thick, intuitiongreen] (0.5, -0.2) -- (1.5, 1.4) node[right, font=\scriptsize, align=left] {Hyperplane:\\$h_1 - 2h_2 = 0.5$};

    \node[font=\footnotesize\bfseries, align=center, text=intuitiongreen] at (1.0, -0.7) {Hidden Space $\mathcal{H}$:\\Linearly Separable!};
  \end{tikzpicture}
\end{minipage}
\caption{The geometric mechanism of representation learning resolving the XOR problem. Left: In the raw input space $\mathcal{X}$, no single line can separate the positive stars from the negative circles. Right: The non-linear hidden layer transformation $\mathbf{h} = \operatorname{ReLU}(\mathbf{W}_1 \mathbf{x} + \mathbf{b}_1)$ collapses both positive points $[0, 1]^\top$ and $[1, 0]^\top$ onto the single coordinate $[1, 0]^\top$, while mapping $[1, 1]^\top$ to $[2, 1]^\top$. In this hidden space $\mathcal{H}$, the classes are cleanly separated by a single linear hyperplane $h_1 - 2h_2 = 0.5$.}
\label{fig:xor_representation_warping}
\end{figure}

Figure~\ref{fig:xor_representation_warping} illuminates the deep geometric principle at work: \textbf{the hidden layer performs representation learning}. Instead of relying on handcrafted features, the network automatically warps and projects the input space into a latent coordinate system $\mathcal{H}$ where previously non-separable distributions become strictly linearly separable!

\subsection{From individual perceptrons to deep fully connected networks (FCNs)}
\label{sec:perceptron_to_fcn}

Building upon the single-hidden-layer MLP, we generalize this construction to arbitrary depth. By cascading multiple layers of perceptrons such that every neuron in layer $l-1$ connects to every neuron in layer $l$, we obtain a \textbf{Fully Connected Network (FCN)}, also termed a \textbf{Dense Feedforward Neural Network}.

Formally, an $L$-layer Fully Connected Network computes a sequence of coupled affine-non-linear vector recurrences:
\begin{align}
\mathbf{a}^{(0)} &\equiv \mathbf{x} \in \R^{d_0}, \label{eq:fcn_input} \\
\mathbf{z}^{(l)} &= \mathbf{W}^{(l)} \mathbf{a}^{(l-1)} + \mathbf{b}^{(l)}, \quad \mathbf{W}^{(l)} \in \R^{d_l \times d_{l-1}}, \, \mathbf{b}^{(l)} \in \R^{d_l}, \quad \forall l \in \{1, \dots, L\}, \label{eq:fcn_preactivation} \\
\mathbf{a}^{(l)} &= \sigma\left( \mathbf{z}^{(l)} \right) \in \R^{d_l}, \label{eq:fcn_activation}
\end{align}
where $\mathbf{z}^{(l)}$ is the layer-$l$ pre-activation vector, $\sigma(\cdot)$ is an element-wise activation function, and $\mathbf{a}^{(l)}$ denotes the layer-$l$ activations (with $\hat{\mathbf{y}} \equiv \mathbf{a}^{(L)}$).

The total parameter count of a dense FCN scales quadratically with layer width:
\begin{equation}
N_{\text{dense parameters}} = \sum_{l=1}^L \left( d_{l-1} \cdot d_l + d_l \right).
\label{eq:dense_parameter_count}
\end{equation}
As illustrated in Figure~\ref{fig:fcn_topology}, the bipartite graph between adjacent layers is complete: every neuron receives input from all neurons of the preceding layer.

\begin{figure}[H]
\centering
\begin{tikzpicture}[scale=0.95, every node/.style={transform shape}]
  % Styles
  \tikzstyle{inputnode}=[circle, draw=deepblue, fill=deepblue!10, thick, inner sep=2.5pt]
  \tikzstyle{hiddennode1}=[circle, draw=intuitiongreen, fill=intuitiongreen!15, thick, inner sep=2.5pt]
  \tikzstyle{hiddennode2}=[circle, draw=politoNavy, fill=politoNavy!15, thick, inner sep=2.5pt]
  \tikzstyle{outputnode}=[circle, draw=examred, fill=examred!20, thick, inner sep=2.5pt]

  % Layer 0: Input (d0 = 4)
  \foreach \i in {1,...,4} {
    \node[inputnode] (L0_\i) at (0, 2.5 - \i*1.0) {$x_\i$};
  }
  \node[font=\footnotesize\bfseries\color{deepblue}] at (0, 2.2) {Layer $\mathbf{a}^{(0)} \in \R^4$};

  % Layer 1: Hidden 1 (d1 = 5)
  \foreach \j in {1,...,5} {
    \node[hiddennode1] (L1_\j) at (3.0, 3.0 - \j*1.0) {$h^{(1)}_\j$};
  }
  \node[font=\footnotesize\bfseries\color{intuitiongreen}] at (3.0, 2.6) {Layer $\mathbf{a}^{(1)} \in \R^5$};

  % Layer 2: Hidden 2 (d2 = 5)
  \foreach \k in {1,...,5} {
    \node[hiddennode2] (L2_\k) at (6.0, 3.0 - \k*1.0) {$h^{(2)}_\k$};
  }
  \node[font=\footnotesize\bfseries\color{politoNavy}] at (6.0, 2.6) {Layer $\mathbf{a}^{(2)} \in \R^5$};

  % Layer 3: Output (d3 = 3)
  \foreach \m in {1,...,3} {
    \node[outputnode] (L3_\m) at (9.0, 2.0 - \m*1.0) {$\hat{y}_\m$};
  }
  \node[font=\footnotesize\bfseries\color{examred}] at (9.0, 1.6) {Layer $\mathbf{a}^{(3)} \in \R^3$};

  % Connect L0 to L1
  \foreach \i in {1,...,4} {
    \foreach \j in {1,...,5} {
      \draw[-{Stealth}, gray!40, thin] (L0_\i) -- (L1_\j);
    }
  }

  % Connect L1 to L2
  \foreach \j in {1,...,5} {
    \foreach \k in {1,...,5} {
      \draw[-{Stealth}, gray!40, thin] (L1_\j) -- (L2_\k);
    }
  }

  % Connect L2 to L3
  \foreach \k in {1,...,5} {
    \foreach \m in {1,...,3} {
      \draw[-{Stealth}, gray!40, thin] (L2_\k) -- (L3_\m);
    }
  }

  % Annotations for weight matrices
  \node[font=\scriptsize\bfseries\color{gray!80}] at (1.5, -2.5) {$\mathbf{W}^{(1)} \in \R^{5 \times 4}$};
  \node[font=\scriptsize\bfseries\color{gray!80}] at (4.5, -2.5) {$\mathbf{W}^{(2)} \in \R^{5 \times 5}$};
  \node[font=\scriptsize\bfseries\color{gray!80}] at (7.5, -2.5) {$\mathbf{W}^{(3)} \in \R^{3 \times 5}$};
\end{tikzpicture}
\caption{The architectural topology of a deep Fully Connected Network (FCN). Each layer $l$ performs an affine transformation followed by an element-wise non-linear activation: $\mathbf{a}^{(l)} = \sigma(\mathbf{W}^{(l)} \mathbf{a}^{(l-1)} + \mathbf{b}^{(l)})$. In this dense topology, the bipartite graph between adjacent layers is complete: every neuron receives input from all neurons of the preceding layer.}
\label{fig:fcn_topology}
\end{figure}

While fully connected architectures possess universal approximation properties, Equation~\eqref{eq:dense_parameter_count} reveals the seed of their catastrophic failure when confronted with high-dimensional image tensors.

\subsection{Activation functions and the deep ReLU revolution}
\label{sec:activation_functions_and_relu}

The choice of non-linear activation function $\sigma(\cdot)$ governs how gradients propagate across depth.

\paragraph{1. The Vanishing Gradient Bottleneck of Sigmoid and Tanh.}
Historically, early neural networks employed the logistic sigmoid $\sigma(z) = \frac{1}{1 + e^{-z}}$ or the hyperbolic tangent:
\begin{equation}
\tanh(z) = \frac{e^z - e^{-z}}{e^z + e^{-z}} = 2\sigma(2z) - 1.
\end{equation}
Both functions are smooth and bounded. However, their derivatives saturate asymptotically:
\begin{equation}
\frac{d\sigma(z)}{dz} = \sigma(z)(1 - \sigma(z)) \le 0.25, \quad \frac{d\tanh(z)}{dz} = 1 - \tanh^2(z) \le 1.0.
\end{equation}
Whenever a neuron's activation falls into the saturated tail regime ($|z| \gg 0$), the local derivative is nearly zero ($\sigma'(z) \approx 0$). In a network with $L$ cascading layers, backpropagating the gradient requires multiplying $L$ derivative terms by the chain rule:
\begin{equation}
\frac{\partial \mathcal{L}}{\partial \mathbf{z}^{(1)}} \propto \prod_{l=1}^L \mathbf{W}^{(l)\top} \operatorname{diag}\left( \sigma'(\mathbf{z}^{(l)}) \right).
\end{equation}
Because $\sigma'(z) \le 0.25$, this product decays exponentially with depth: $0.25^L \to 0$. By layer 4 or 5, the gradient signal vanishes completely ($\|\nabla_{\mathbf{W}^{(1)}} \mathcal{L}\| \approx 0$). Early layers stop learning entirely, freezing at their random initialization.

\paragraph{2. The Rectified Linear Unit (ReLU) Revolution.}
To solve the vanishing gradient problem, modern deep learning adopted the \textbf{Rectified Linear Unit (ReLU)}~\cite{glorot2011deep,nair2010rectified}:
\begin{equation}
\operatorname{ReLU}(z) \triangleq \max(0, z) = \begin{cases} z & \text{if } z \ge 0, \\ 0 & \text{if } z < 0. \end{cases}
\label{eq:relu_formula}
\end{equation}
Its derivative is piecewise constant:
\begin{equation}
\frac{d}{dz}\operatorname{ReLU}(z) = \begin{cases} 1 & \text{if } z > 0, \\ 0 & \text{if } z < 0. \end{cases}
\end{equation}
ReLU provides two transformative mathematical advantages:
\begin{itemize}[noitemsep]
  \item \textbf{Non-Saturating Gradient}: For all active neurons ($z > 0$), the derivative is exactly $1$, completely eliminating vanishing gradients regardless of network depth ($1^L = 1$).
  \item \textbf{Computational Efficiency}: Computing $\max(0, z)$ requires a single comparison instruction on modern CPUs/GPUs, avoiding expensive transcendental exponential evaluations ($e^{-z}$).
  \item \textbf{Sparse Representations}: Neurons with negative inputs output an exact zero, inducing sparse activation patterns where only a subset of pathways fire for any given input.
\end{itemize}

\paragraph{3. The ``Dying ReLU'' Pathological Failure: Mathematical Mechanics.}
Despite its non-saturating advantage for active units ($z > 0$), standard ReLU introduces an insidious optimization vulnerability known as the \textbf{``Dying ReLU''} problem. Under specific gradient dynamics, a single large update can permanently knock a neuron into the negative pre-activation regime across the entire data support, rendering it permanently inert and incapable of ever updating again.

We formalize the exact mechanical sequence of this failure:
\begin{enumerate}[leftmargin=*]
  \item \textbf{Neuron Mechanics and Gradient Transmission}: Consider a neuron parameterized by weight vector $\mathbf{w}$ and scalar bias $b$, evaluating input $\mathbf{x} \in \mathbb{R}^d$:
  \begin{equation}
  z = \mathbf{w}^\top \mathbf{x} + b, \qquad a = \operatorname{ReLU}(z) = \max(0, z).
  \end{equation}
  During backpropagation, the gradient of the scalar loss $\mathcal{L}$ with respect to the pre-activation $z$ is:
  \begin{equation}
  \frac{\partial \mathcal{L}}{\partial z} = \frac{\partial \mathcal{L}}{\partial a} \cdot \frac{da}{dz} = \begin{cases} \frac{\partial \mathcal{L}}{\partial a} & \text{if } z > 0, \\ 0 & \text{if } z < 0. \end{cases}
  \label{eq:relu_gradient_cases}
  \end{equation}
  Parameter updates under stochastic gradient descent (learning rate $\eta > 0$) evaluate to:
  \begin{equation}
  \mathbf{w}^{(t+1)} = \mathbf{w}^{(t)} - \eta \left( \frac{\partial \mathcal{L}}{\partial z} \mathbf{x} \right), \qquad b^{(t+1)} = b^{(t)} - \eta \frac{\partial \mathcal{L}}{\partial z}.
  \end{equation}

  \item \textbf{The Destructive Trigger (The Large Positive Gradient Step)}: Suppose the neuron is initially active ($z > 0$). If a sudden large error is backpropagated to this unit---triggered by an excessively high learning rate $\eta$, an unnormalized outlier batch, or unfortunate weight initialization---the incoming gradient signal $\frac{\partial \mathcal{L}}{\partial z}$ can be a massive positive scalar:
  \begin{equation}
  \frac{\partial \mathcal{L}}{\partial z} \gg 0.
  \end{equation}
  The gradient step subtracts $\eta \frac{\partial \mathcal{L}}{\partial z}$ directly from the bias parameter, driving it to an extreme negative value:
  \begin{equation}
  b^{(t+1)} = b^{(t)} - \eta \frac{\partial \mathcal{L}}{\partial z} \ll 0.
  \end{equation}
  Simultaneously, for any positive input features ($x_j > 0$), the weight update subtracts $\eta \frac{\partial \mathcal{L}}{\partial z} x_j$, driving the weights $\mathbf{w}$ strongly negative as well.

  \item \textbf{The Inactive Trap ($z < 0$ Across the Entire Data Distribution)}: Consider the neuron's pre-activation evaluated on any arbitrary sample $\mathbf{x}$ from the data distribution $\mathcal{D}$:
  \begin{equation}
  z^{(t+1)}(\mathbf{x}) = \mathbf{w}^{(t+1)\top} \mathbf{x} + b^{(t+1)}.
  \end{equation}
  Because $b^{(t+1)}$ is now deeply negative, the scalar offset dominates the bounded inner product $\mathbf{w}^{(t+1)\top} \mathbf{x}$. As a result, the pre-activation remains strictly negative across the entire training dataset:
  \begin{equation}
  z^{(t+1)}(\mathbf{x}) < 0, \quad \forall \mathbf{x} \sim \mathcal{D}.
  \end{equation}
  Consequently, the neuron emits zero activation for every sample:
  \begin{equation}
  a(\mathbf{x}) = \operatorname{ReLU}\left( z^{(t+1)}(\mathbf{x}) \right) = 0, \quad \forall \mathbf{x} \sim \mathcal{D}.
  \end{equation}

  \item \textbf{The Permanence: Zero-Gradient Lockout}: Why is the neuron unable to recover on subsequent iterations? On any future iteration $t+k$, for every incoming training sample $\mathbf{x}$:
  \begin{equation}
  z(\mathbf{x}) < 0 \implies \frac{da}{dz} = 0 \implies \frac{\partial \mathcal{L}}{\partial z} = \frac{\partial \mathcal{L}}{\partial a} \cdot 0 = 0.
  \end{equation}
  Because the local pre-activation gradient vanishes identically for all samples in every batch:
  \begin{equation}
  \nabla_\mathbf{w} \mathcal{L} = \frac{\partial \mathcal{L}}{\partial z} \mathbf{x} = \mathbf{0}, \qquad \frac{\partial \mathcal{L}}{\partial b} = \frac{\partial \mathcal{L}}{\partial z} = 0.
  \end{equation}
  The gradient update evaluates to a permanent identity:
  \begin{equation}
  \mathbf{w}^{(t+k+1)} = \mathbf{w}^{(t+k)} - \eta \mathbf{0} = \mathbf{w}^{(t+k)}, \qquad b^{(t+k+1)} = b^{(t+k)} - \eta \cdot 0 = b^{(t+k)}.
  \end{equation}
  Because the gradients are identically zero across all possible training batches, no gradient-based update can ever modify $\mathbf{w}$ or $b$ again. The neuron is permanently dead. When multiple neurons in a hidden layer suffer this lockout, the representational capacity of the network permanently and irreversibly collapses.
\end{enumerate}

\paragraph{Why Momentum and Weight Decay Fail to Resurrect Dead Neurons.}
Two standard optimizer mechanics fail to remedy this lockout:
\begin{itemize}[noitemsep]
  \item \textbf{Momentum}: The velocity vector updates as $\mathbf{v}_{t+1} = \beta \mathbf{v}_t + \eta \nabla_\theta \mathcal{L}_t$. Once gradients vanish identically across batches ($\nabla_\theta \mathcal{L} = \mathbf{0}$), velocity decays exponentially ($\mathbf{v}_{t+k} = \beta^k \mathbf{v}_t \to \mathbf{0}$), freezing parameters permanently.
  \item \textbf{Weight Decay ($L_2$ Regularization)}: Weight decay pulls weights toward zero ($\mathbf{w} \leftarrow (1 - \eta \lambda)\mathbf{w}$). However, biases $b$ are standardly excluded from weight decay to prevent underfitting; furthermore, even if $\mathbf{w} \to \mathbf{0}$, the pre-activation remains $z = \mathbf{0}^\top \mathbf{x} + b = b \ll 0$, which remains strictly negative and permanently locked.
\end{itemize}

\paragraph{4. Mathematical Resolution via Modern ReLU Variants.}
To prevent zero-gradient lockout, alternative non-linearities guarantee a strictly non-zero derivative in the negative regime ($z < 0$), ensuring that $\frac{da}{dz} \ne 0$ so corrective gradient signals can flow and pull deeply negative biases back up:
\begin{itemize}[leftmargin=*]
  \item \textbf{Leaky ReLU}~\cite{maas2013rectifier}: Replaces the zero-gradient flatline with a small positive slope $\alpha \approx 0.01$:
  \begin{equation}
  f(z) = \max(\alpha z, z) = \begin{cases} z & \text{if } z \ge 0, \\ \alpha z & \text{if } z < 0, \end{cases} \implies \frac{df}{dz} = \begin{cases} 1 & \text{if } z \ge 0, \\ \alpha & \text{if } z < 0. \end{cases}
  \end{equation}
  Even when knocked into $z < 0$, the neuron receives a live gradient $\frac{\partial \mathcal{L}}{\partial z} = \alpha \frac{\partial \mathcal{L}}{\partial a} \ne 0$, producing corrective parameter updates $b \leftarrow b - \eta \alpha \frac{\partial \mathcal{L}}{\partial a}$ that can restore the unit to active status.
  \item \textbf{Parametric ReLU (PReLU)}~\cite{he2015delving}: Treats the negative slope $\alpha$ as a learnable parameter optimized end-to-end via backpropagation ($\frac{\partial \mathcal{L}}{\partial \alpha} = \sum_{\mathbf{x}} \frac{\partial \mathcal{L}}{\partial a} \min(0, z)$).
  \item \textbf{Exponential Linear Unit (ELU)}~\cite{clevert2015fast}: Employs a saturating exponential curve for negative inputs:
  \begin{equation}
  f(z) = \begin{cases} z & \text{if } z \ge 0, \\ \alpha (e^z - 1) & \text{if } z < 0, \end{cases} \implies \frac{df}{dz} = \begin{cases} 1 & \text{if } z \ge 0, \\ f(z) + \alpha & \text{if } z < 0. \end{cases}
  \end{equation}
  ELU provides $C^1$ smoothness at $z = 0$ (when $\alpha = 1$) and drives mean activations closer to zero, accelerating convergence.
  \item \textbf{Architectural Safeguards}: In standard ReLU architectures, dying neurons are prevented by: (i) conservative initial learning rates with warm-up schedules; (ii) He/Kaiming initialization ($\operatorname{Var}(w) = 2/d_{\text{in}}$)~\cite{he2015delving}; and (iii) \textbf{Batch Normalization}, which re-centers and re-scales pre-activations before non-linearities, mathematically guaranteeing that roughly $50\%$ of inputs produce positive activations ($z > 0$).
\end{itemize}

\subsection{The Universal Approximation Theorem: Expressive power vs. learnability}
\label{sec:universal_approximation_theorem}

How expressive are neural networks? This fundamental question is answered by the \textbf{Universal Approximation Theorem}, independently proved by George Cybenko (1989)~\cite{cybenko1989approximation} for sigmoidal functions and generalized by Kurt Hornik (1991)~\cite{hornik1991approximation} for arbitrary non-linear activations.

\begin{theorem}[Universal Approximation Theorem~\cite{cybenko1989approximation,hornik1991approximation}]
Let $\sigma: \R \to \R$ be any continuous, non-constant, bounded activation function. Let $K \subset \R^d$ be a compact (closed and bounded) subset. Then, for any continuous target function $f \in C(K)$ and any arbitrary error tolerance $\epsilon > 0$, there exists an integer $M \in \mathbb{N}$, weight vectors $\mathbf{w}_j \in \R^d$, biases $b_j \in \R$, and output weights $\alpha_j \in \R$ such that the single-hidden-layer network:
\begin{equation}
F(\mathbf{x}) = \sum_{j=1}^M \alpha_j \sigma(\mathbf{w}_j^\top \mathbf{x} + b_j),
\label{eq:universal_approximator_form}
\end{equation}
satisfies uniform approximation over all $\mathbf{x} \in K$:
\begin{equation}
\sup_{\mathbf{x} \in K} |F(\mathbf{x}) - f(\mathbf{x})| < \epsilon.
\label{eq:universal_approximation_bound}
\end{equation}
\end{theorem}

\paragraph{Critical oral exam insight: Representation vs. learnability.}
In oral examinations, a pervasive misconception is asserting that the Universal Approximation Theorem guarantees that a single-hidden-layer network will solve any learning problem. As Yoshua Bengio emphasizes in \textit{Deep Learning}~\cite{goodfellow2016deep}, this is false:
\begin{enumerate}
  \item \textbf{Existence vs. Optimization}: The theorem is purely an \textbf{existential theorem}. It proves that valid parameters $(\mathbf{W}, \mathbf{b})$ exist, but provides zero guarantee that iterative optimization (gradient descent) can ever locate them in a non-convex landscape.
  \item \textbf{Exponential Width Bottleneck}: The required number of hidden units $M$ can scale exponentially with input dimensionality ($M \sim \mathcal{O}(2^d)$), rendering shallow approximation computationally intractable for high-dimensional data.
  \item \textbf{The Power of Depth}: Deep networks with $L$ cascading layers can represent compositional functions (such as hierarchical visual geometries) using a \textbf{polynomial number of parameters} $\mathcal{O}(d)$, whereas achieving identical approximation in a shallow single-layer network requires an exponential number of neurons $\mathcal{O}(2^d)$. Depth is statistically and computationally superior to width.
\end{enumerate}

---

\section{Algorithmic differentiation, backpropagation, and why we need gradient descent}
\label{sec:autodiff_backprop_gradient_descent}

\subsection{Why we need gradient descent: The impossibility of closed-form solutions}
\label{sec:why_we_need_gradient_descent}

In Section~\ref{sec:linear_predictors_and_least_squares}, we demonstrated that for linear regression, the optimal parameter vector admits a direct analytical matrix inversion:
\begin{equation}
\mathbf{w}^* = (\mathbf{X}^\top \mathbf{X})^{-1} \mathbf{X}^\top \mathbf{y}.
\end{equation}
Why can we not apply a similar closed-form algebraic solution to Multi-Layer Perceptrons and Deep Neural Networks?

Consider an $L$-layer neural network hypothesis:
\begin{equation}
f(\mathbf{x}; \boldsymbol{\theta}) = \mathbf{W}^{(L)} \sigma\left( \mathbf{W}^{(L-1)} \sigma\left( \dots \sigma\left( \mathbf{W}^{(1)} \mathbf{x} + \mathbf{b}^{(1)} \right) \dots \right) + \mathbf{b}^{(L-1)} \right) + \mathbf{b}^{(L)},
\label{eq:deep_composition_hypothesis}
\end{equation}
evaluated across training set $\mathcal{D} = \{(\mathbf{x}_i, y_i)\}_{i=1}^N$ under empirical risk objective:
\begin{equation}
\hat{\mathcal{R}}_N(\boldsymbol{\theta}) = \frac{1}{N} \sum_{i=1}^N \mathcal{L}\left( f(\mathbf{x}_i; \boldsymbol{\theta}), \, y_i \right).
\end{equation}

To find the minimizer analytically, one must equate the gradient with respect to all parameter matrices $\mathbf{W}^{(l)}$ to zero:
\begin{equation}
\nabla_{\boldsymbol{\theta}} \hat{\mathcal{R}}_N(\boldsymbol{\theta}) = \mathbf{0}.
\label{eq:stationary_condition_deep}
\end{equation}
Equation~\eqref{eq:stationary_condition_deep} represents a system of millions of coupled, transcendental non-linear algebraic equations. Because parameters $\mathbf{W}^{(1)}$ are deeply entangled inside cascading non-linear activation functions $\sigma(\cdot)$, they \textbf{cannot be decoupled, factorized, or inverted through linear algebra}.

Furthermore, as established in Chapter~\ref{sec:optimization_and_gradient_descent}, the empirical loss surface of a deep neural network is non-convex, characterized by uncountably many saddle points, local valleys, and permutation symmetries (weight space symmetries where reordering hidden neurons yields identical loss).

Because closed-form solutions are mathematically impossible, \textbf{iterative numerical optimization—specifically first-order Gradient Descent—is the only viable mathematical vehicle for training deep learning models}.

\subsection{Computational graphs and automatic differentiation: Forward vs. reverse mode}
\label{sec:computational_graphs_autodiff}

To execute gradient descent, we must compute the exact gradient vector $\nabla_{\boldsymbol{\theta}} \mathcal{L}$ for networks containing millions to billions of parameters. 

\paragraph{1. The Failure of Numerical Finite Differences.}
One naive approach is numerical perturbation (finite differences):
\begin{equation}
\frac{\partial \mathcal{L}}{\partial \theta_j} \approx \frac{\mathcal{L}(\boldsymbol{\theta} + \epsilon \mathbf{e}_j) - \mathcal{L}(\boldsymbol{\theta})}{\epsilon}.
\end{equation}
For a model with $W$ total parameters, evaluating this quotient requires $W + 1$ forward passes through the network. In modern networks where $W \sim 10^8$ (AlexNet, ResNet), computing a single parameter update step would require 100 million network evaluations, leading to an intractable computational complexity of $\mathcal{O}(W^2)$.

\paragraph{2. The Computational Graph.}
Modern frameworks (PyTorch, JAX) formulate deep networks as \textbf{Directed Acyclic Graphs (DAGs)} where nodes represent variables (tensors) and directed edges represent primitive differentiable operations ($+$, $\times$, $\sigma$, $\ln$).

\begin{figure}[H]
\centering
\begin{tikzpicture}[scale=1.05]
  % Nodes
  \node[circle, draw=deepblue, fill=deepblue!15, thick, inner sep=3pt] (x) at (0, 1.2) {$x$};
  \node[circle, draw=deepblue, fill=deepblue!15, thick, inner sep=3pt] (w) at (0, -0.2) {$w$};
  \node[circle, draw=politoNavy, fill=politoNavy!20, thick, inner sep=3pt] (mul) at (2.0, 0.5) {$\times$};
  \node[circle, draw=deepblue, fill=deepblue!15, thick, inner sep=3pt] (b) at (2.0, -1.0) {$b$};
  \node[circle, draw=politoNavy, fill=politoNavy!20, thick, inner sep=3pt] (add) at (4.0, 0.0) {$+$};
  \node[circle, draw=intuitiongreen, fill=intuitiongreen!20, thick, inner sep=3pt] (sig) at (6.0, 0.0) {$\sigma$};
  \node[circle, draw=deepblue, fill=deepblue!15, thick, inner sep=3pt] (y) at (6.0, -1.2) {$y$};
  \node[circle, draw=examred, fill=examred!25, thick, inner sep=3pt] (loss) at (8.0, -0.6) {$\mathcal{L}$};

  % Forward Edges (Blue)
  \draw[-{Stealth}, very thick, deepblue] (x) -- node[above, font=\scriptsize] {$x$} (mul);
  \draw[-{Stealth}, very thick, deepblue] (w) -- node[below, font=\scriptsize] {$w$} (mul);
  \draw[-{Stealth}, very thick, deepblue] (mul) -- node[above, font=\scriptsize] {$wx$} (add);
  \draw[-{Stealth}, very thick, deepblue] (b) -- node[below, font=\scriptsize] {$b$} (add);
  \draw[-{Stealth}, very thick, deepblue] (add) -- node[above, font=\scriptsize] {$z$} (sig);
  \draw[-{Stealth}, very thick, deepblue] (sig) -- node[above, font=\scriptsize] {$\hat{y}$} (loss);
  \draw[-{Stealth}, very thick, deepblue] (y) -- node[below, font=\scriptsize] {$y$} (loss);

  % Backward Adjoint Flow (Red Dashed)
  \draw[-{Stealth}, thick, examred, dashed] (loss) to[bend left=30] node[above, font=\tiny\bfseries] {$\frac{\partial \mathcal{L}}{\partial \hat{y}}$} (sig);
  \draw[-{Stealth}, thick, examred, dashed] (sig) to[bend left=30] node[above, font=\tiny\bfseries] {$\frac{\partial \mathcal{L}}{\partial z}$} (add);
  \draw[-{Stealth}, thick, examred, dashed] (add) to[bend left=30] node[above, font=\tiny\bfseries] {$\frac{\partial \mathcal{L}}{\partial (wx)}$} (mul);
  \draw[-{Stealth}, thick, examred, dashed] (add) to[bend left=30] node[below, font=\tiny\bfseries] {$\frac{\partial \mathcal{L}}{\partial b}$} (b);
  \draw[-{Stealth}, thick, examred, dashed] (mul) to[bend left=30] node[below, font=\tiny\bfseries] {$\frac{\partial \mathcal{L}}{\partial w}$} (w);

  % Legend
  \node[draw=gray!60, fill=white, rounded corners, font=\scriptsize, align=left] at (4.0, 2.0) {
    \textcolor{deepblue}{\textbf{--- Forward Pass}}: Activations evaluated and cached\\
    \textcolor{examred}{\textbf{- - Backward Pass}}: Adjoint sensitivities propagated via chain rule
  };
\end{tikzpicture}
\caption{A Directed Acyclic Graph (DAG) representation of a computational graph for an artificial neuron with scalar loss $\mathcal{L}(\sigma(wx + b), y)$. Forward evaluation (solid blue lines) evaluates and caches intermediate variables. Reverse-mode automatic differentiation (dashed red lines) propagates adjoint sensitivities backwards from the scalar loss.}
\label{fig:computational_graph}
\end{figure}

\paragraph{3. Forward-Mode vs. Reverse-Mode Automatic Differentiation.}
Let $f: \R^n \to \R^m$ be a differentiable function computed on a graph:
\begin{itemize}
  \item \textbf{Forward-Mode Autodiff (Tangent Propagation)}: Propagates directional derivatives $\dot{v} \triangleq \frac{\partial v}{\partial x_i}$ forward along the graph alongside values $v$. To compute the full Jacobian $\mathbf{J} \in \R^{m \times n}$, forward-mode requires $n$ independent sweeps. It is computationally optimal when $n \ll m$ (few inputs, many outputs).
  \item \textbf{Reverse-Mode Autodiff (Adjoint Propagation)}: First evaluates the forward pass to compute and cache all intermediate variables $v$. Then, starting from the final scalar objective, it propagates adjoint sensitivity variables backwards:
    \begin{equation}
    \bar{v} \triangleq \frac{\partial \mathcal{L}}{\partial v}.
    \label{eq:adjoint_definition}
    \end{equation}
    Reverse-mode requires $m$ backward sweeps. 
\end{itemize}

In deep learning, the loss function is a \textbf{scalar} ($\mathcal{L} \in \R$, so $m = 1$), while the parameter count is massive ($n = W \sim 10^8$). Reverse-mode automatic differentiation evaluates the \textbf{entire parameter gradient $\nabla_{\boldsymbol{\theta}} \mathcal{L} \in \R^W$ in exactly ONE backward pass}, reducing computational complexity from $\mathcal{O}(W^2)$ to $\mathcal{O}(W)$! This single-pass reverse mode is known as \textbf{Backpropagation}.

\subsection{The Backpropagation algorithm: Full mathematical derivation}
\label{sec:backpropagation_derivation}

We now derive the backpropagation algorithm from first principles across an arbitrary $L$-layer network.

\paragraph{Forward Pass Equations.}
For layers $l = 1, \dots, L$:
\begin{align}
\mathbf{z}^{(l)} &= \mathbf{W}^{(l)} \mathbf{a}^{(l-1)} + \mathbf{b}^{(l)}, \quad \text{where } \mathbf{z}^{(l)} \in \R^{d_l}, \, \mathbf{W}^{(l)} \in \R^{d_l \times d_{l-1}}, \label{eq:forward_z} \\
\mathbf{a}^{(l)} &= \sigma\left( \mathbf{z}^{(l)} \right) \in \R^{d_l}, \label{eq:forward_a}
\end{align}
with input initialization $\mathbf{a}^{(0)} \equiv \mathbf{x} \in \R^{d_0}$, and final scalar loss $\mathcal{L}(\mathbf{a}^{(L)}, \mathbf{y}) \in \R$.

\paragraph{Backward Pass: The Error Adjoint Vector $\boldsymbol{\delta}^{(l)}$.}
We define the layer error vector $\boldsymbol{\delta}^{(l)}$ as the partial derivative of the scalar loss with respect to the pre-activation vector $\mathbf{z}^{(l)}$:
\begin{equation}
\boldsymbol{\delta}^{(l)} \triangleq \frac{\partial \mathcal{L}}{\partial \mathbf{z}^{(l)}} \in \R^{d_l}.
\label{eq:delta_definition}
\end{equation}

\paragraph{1. Base Case: Output Layer Error $\boldsymbol{\delta}^{(L)}$.}
Applying the multivariate chain rule to the final layer:
\begin{equation}
\delta_k^{(L)} = \frac{\partial \mathcal{L}}{\partial z_k^{(L)}} = \sum_{j=1}^{d_L} \frac{\partial \mathcal{L}}{\partial a_j^{(L)}} \frac{\partial a_j^{(L)}}{\partial z_k^{(L)}} = \frac{\partial \mathcal{L}}{\partial a_k^{(L)}} \sigma'\left( z_k^{(L)} \right).
\end{equation}
In vector notation:
\begin{equation}
\boldsymbol{\delta}^{(L)} = \nabla_{\mathbf{a}^{(L)}} \mathcal{L} \odot \sigma'\left( \mathbf{z}^{(L)} \right),
\label{eq:delta_output_layer}
\end{equation}
where $\odot$ denotes the Hadamard (element-wise) product. For cross-entropy coupled with softmax or sigmoid, this simplifies directly to $\boldsymbol{\delta}^{(L)} = \hat{\mathbf{y}} - \mathbf{y}$.

\paragraph{2. Inductive Step: Error Backpropagation Recurrence $\boldsymbol{\delta}^{(l)}$.}
Suppose the error vector $\boldsymbol{\delta}^{(l+1)} \in \R^{d_{l+1}}$ at layer $l+1$ is known. We compute $\boldsymbol{\delta}^{(l)} \in \R^{d_l}$ by differentiating through layer $l+1$. By the multivariate chain rule:
\begin{equation}
\delta_j^{(l)} = \frac{\partial \mathcal{L}}{\partial z_j^{(l)}} = \sum_{k=1}^{d_{l+1}} \frac{\partial \mathcal{L}}{\partial z_k^{(l+1)}} \frac{\partial z_k^{(l+1)}}{\partial z_j^{(l)}}.
\label{eq:delta_chain_component}
\end{equation}
From Equation~\eqref{eq:forward_z}, $z_k^{(l+1)} = \sum_{m=1}^{d_l} W_{km}^{(l+1)} a_m^{(l)} + b_k^{(l+1)} = \sum_{m=1}^{d_l} W_{km}^{(l+1)} \sigma(z_m^{(l)}) + b_k^{(l+1)}$. Differentiating with respect to $z_j^{(l)}$:
\begin{equation}
\frac{\partial z_k^{(l+1)}}{\partial z_j^{(l)}} = W_{kj}^{(l+1)} \sigma'\left( z_j^{(l)} \right).
\end{equation}
Substituting this into Equation~\eqref{eq:delta_chain_component} and recalling that $\delta_k^{(l+1)} = \frac{\partial \mathcal{L}}{\partial z_k^{(l+1)}}$:
\begin{equation}
\delta_j^{(l)} = \sum_{k=1}^{d_{l+1}} \delta_k^{(l+1)} W_{kj}^{(l+1)} \sigma'\left( z_j^{(l)} \right) = \left( \sum_{k=1}^{d_{l+1}} W_{kj}^{(l+1)} \delta_k^{(l+1)} \right) \sigma'\left( z_j^{(l)} \right).
\end{equation}
In compact vector matrix form:
\begin{equation}
\boldsymbol{\delta}^{(l)} = \left( (\mathbf{W}^{(l+1)})^\top \boldsymbol{\delta}^{(l+1)} \right) \odot \sigma'\left( \mathbf{z}^{(l)} \right).
\label{eq:backprop_recurrence_matrix}
\end{equation}
Equation~\eqref{eq:backprop_recurrence_matrix} is the celebrated \textbf{Backpropagation Recurrence}: the error signal flows backwards by multiplying with the transposed weight matrix $(\mathbf{W}^{(l+1)})^\top$, followed by element-wise modulation by the local derivative $\sigma'(\mathbf{z}^{(l)})$.

\paragraph{3. Parameter Gradients.}
Once $\boldsymbol{\delta}^{(l)}$ is computed, the gradients with respect to parameters at layer $l$ follow directly:
\begin{align}
\frac{\partial \mathcal{L}}{\partial W_{jk}^{(l)}} &= \frac{\partial \mathcal{L}}{\partial z_j^{(l)}} \frac{\partial z_j^{(l)}}{\partial W_{jk}^{(l)}} = \delta_j^{(l)} a_k^{(l-1)} \implies \frac{\partial \mathcal{L}}{\partial \mathbf{W}^{(l)}} = \boldsymbol{\delta}^{(l)} (\mathbf{a}^{(l-1)})^\top \in \R^{d_l \times d_{l-1}}, \label{eq:grad_weight_matrix} \\
\frac{\partial \mathcal{L}}{\partial b_j^{(l)}} &= \frac{\partial \mathcal{L}}{\partial z_j^{(l)}} \frac{\partial z_j^{(l)}}{\partial b_j^{(l)}} = \delta_j^{(l)} \cdot 1 \implies \frac{\partial \mathcal{L}}{\partial \mathbf{b}^{(l)}} = \boldsymbol{\delta}^{(l)} \in \R^{d_l}. \label{eq:grad_bias_vector}
\end{align}

\paragraph{4. Memory Bottleneck: Forward Activation Caching.}
Observe the structure of Equation~\eqref{eq:grad_weight_matrix}: evaluating $\frac{\partial \mathcal{L}}{\partial \mathbf{W}^{(l)}}$ requires the forward activation vector $\mathbf{a}^{(l-1)}$. Consequently, all intermediate activations computed during the forward pass must be stored in memory (GPU VRAM) until the backward pass reaches that layer. This requirement explains why training deep networks consumes vastly more memory than inference.

---

\section{The visual challenge and the structural failure of fully connected networks}
\label{sec:failure_of_fc_on_images}

Having derived Fully Connected Networks (FCNs) in Section~\ref{sec:perceptron_to_fcn}, we now confront an empirical crisis: \textit{why does deploying dense fully connected networks on high-dimensional visual sensory data fail completely?}

\subsection{Why MLPs fail on high-dimensional visual data}
\label{sec:why_mlps_fail_images}

\begin{figure}[H]
\centering
\begin{minipage}{0.48\linewidth}
  \centering
  \begin{tikzpicture}[scale=0.88]
    \node[font=\footnotesize\bfseries\color{politoNavy}] at (2.2, 3.4) {(a) Dense Layer (Flattened Vector)};

    % 2D Image Grid
    \draw[fill=gray!10, draw=gray!60, thick] (0, 0.5) rectangle (1.6, 2.1);
    \node[font=\tiny, align=center] at (0.8, 1.3) {$32 \times 32 \times 3$\\Image};

    % Flatten Arrow
    \draw[-{Stealth}, thick, gray!70] (1.7, 1.3) -- (2.3, 1.3) node[midway, above, font=\tiny] {Flatten};

    % 1D Vector (d = 3072)
    \draw[fill=deepblue!20, draw=deepblue, thick] (2.5, -0.4) rectangle (2.9, 3.0);
    \node[font=\tiny, rotate=90] at (2.7, 1.3) {Vector: $3072 \times 1$};

    % Dense Connections
    \foreach \y in {0.0, 0.8, 1.6, 2.4} {
      \draw[gray!40, thin] (2.9, \y) -- (4.2, 2.0);
      \draw[gray!40, thin] (2.9, \y) -- (4.2, 1.2);
      \draw[gray!40, thin] (2.9, \y) -- (4.2, 0.4);
    }

    % Hidden Vector
    \draw[fill=intuitiongreen!20, draw=intuitiongreen, thick] (4.2, 0.0) rectangle (4.6, 2.4);
    \node[font=\tiny, rotate=90] at (4.4, 1.2) {Hidden: $1024 \times 1$};

    \node[font=\scriptsize\color{examred}, align=center] at (2.2, -1.0) {Destroys 2D topology\\$3.15 \times 10^6$ weights};
  \end{tikzpicture}
\end{minipage}\hfill
\begin{minipage}{0.48\linewidth}
  \centering
  \begin{tikzpicture}[scale=0.88]
    \node[font=\footnotesize\bfseries\color{politoNavy}] at (2.2, 3.4) {(b) Convolutional Layer (3D Volume)};

    % 3D Input Volume: 32x32x3
    \draw[fill=deepblue!15, draw=deepblue, thick] (0, 0) -- (1.6, 0) -- (1.6, 1.6) -- (0, 1.6) -- cycle;
    \draw[fill=deepblue!25, draw=deepblue, thick] (1.6, 0) -- (2.1, 0.5) -- (2.1, 2.1) -- (1.6, 1.6) -- cycle;
    \draw[fill=deepblue!10, draw=deepblue, thick] (0, 1.6) -- (0.5, 2.1) -- (2.1, 2.1) -- (1.6, 1.6) -- cycle;
    \node[font=\tiny, align=center] at (0.8, 0.8) {$32 \times 32$\\$\times 3$};

    % Small 3D Receptive Field Filter: 5x5x3
    \draw[fill=examred!30, draw=examred, thick] (0.4, 0.6) rectangle (0.9, 1.1);
    \node[font=\tiny\color{examred}] at (0.65, 0.85) {$F$};

    % Filter arrow
    \draw[-{Stealth}, very thick, examred] (1.1, 0.85) -- (2.8, 0.85);

    % 3D Output Feature Volume
    \draw[fill=intuitiongreen!15, draw=intuitiongreen, thick] (3.0, 0.2) -- (4.2, 0.2) -- (4.2, 1.4) -- (3.0, 1.4) -- cycle;
    \draw[fill=intuitiongreen!25, draw=intuitiongreen, thick] (4.2, 0.2) -- (4.6, 0.6) -- (4.6, 1.8) -- (4.2, 1.4) -- cycle;
    \draw[fill=intuitiongreen!10, draw=intuitiongreen, thick] (3.0, 1.4) -- (3.4, 1.8) -- (4.6, 1.8) -- (4.2, 1.4) -- cycle;
    \node[font=\tiny, align=center] at (3.6, 0.8) {$28 \times 28$\\$\times 6$};

    \node[font=\scriptsize\color{intuitiongreen}, align=center] at (2.2, -1.0) {Preserves 3D spatial grid\\Reuses shared weights};
  \end{tikzpicture}
\end{minipage}
\caption{The structural contrast between dense and convolutional representations. (a) A dense fully connected layer flattens an image into an unrolled 1D vector, destroying all 2D spatial adjacency and requiring an independent weight for every pixel-neuron pair. (b) A convolutional layer preserves the native 3D tensor volume, sliding localized spatial filters across coordinates with shared parameters.}
\label{fig:dense_vs_conv_geometry}
\end{figure}

As illustrated in Figure~\ref{fig:dense_vs_conv_geometry}, standard fully connected networks suffer from three crippling structural flaws when applied to visual data:

\paragraph{1. The Vectorization Requirement and Destruction of 2D Spatial Topology.}
To ingest visual sensory data into a Multi-Layer Perceptron, one must reconcile a fundamental dimensional mismatch. An artificial neuron computes an affine inner product followed by an element-wise activation:
\begin{equation}
z = \mathbf{w}^\top \mathbf{x} + b = \sum_{k=1}^d w_k x_k + b.
\end{equation}
Because the canonical Euclidean inner product and standard matrix-vector multiplication $\mathbf{W}\mathbf{x}$ are defined exclusively over one-dimensional coordinate vectors $\mathbf{x} \in \R^d$, a dense layer \textbf{cannot natively ingest multidimensional array grids}. 

Consequently, any multi-channel visual input tensor $\mathbf{X} \in \R^{H \times W \times C}$ must undergo mandatory \textbf{vectorization (flattening)} into an unrolled 1D column vector:
\begin{equation}
\mathbf{x} = \operatorname{vec}(\mathbf{X}) \in \R^{d}, \quad \text{where } d = H \cdot W \cdot C,
\label{eq:flattening_transformation}
\end{equation}
under which each spatial coordinate $(i, j)$ in channel $c$ is mapped to a serialized linear index $k(i, j, c) = c \cdot (HW) + i \cdot W + j$.

This flattening operation fundamentally shatters the geometrical structure of the image:
\begin{itemize}[leftmargin=*]
  \item \textbf{Metric Space Disruption}: In the native continuous 2D image plane, geometry is governed by the standard Euclidean distance metric between pixel coordinates:
  \begin{equation}
  d_{\text{spatial}}(\mathbf{p}_1, \mathbf{p}_2) = \|\mathbf{p}_1 - \mathbf{p}_2\|_2 = \sqrt{(i_1 - i_2)^2 + (j_1 - j_2)^2}.
  \end{equation}
  Pixels that are immediate vertical neighbors, such as $(i, j)$ and $(i+1, j)$, have a spatial distance of $1$. However, in the flattened vector $\mathbf{x}$, their serialized indices are separated by an entire row of $W$ elements: $|k_1 - k_2| = W$. Conversely, pixels at the end of one row $(i, W)$ and the start of the next row $(i+1, 1)$ appear adjacent in 1D index space ($|k_1 - k_2| = 1$), despite residing on opposite physical edges of the image!
  \item \textbf{Loss of Locality and Spatial Adjacency}: Natural visual semantics (edges, textures, object contours) emerge entirely from \textbf{local spatial coherence}: adjacent pixels share physical illumination, surface reflectance, and material properties. By discarding the 2D grid, an MLP discards the concept of a neighborhood $\mathcal{N}(i, j)$, treating a pixel's relationship with its immediate neighbor identically to its relationship with a distant pixel hundreds of scanlines away.
  \item \textbf{The Permutation Invariance Flaw}: Because a dense weight matrix $\mathbf{W} \in \R^{M \times d}$ assigns an unconstrained independent parameter $W_{m, k}$ to every coordinate $k$, an MLP possesses \textbf{zero spatial inductive bias}. To demonstrate this formally, let $\mathbf{P}_\pi \in \{0, 1\}^{d \times d}$ be an arbitrary fixed permutation matrix that randomly shuffles the $d$ pixel indices, transforming every image into humanly unrecognizable white noise static. If we simultaneously permute the columns of the weight matrix by $\mathbf{W}' = \mathbf{W}\mathbf{P}_\pi^\top$, then for any scrambled input $\mathbf{x}' = \mathbf{P}_\pi \mathbf{x}$:
  \begin{equation}
  \mathbf{W}' \mathbf{x}' = (\mathbf{W}\mathbf{P}_\pi^\top)(\mathbf{P}_\pi \mathbf{x}) = \mathbf{W}(\mathbf{P}_\pi^\top \mathbf{P}_\pi)\mathbf{x} = \mathbf{W}\mathbf{x}.
  \end{equation}
  The MLP produces the \textit{exact same forward pre-activations, the exact same backpropagated gradients, and the exact same loss trajectory} on scrambled white noise as it does on structured natural images! An MLP cannot tell whether its input was a photographic scene or a scrambled column of tabular database numbers.
\end{itemize}

\paragraph{2. Zero Weight Sharing and Total Lack of Translation Equivariance.}
In an unconstrained dense layer, every hidden neuron $m$ connects to every input coordinate $k = (i, j, c)$ via a unique, dedicated weight parameter $W_{m, (i, j, c)}$:
\begin{equation}
z_m = \sum_{c=1}^C \sum_{i=1}^H \sum_{j=1}^W W_{m, (i, j, c)} X_{i, j, c} + b_m.
\label{eq:dense_coordinate_binding}
\end{equation}
This formulation has two catastrophic consequences:
\begin{itemize}[leftmargin=*]
  \item \textbf{Coordinate-Specific Parameter Binding (Zero Weight Sharing)}: In Equation~\eqref{eq:dense_coordinate_binding}, the weight $W_{m, (i, j, c)}$ is bound to the absolute spatial coordinate $(i, j)$ on the canvas. There is \textbf{zero parameter sharing}: the weight that examines coordinate $(10, 10)$ shares no algebraic relationship with the weight examining coordinate $(150, 150)$. To share parameters across spatial locations, the matrix $\mathbf{W}$ would have to be constrained to have a doubly block-circulant (or Toeplitz) structure---which is precisely the mathematical definition of a discrete convolution! Unconstrained MLPs lack this structure entirely.
  \item \textbf{Breakdown of Translation Equivariance}: Physical objects satisfy \textbf{translational symmetry}: a cat, a pedestrian, or a traffic sign retains identical identity, scale, and visual features regardless of whether it appears in the upper-left corner, the center, or the lower-right corner of the visual field. Mathematically, a representation mapping $f: \mathcal{X} \to \mathcal{H}$ is \textbf{translation-equivariant} with respect to spatial shift $T_{\mathbf{v}}$ if:
  \begin{equation}
  f(T_{\mathbf{v}}(\mathbf{X})) = T_{\mathbf{v}}(f(\mathbf{X})), \quad \forall \mathbf{v} \in \mathbb{Z}^2.
  \end{equation}
  In an MLP, because weights are rigidly tied to absolute pixel coordinates, $f(T_{\mathbf{v}}(\mathbf{x})) \ne T_{\mathbf{v}}(f(\mathbf{x}))$. If a hidden neuron learns a set of weights that successfully detects an eye at coordinates $(10, 10)$, and in a new test image the subject is shifted so the eye appears at $(150, 150)$, the neuron at $(10, 10)$ fails to activate, and the neuron at $(150, 150)$ has never seen an eye during training! The network possesses zero ability to transfer learned patterns across spatial coordinates.
  \item \textbf{Exponential Sample Complexity}: Because an MLP cannot share parameters across space, it cannot generalize across translations. To learn that an object can appear anywhere on the canvas, the MLP must observe that object at \textbf{every possible spatial coordinate $(i, j)$ during training}. This multiplies the required training dataset size by a factor proportional to the spatial resolution $H \times W$, making sample-efficient visual learning impossible.
\end{itemize}

\paragraph{3. Catastrophic Parameter Explosion and Overfitting.}
The combination of all-to-all dense connectivity and zero parameter sharing inevitably triggers a catastrophic parameter explosion:
\begin{equation}
N_{\text{weights}} = d_{\text{in}} \times d_{\text{out}} = (H \cdot W \cdot C) \times M,
\end{equation}
where $M$ is the number of neurons in the first hidden layer.
\begin{itemize}[noitemsep]
  \item \textbf{Small Resolution (CIFAR-10)}: For a tiny $32 \times 32 \times 3$ thumbnail ($d = 3,072$), connecting to a hidden layer of $M = 1,024$ neurons requires $3,072 \times 1,024 \approx 3.15 \times 10^6$ parameters.
  \item \textbf{Medium Resolution ($256 \times 256 \times 3$)}: Here, $d = 196,608$. Connecting to a modest hidden layer of $M = 2,048$ neurons requires:
    \begin{equation}
    196,608 \times 2,048 = 402,653,184 \text{ parameters (over } 402 \text{ million weights!)}
    \end{equation}
    for the \textbf{first layer alone}! Storing this single weight matrix in 32-bit floating point precision consumes over $1.61\text{ GB}$ of VRAM.
  \item \textbf{High Resolution (Full HD $1080 \times 1920 \times 3$)}: $d \approx 6.22 \times 10^6$. A single hidden layer of $2,048$ neurons would require over $1.27 \times 10^{10}$ parameters (12.7 billion weights, $\sim 50\text{ GB}$ of memory).
\end{itemize}
By Vapnik-Chervonenkis (VC) generalization theory (Chapter~\ref{sec:pac_learning_theory}), the sample complexity required to prevent overfitting scales linearly with the number of free parameters ($N \propto \mathcal{O}(|\boldsymbol{\theta}|)$). Training an unconstrained model with hundreds of millions of parameters on typical dataset sizes ($N \sim 10^4$--$10^6$) guarantees catastrophic overfitting and failure to generalize.

\subsection{Natural image statistics: Why do all images share universal patterns?}
\label{sec:natural_image_statistics}

Why are dense networks structurally mismatched to images, and why can we design a far more compact, universal visual architecture? As Prof.~Tatiana Tommasi emphasized during lectures, the answer lies in the fundamental statistical physics of natural scenes.

\paragraph{1. Natural Images vs. White Noise.}
An arbitrary $256 \times 256$ RGB image can take $(256)^{256 \times 256 \times 3} \approx 10^{473,000}$ possible configurations. However, the set of \textbf{natural images} occupies an infinitesimal sub-manifold within this immense space. If pixels were drawn independently from a Gaussian white noise distribution, images would resemble uniform television static, containing zero discernible structures.

Natural images are profoundly non-random because they are optical projections of a \textbf{physical 3D universe}:
\begin{itemize}[noitemsep]
  \item Solid physical objects consist of continuous matter with uniform surface reflectance.
  \item Optical illumination casts smooth gradients across surfaces.
  \item Object occlusions and material boundaries produce sharp, localized step discontinuities.
\end{itemize}

\paragraph{2. Universal Statistical Symmetries of Natural Scenes.}
Empirical investigations into natural image statistics (Ruderman, Simoncelli, Olshausen \& Field) reveal that virtually all natural images share identical mathematical regularities:
\begin{enumerate}
  \item \textbf{High Local Autocorrelation}: The correlation between two pixels $\mathbf{I}(\mathbf{x})$ and $\mathbf{I}(\mathbf{x} + \Delta \mathbf{x})$ approaches $1$ as $\|\Delta \mathbf{x}\|_2 \to 0$. Adjacent pixels almost always possess identical color and luminance; significant differences occur only at boundaries.
  \item \textbf{Scale Invariance and the $1/f^2$ Power Spectrum}: The power spectral density $S(f)$ of natural scenes universally follows a power law:
    \begin{equation}
    S(f) \propto \frac{1}{f^{2 - \eta}}, \quad \text{with } \eta \approx 0.2.
    \label{eq:power_spectrum_law}
    \end{equation}
    This scale-invariant property implies that natural images exhibit identical statistical structure at all spatial scales: low spatial frequencies (broad smooth areas) dominate the energy budget, while high spatial frequencies (edges and fine textures) represent a sparse, highly informative fraction.
  \item \textbf{Universality of Edge Primitives}: Across entirely unrelated domains—whether analyzing portraits, satellite topography, medical histology, or urban streetscapes—the visual world is universally built from the same primitive visual tokens: \textbf{oriented edges, localized step transitions, and corner junctions}.
\end{enumerate}

Because all natural images share these identical low-level statistical regularities, an optimal visual machine learning model should not learn arbitrary, unconstrained dense weights. Instead, it must embody inductive biases that exploit locality, stationarity, and hierarchical compositionality!

\begin{figure}[H]
\centering
\begin{tikzpicture}[scale=0.95]
  % Stage 1: Pixels
  \node[rectangle, draw=gray!80, fill=gray!10, thick, inner sep=6pt, font=\footnotesize, align=center] (s1) at (0, 0) {
    \textbf{Level 1: Raw Pixels}\\
    Coordinate grid $\mathbf{I}(i, j) \in [0, 255]^3$
  };

  % Stage 2: Edges & Contours
  \node[rectangle, draw=deepblue, fill=deepblue!15, thick, inner sep=6pt, font=\footnotesize, align=center] (s2) at (3.5, 0) {
    \textbf{Level 2: Low-Level Primitives}\\
    Oriented edges, bars, Gabor wavelets
  };

  % Stage 3: Motifs & Textures
  \node[rectangle, draw=intuitiongreen, fill=intuitiongreen!15, thick, inner sep=6pt, font=\footnotesize, align=center] (s3) at (7.2, 0) {
    \textbf{Level 3: Mid-Level Motifs}\\
    Corners, junctions, textures, shapes
  };

  % Stage 4: Parts & Objects
  \node[rectangle, draw=politoNavy, fill=politoNavy!20, thick, inner sep=6pt, font=\footnotesize, align=center] (s4) at (11.0, 0) {
    \textbf{Level 4: High-Level Semantics}\\
    Object parts (eyes, wheels) $\to$ Full objects
  };

  % Directed Arrows between stages
  \draw[-{Stealth}, very thick, politoNavy] (s1) -- (s2);
  \draw[-{Stealth}, very thick, politoNavy] (s2) -- (s3);
  \draw[-{Stealth}, very thick, politoNavy] (s3) -- (s4);

  % Receptive Field expansion annotation below
  \draw[thick, deepdiveTeal, {Stealth}-{Stealth}] (0, -0.9) -- (11.0, -0.9);
  \node[font=\scriptsize\bfseries\color{deepdiveTeal}, below=2pt] at (5.5, -0.9) {Increasing Receptive Field \& Semantic Abstraction $\longrightarrow$};
\end{tikzpicture}
\caption{The hierarchical compositionality of visual representations. Natural images are formed compositionally: local pixels combine into oriented edges, edges aggregate into textures and geometric motifs, which assemble into object parts, and ultimately into semantic categories. Convolutional networks emulate this visual hierarchy across cascading layers.}
\label{fig:visual_feature_hierarchy}
\end{figure}

As depicted in Figure~\ref{fig:visual_feature_hierarchy}, neurophysiologists David Hubel and Torsten Wiesel~\cite{hubel1959receptive,hubel1962receptive} discovered in 1959 that the mammalian visual cortex (area V1) is organized precisely according to this compositional principle: \textbf{Simple Cells} respond to localized oriented bars within small receptive fields, while \textbf{Complex Cells} pool over simple cells to achieve local translation tolerance. 
This biological architecture inspired Kunihiko Fukushima's \textbf{Neocognitron} (1980)~\cite{fukushima1980neocognitron} and Yann LeCun's \textbf{LeNet} (1989)~\cite{lecun1989backpropagation}, establishing the mathematical paradigm of \textbf{Convolutional Neural Networks}.

---

\section{Convolutional neural networks: Mechanics, geometry, and design principles}
\label{sec:cnn_mechanics_and_geometry}

Having demonstrated the structural failure of dense MLPs and the statistical regularities of natural images, we now introduce the mathematical framework engineered to operationalize these visual priors: \textbf{Convolutional Neural Networks (CNNs)}.

By replacing unconstrained matrix multiplications with discrete cross-correlations, CNNs enforce three foundational physical priors directly into the network architecture:
\begin{enumerate}
  \item \textbf{Sparse Interactions (Locality)}: Rather than connecting each hidden unit to all $H \times W$ pixels, a convolutional filter interacts only with a small, localized spatial patch of size $F \times F$ (typically $3 \times 3$ or $5 \times 5$).
  \item \textbf{Parameter Sharing (Tied Weights)}: Rather than learning separate parameters for every spatial location, the identical filter kernel is reused across every receptive field in the visual grid.
  \item \textbf{Translation Equivariance}: Because the filter slides uniformly across space, translating the input image translates the output representation identically:
    \begin{equation}
    f\left( \mathcal{T}_{(\Delta x, \Delta y)}[\mathbf{X}] \right) = \mathcal{T}_{(\Delta x, \Delta y)}[f(\mathbf{X})].
    \label{eq:equivariance_def}
    \end{equation}
\end{enumerate}

\subsection{The continuous convolution and discrete 2D cross-correlation}
\label{sec:convolution_mathematical_definition}

In functional analysis, the \textbf{continuous convolution} of two functions $f, g: \R \to \R$ is defined by the integral:
\begin{equation}
(f * g)(t) \triangleq \int_{-\infty}^{\infty} f(\tau) g(t - \tau) \, d\tau.
\label{eq:continuous_convolution}
\end{equation}
In continuous signal processing, convolution is commutative ($f * g = g * f$), associative, and distributive. The term $g(t - \tau)$ implies that the kernel function $g$ is \textbf{flipped (reflected)} prior to shifting.

\paragraph{Discrete 2D Cross-Correlation in Deep Learning.}
In digital image processing, images are discrete 2D grids of pixel intensities $\mathbf{I} \in \R^{H \times W}$, and filters are small discrete 2D kernels $\mathbf{K} \in \R^{F \times F}$. In deep learning libraries (PyTorch, TensorFlow, JAX), the operation implemented under the name "convolution" is, strictly speaking, \textbf{discrete 2D cross-correlation} (convolution without kernel flipping)~\cite{bishop2024deep,goodfellow2016deep}:
\begin{equation}
S(i, j) = (\mathbf{I} * \mathbf{K})(i, j) \triangleq \sum_{m=0}^{F-1} \sum_{n=0}^{F-1} \mathbf{I}(i + m, \, j + n) \, \mathbf{K}(m, n).
\label{eq:discrete_cross_correlation}
\end{equation}
Because the kernel weights $\mathbf{K}(m, n)$ are learned parameters optimized directly via backpropagation from data, flipping the kernel mathematically would simply cause the optimizer to learn the mirrored weights with zero difference in representational power. Following machine learning convention, we refer to Equation~\eqref{eq:discrete_cross_correlation} as convolution.

\begin{figure}[H]
\centering
\begin{tikzpicture}[scale=0.92]
  % Top: Dense Connectivity
  \begin{scope}[shift={(0,2.8)}]
    \node[font=\footnotesize\bfseries\color{politoNavy}] at (2.5, 1.2) {(a) Dense Connectivity (Independent Weights: $5 \times 5 = 25$)};
    \foreach \i in {1,...,5} {
      \node[circle, draw=deepblue, fill=deepblue!15, inner sep=2pt, font=\scriptsize] (in_\i) at (\i, 0) {$x_\i$};
      \node[circle, draw=politoNavy, fill=politoNavy!15, inner sep=2pt, font=\scriptsize] (out_\i) at (\i, -1.5) {$y_\i$};
    }
    \foreach \i in {1,...,5} {
      \foreach \j in {1,...,5} {
        \draw[gray!30, thin] (in_\i) -- (out_\j);
      }
    }
  \end{scope}

  % Bottom: Convolutional Sparse Interactions & Tied Weights
  \begin{scope}[shift={(0,0)}]
    \node[font=\footnotesize\bfseries\color{intuitiongreen}] at (2.5, 0.9) {(b) Convolutional Sparse Interactions (Kernel Size $F=3$, Tied Weights: $3$)};
    \foreach \i in {1,...,5} {
      \node[circle, draw=deepblue, fill=deepblue!15, inner sep=2pt, font=\scriptsize] (cin_\i) at (\i, 0) {$x_\i$};
    }
    \foreach \j in {2,3,4} {
      \node[circle, draw=intuitiongreen, fill=intuitiongreen!20, very thick, inner sep=2pt, font=\scriptsize] (cout_\j) at (\j, -1.5) {$y_\j$};
    }

    % Shared weights color coded
    % Kernel: k_1 (red), k_2 (green), k_3 (blue)
    \foreach \j in {2,3,4} {
      \pgfmathtruncatemacro{\leftin}{\j-1}
      \pgfmathtruncatemacro{\midin}{\j}
      \pgfmathtruncatemacro{\rightin}{\j+1}
      \draw[-{Stealth}, very thick, examred] (cin_\leftin) -- (cout_\j);
      \draw[-{Stealth}, very thick, intuitiongreen] (cin_\midin) -- (cout_\j);
      \draw[-{Stealth}, very thick, politoNavy] (cin_\rightin) -- (cout_\j);
    }

    \node[font=\scriptsize, align=center] at (6.8, -0.7) {
      \textcolor{examred}{$\mathbf{w}_1$ (Left Tap)}\\
      \textcolor{intuitiongreen}{$\mathbf{w}_2$ (Center Tap)}\\
      \textcolor{politoNavy}{$\mathbf{w}_3$ (Right Tap)}
    };
  \end{scope}
\end{tikzpicture}
\caption{Sparse interactions and parameter sharing in convolutional layers compared to dense layers. (a) A dense layer connects every input unit to every output unit with $25$ independent weights. (b) A convolutional layer connects each output unit only to a local spatial receptive field of size $F=3$. The identical weights (color-coded red, green, blue) are reused across all spatial positions.}
\label{fig:sparse_vs_dense_connectivity}
\end{figure}

As shown in Figure~\ref{fig:sparse_vs_dense_connectivity}, discrete convolution mathematically enforces sparse connectivity and parameter sharing. The number of parameters is \textbf{completely independent of image resolution}, determined strictly by kernel size $F \times F \times C$.

\subsection{Translation equivariance and algebraic sample complexity advantages}
\label{sec:translation_equivariance_sample_complexity}

Having established discrete cross-correlation, we examine a foundational question emphasized in both theoretical computer science and oral examinations: \textit{why cannot a standard dense Multi-Layer Perceptron (MLP) simply learn to process images by training on large datasets, and what formal advantage does the algebraic inductive bias of convolution provide?}~\cite{prince2023understanding,goodfellow2016deep}

\paragraph{Mathematical Definitions: Invariance vs. Equivariance.}
Let $\mathcal{X} = \R^{H \times W}$ represent the space of 2D images. We define the 2D discrete translation operator $\mathcal{T}_{\boldsymbol{\delta}}$ parameterized by spatial displacement vector $\boldsymbol{\delta} = (\Delta x, \Delta y) \in \mathbb{Z}^2$:
\begin{equation}
(\mathcal{T}_{\boldsymbol{\delta}} \mathbf{X})_{i, j} \triangleq \mathbf{X}_{i - \Delta x, \, j - \Delta y}.
\label{eq:translation_operator_def}
\end{equation}
As formalized by Prince~\cite{prince2023understanding}, visual recognition tasks impose two distinct operational symmetries:
\begin{enumerate}[leftmargin=*]
  \item \textbf{Translation Invariance}: A mapping $f: \mathcal{X} \to \mathcal{Y}$ is translation-invariant if translating the input leaves the output completely unchanged:
    \begin{equation}
    f\left( \mathcal{T}_{\boldsymbol{\delta}} \mathbf{X} \right) = f(\mathbf{X}), \quad \forall \boldsymbol{\delta} \in \mathbb{Z}^2.
    \label{eq:translation_invariance_math}
    \end{equation}
    Translation invariance is the mathematically required property for whole-image categorical classification: the semantic label (e.g. \textit{``Cat''}) must remain identical regardless of whether the feline sits in the top-left or bottom-right corner of the canvas.
  
  \item \textbf{Translation Equivariance (Covariance)}: A mapping $f: \mathcal{X} \to \mathcal{Z}$ is translation-equivariant with respect to $\mathcal{T}_{\boldsymbol{\delta}}$ if transforming the input transforms the output identically, commuting with the translation operator:
    \begin{equation}
    f\left( \mathcal{T}_{\boldsymbol{\delta}} \mathbf{X} \right) = \mathcal{T}_{\boldsymbol{\delta}} \big[ f(\mathbf{X}) \big], \quad \forall \boldsymbol{\delta} \in \mathbb{Z}^2.
    \label{eq:translation_equivariance_math}
    \end{equation}
    Translation equivariance is the indispensable property for internal feature representations, dense spatial localization, and semantic segmentation: if an object shifts five pixels to the right, all extracted feature maps must shift by exactly five pixels to preserve spatial geometry.
\end{enumerate}

\paragraph{The Catastrophic Sample Complexity of Dense MLPs.}
Consider an MLP processing an image $\mathbf{X} \in \R^{H \times W}$ flattened into vector $\mathbf{x} = \operatorname{vec}(\mathbf{X}) \in \R^{D}$ with $D = H \cdot W$. A fully connected layer computes hidden activations via an unconstrained dense weight matrix $\mathbf{W} \in \R^{D_{\text{out}} \times D}$:
\begin{equation}
h_p = \sigma\left( \sum_{i=1}^H \sum_{j=1}^W W_{p, (i, j)} X_{i, j} + b_p \right).
\label{eq:mlp_dense_coordinate_dot}
\end{equation}
Notice that in Equation~\eqref{eq:mlp_dense_coordinate_dot}, every input coordinate $(i, j)$ is governed by a completely independent, decoupled parameter $W_{p, (i, j)}$.

Suppose a visual pattern $\mathcal{P}$ (e.g., an eye or a sharp vertical edge) is presented at spatial coordinate $(i_1, j_1)$ during training. Backpropagation calculates the gradient of the objective $\mathcal{L}$ with respect to the weights connected to that specific coordinate:
\begin{equation}
\frac{\partial \mathcal{L}}{\partial W_{p, (i_1, j_1)}} = \delta_p X_{i_1, j_1} \neq 0.
\end{equation}
However, for any other spatial coordinate $(i_2, j_2) \neq (i_1, j_1)$ where the pattern was not presented, the local input is zero or background, yielding:
\begin{equation}
\frac{\partial \mathcal{L}}{\partial W_{p, (i_2, j_2)}} = 0.
\end{equation}
The weights connected to coordinate $(i_2, j_2)$ receive \textbf{identically zero gradient updates}! Consequently, if the object appears at position $(i_2, j_2)$ at test time, the MLP possesses zero ability to recognize it.

To learn translation equivariance empirically, the MLP must observe every visual concept $\mathcal{P}$ instantiated at every coordinate in the spatial translation group $\mathcal{T} = \{(\Delta x, \Delta y) \mid 0 \le \Delta x < W, \, 0 \le \Delta y < H\}$. Because the cardinality of this group is $|\mathcal{T}| = H \times W$, the required training sample complexity scales as:
\begin{equation}
N_{\text{MLP}} = \Omega\left( |\mathcal{T}| \cdot N_0 \right) = \Omega\left( H \cdot W \cdot N_0 \right),
\label{eq:sample_complexity_mlp}
\end{equation}
where $N_0$ is the sample complexity required to learn the pattern at a single fixed location. For a standard $256 \times 256$ input image, $|\mathcal{T}| = 65,536$: an MLP requires up to a \textbf{$65,000\times$ larger dataset} purely to brute-force learn translation invariance through exhaustive spatial data augmentation!

\begin{figure}[H]
\centering
\begin{tikzpicture}[scale=0.90]
  % -------------------------------------------------------------
  % PANEL (a): MLP Failure & High Sample Complexity
  % -------------------------------------------------------------
  \node[font=\footnotesize\bfseries, text=examred, align=center] at (2.6, 3.6) {
    (a) Dense Multi-Layer Perceptron\\[2pt]
    {\scriptsize\normalfont\color{examred!90!black}Spatial Decoupling: $\Omega(H \cdot W \cdot N_0)$}
  };
  
  % Input Grid 1 (Pattern at Top-Left)
  \draw[fill=gray!12, draw=gray!50, rounded corners=1pt] (0.4, 1.2) rectangle (1.8, 2.6);
  \draw[step=0.35, gray!25, thin] (0.4, 1.2) grid (1.8, 2.6);
  \node[fill=examred!25, draw=examred, thick, rounded corners=1pt, inner sep=1.8pt, font=\tiny\bfseries\color{examred}] at (0.58, 2.42) {Obj};
  \node[font=\tiny, text=gray!80, anchor=north] at (1.1, 1.05) {Position $(i_1, j_1)$};

  % MLP Layer 1 (Hidden Neuron h1)
  \node[circle, draw=politoNavy, fill=politoNavy!15, thick, inner sep=2.5pt, font=\scriptsize\bfseries\color{politoNavy}] (h1) at (5.0, 1.9) {$h_1$};
  \draw[-{Stealth[scale=1.0]}, thick, examred] (1.85, 2.1) -- (h1) 
    node[midway, above=2pt, font=\tiny\bfseries, text=examred] {$\mathbf{w}_1$ (Learned)};

  % Input Grid 2 (Pattern Shifted to Bottom-Right)
  \draw[fill=gray!12, draw=gray!50, rounded corners=1pt] (0.4, -0.8) rectangle (1.8, 0.6);
  \draw[step=0.35, gray!25, thin] (0.4, -0.8) grid (1.8, 0.6);
  \node[fill=deepblue!25, draw=deepblue, thick, rounded corners=1pt, inner sep=1.8pt, font=\tiny\bfseries\color{deepblue}] at (1.62, -0.42) {Obj};
  \node[font=\tiny, text=gray!80, anchor=north] at (1.1, -0.95) {Position $(i_2, j_2)$};

  % MLP Layer 2 (Hidden Neuron h2)
  \node[circle, draw=politoNavy, fill=politoNavy!15, thick, inner sep=2.5pt, font=\scriptsize\bfseries\color{politoNavy}] (h2) at (5.0, -0.1) {$h_2$};
  \draw[-{Stealth[scale=1.0]}, thick, dashed, gray!60] (1.85, -0.3) -- (h2) 
    node[midway, above=2pt, font=\tiny\bfseries, text=gray!70] {$\mathbf{w}_2$ (Untrained!)};

  % Panel (a) Consequence Callout Box (Narrower with ample margins)
  \node[draw=examred!50, fill=examback, rounded corners=2pt, inner sep=3.5pt, align=center, font=\tiny, text=examred!90!black] at (2.6, -1.95) {
    $\implies$ Requires observing pattern at all $H \times W$ locations\\
    $\partial \mathcal{L} / \partial W_{p,(i_2, j_2)} = 0 \implies$ Zero transfer across space
  };

  % -------------------------------------------------------------
  % CENTRAL DIVIDER
  % -------------------------------------------------------------
  \draw[dashed, gray!40, thick] (5.9, -2.4) -- (5.9, 3.9);

  % -------------------------------------------------------------
  % PANEL (b): CNN Weight Tying & Exact Equivariance
  % -------------------------------------------------------------
  \node[font=\footnotesize\bfseries, text=intuitiongreen, align=center] at (9.4, 3.6) {
    (b) Convolutional Neural Network\\[2pt]
    {\scriptsize\normalfont\color{intuitiongreen!90!black}Shift Equivariance: $O(N_0)$ samples}
  };

  % Conv input with pattern at Top-Left
  \draw[fill=gray!12, draw=gray!50, rounded corners=1pt] (6.6, 1.2) rectangle (8.0, 2.6);
  \draw[step=0.35, gray!25, thin] (6.6, 1.2) grid (8.0, 2.6);
  \node[fill=examred!25, draw=examred, thick, rounded corners=1pt, inner sep=1.8pt, font=\tiny\bfseries\color{examred}] at (6.78, 2.42) {Obj};
  \node[font=\tiny, text=gray!80, anchor=north] at (7.3, 1.05) {Position $(i_1, j_1)$};

  % Conv output map 1
  \draw[fill=intuitiongreen!10, draw=intuitiongreen!80, rounded corners=1pt] (10.8, 1.2) rectangle (12.2, 2.6);
  \draw[step=0.35, intuitiongreen!20, thin] (10.8, 1.2) grid (12.2, 2.6);
  \node[fill=intuitiongreen!40, draw=intuitiongreen, thick, rounded corners=1pt, inner sep=1.8pt, font=\tiny\bfseries, text=intuitiongreen!90!black] at (10.98, 2.42) {Act};
  \node[font=\tiny, text=gray!80, anchor=north] at (11.5, 1.05) {Response at $(i_1, j_1)$};

  \draw[-{Stealth[scale=1.0]}, thick, intuitiongreen] (8.2, 1.9) -- (10.6, 1.9) 
    node[midway, above=2pt, font=\tiny\bfseries, text=intuitiongreen] {Kernel $\mathbf{K}$};

  % Conv input with pattern shifted to Bottom-Right
  \draw[fill=gray!12, draw=gray!50, rounded corners=1pt] (6.6, -0.8) rectangle (8.0, 0.6);
  \draw[step=0.35, gray!25, thin] (6.6, -0.8) grid (8.0, 0.6);
  \node[fill=deepblue!25, draw=deepblue, thick, rounded corners=1pt, inner sep=1.8pt, font=\tiny\bfseries\color{deepblue}] at (7.82, -0.42) {Obj};
  \node[font=\tiny, text=gray!80, anchor=north] at (7.3, -0.95) {Position $(i_2, j_2)$};

  % Conv output map 2
  \draw[fill=intuitiongreen!10, draw=intuitiongreen!80, rounded corners=1pt] (10.8, -0.8) rectangle (12.2, 0.6);
  \draw[step=0.35, intuitiongreen!20, thin] (10.8, -0.8) grid (12.2, 0.6);
  \node[fill=intuitiongreen!40, draw=intuitiongreen, thick, rounded corners=1pt, inner sep=1.8pt, font=\tiny\bfseries, text=intuitiongreen!90!black] at (12.02, -0.42) {Act};
  \node[font=\tiny, text=gray!80, anchor=north] at (11.5, -0.95) {Response at $(i_2, j_2)$};

  \draw[-{Stealth[scale=1.0]}, thick, intuitiongreen] (8.2, -0.1) -- (10.6, -0.1) 
    node[midway, above=2pt, font=\tiny\bfseries, text=intuitiongreen] {Same $\mathbf{K}$!};

  % Panel (b) Consequence Callout Box (Narrower with ample margins)
  \node[draw=intuitiongreen!50, fill=intuitionback, rounded corners=2pt, inner sep=3.5pt, align=center, font=\tiny, text=intuitiongreen!90!black] at (9.4, -1.95) {
    $\implies$ Exact translation equivariance: $f(\mathcal{T}_{\boldsymbol{\delta}}\mathbf{X}) = \mathcal{T}_{\boldsymbol{\delta}}[f(\mathbf{X})]$\\
    Sample complexity $O(N_0)$ independent of grid size
  };
\end{tikzpicture}
\caption{Sample complexity contrast between dense MLPs and Convolutional Networks under 2D spatial translation. (a) In an MLP, disjoint weights connect to each spatial pixel; learning a visual concept at position $(i_1, j_1)$ fails to generalize to position $(i_2, j_2)$. (b) In a CNN, parameter sharing ties all spatial positions to identical kernel weights $\mathbf{K}$, guaranteeing exact translation equivariance and sample complexity scaling independent of grid resolution.}
\label{fig:mlp_vs_cnn_sample_complexity}
\end{figure}

\paragraph{The CNN Algebraic Inductive Bias.}
As illustrated in Figure~\ref{fig:mlp_vs_cnn_sample_complexity}, a Convolutional Neural Network resolves this dilemma by hardcoding translation equivariance directly into its algebraic structure through \textbf{parameter sharing (tied weights)}. Convolving input $\mathbf{X}$ with shared kernel $\mathbf{K} \in \R^{F \times F}$ guarantees commutativity with translation by algebraic identity:
\begin{equation}
\begin{aligned}
\big[ (\mathcal{T}_{\boldsymbol{\delta}} \mathbf{X}) * \mathbf{K} \big](i, j) &= \sum_{m=0}^{F-1} \sum_{n=0}^{F-1} (\mathcal{T}_{\boldsymbol{\delta}} \mathbf{X})(i+m, \, j+n) \, K_{m, n} \\
&= \sum_{m=0}^{F-1} \sum_{n=0}^{F-1} \mathbf{X}(i - \Delta x + m, \, j - \Delta y + n) \, K_{m, n} \\
&= (\mathbf{X} * \mathbf{K})(i - \Delta x, \, j - \Delta y) = \mathcal{T}_{\boldsymbol{\delta}} \big[ (\mathbf{X} * \mathbf{K}) \big](i, j).
\end{aligned}
\label{eq:equivariance_proof}
\end{equation}
Because Equation~\eqref{eq:equivariance_proof} holds for \textit{any} kernel values $\mathbf{K}$, translation equivariance is strictly hardcoded prior to observing any training data. Observing a visual feature at a single location $(i_1, j_1)$ produces a gradient that updates the shared kernel $\mathbf{K}$; this single parameter update immediately generalizes detection across all $H \times W$ locations simultaneously:
\begin{equation}
N_{\text{CNN}} = O(N_0).
\label{eq:sample_complexity_cnn}
\end{equation}
As Simon Prince highlights~\cite{prince2023understanding}, weight sharing acts as an \textbf{infinite Bayesian prior} (or Dirac delta regularization) that penalizes all non-translation-equivariant linear mappings to infinity, slashing hypothesis complexity without sacrificing expressive power.

\subsection{Biological visual cortex foundations and spontaneous Gabor emergence}
\label{sec:biological_visual_cortex_and_gabor_emergence}

The hierarchical design of Convolutional Neural Networks is deeply rooted in 20th-century neurobiology, reflecting profound mathematical parallels between mammalian visual processing and statistical gradient optimization~\cite{bishop2024deep}.

\paragraph{1. The Hubel \& Wiesel Discovery (1959, 1962).}
In their Nobel-Prize-winning electrophysiological investigations, David Hubel and Torsten Wiesel recorded action potentials from individual cortical neurons in the feline striate cortex (primary visual cortex, Area 17 / V1) under controlled visual illumination~\cite{hubel1959receptive,hubel1962receptive}. They established that cortical neurons do not respond to uniform, diffuse illumination. Instead, V1 visual processing is organized hierarchically:
\begin{enumerate}[leftmargin=*]
  \item \textbf{Simple Cells}: Simple cells possess spatially localized receptive fields partitioned into antagonistic excitatory (ON) and inhibitory (OFF) sub-regions. They fire vigorously only in response to oriented visual stimuli (slits, bars, or sharp edges) oriented at a specific angle $\theta$ and situated at a precise spatial location $(x_0, y_0)$ in the visual field.
  \item \textbf{Complex Cells}: Deeper within the cortical lamina, complex cells receive convergent inputs from groups of simple cells sharing identical orientation preferences $\theta$ across adjacent spatial positions. Complex cells respond to oriented edges anywhere within a broader receptive field, displaying \textbf{local translation invariance} to small spatial shifts.
\end{enumerate}
This functional division---simple cells detecting localized oriented features followed by complex cells pooling across space---provided the direct biological archetype for Fukushima's \textit{Neocognitron} (1980)~\cite{fukushima1980neocognitron} and LeCun's \textit{LeNet} (1989, 1998)~\cite{lecun1989backpropagation,lecun1998gradient}: \textbf{convolutional layers computationally embody simple cells, while pooling layers embody complex cells}.

\paragraph{2. Mathematical Modeling: 2D Gabor Wavelets.}
Subsequent neurophysiological mapping (Daugman, 1985; Jones \& Palmer, 1987) demonstrated that the spatial receptive field response profile of a V1 simple cell is modeled with exceptional mathematical precision by a \textbf{2D Gabor filter}~\cite{bishop2024deep}: a 2D sinusoidal plane wave modulated by an anisotropic Gaussian decay envelope.

Following Bishop~\cite{bishop2024deep}, the continuous 2D Gabor function $G: \R^2 \to \R$ is defined by:
\begin{equation}
G(x, y) = A \exp\left( -\alpha \tilde{x}^2 - \beta \tilde{y}^2 \right) \sin(\omega \tilde{x} + \phi),
\label{eq:bishop_gabor_formula}
\end{equation}
where $(\tilde{x}, \tilde{y})$ represent coordinates rotated by orientation angle $\theta$ and translated to receptive field center $(x_0, y_0)$:
\begin{align}
\tilde{x} &= (x - x_0) \cos\theta + (y - y_0) \sin\theta, \label{eq:gabor_xtilde} \\
\tilde{y} &= -(x - x_0) \sin\theta + (y - y_0) \cos\theta. \label{eq:gabor_ytilde}
\end{align}
The analytical parameters in Equations~\eqref{eq:bishop_gabor_formula}--\eqref{eq:gabor_ytilde} govern physical feature tuning:
\begin{itemize}[noitemsep]
  \item $\theta \in [0, \pi)$: The preferred orientation angle of the edge detector.
  \item $\omega \in \R_+$: The spatial frequency of the sinusoidal oscillation, tuning the detector to specific edge scales.
  \item $\phi \in [0, 2\pi)$: The phase angle ($\phi = 0$ corresponds to an odd-symmetric, anti-phase edge detector; $\phi = \pi/2$ corresponds to an even-symmetric bar detector).
  \item $\alpha, \beta > 0$: Spatial decay constants determining the Gaussian standard deviations $\sigma_x = 1/\sqrt{2\alpha}$ (across stripes) and $\sigma_y = 1/\sqrt{2\beta}$ (along stripes).
\end{itemize}
In harmonic analysis, 2D Gabor functions possess a remarkable property: they uniquely attain the theoretical lower bound of the \textbf{Heisenberg-Gabor joint uncertainty relation} in space and spatial frequency:
\begin{equation}
\Delta x \cdot \Delta u \ge \frac{1}{4\pi}, \quad \Delta y \cdot \Delta v \ge \frac{1}{4\pi}.
\label{eq:heisenberg_gabor_uncertainty}
\end{equation}
They are the optimal mathematical operators for simultaneously localizing spatial position $(x, y)$ and orientation/frequency $(u, v)$.

\paragraph{3. Spontaneous Gabor Emergence via Backpropagation on Natural Images.}
When a deep convolutional network (such as AlexNet~\cite{krizhevsky2012imagenet} or VGG~\cite{simonyan2014very}) is trained on large-scale datasets (ImageNet) using standard stochastic gradient descent to minimize cross-entropy loss:
\begin{equation}
\min_{\mathbf{W}} \frac{1}{N} \sum_{i=1}^N \mathcal{L}_{\text{CE}}\big( f(\mathbf{X}_i; \mathbf{W}), y_i \big),
\end{equation}
a remarkable phenomenon emerges: \textbf{the first-layer convolutional filters spontaneously organize into oriented Gabor-like wavelets and color-opponent blob filters}~\cite{krizhevsky2012imagenet,bishop2024deep}.

Crucially, \textbf{no biological priors or Gabor templates were engineered into the network or the loss function}. Why does pure gradient descent autonomously rediscover mammalian visual cortex receptive fields?
\begin{itemize}[leftmargin=*]
  \item \textbf{Natural Image Statistics}: As established by Olshausen \& Field (1996)~\cite{olshausen1996emergence} and Bell \& Sejnowski (1997), natural images are not random Gaussian noise; their power spectra follow a universal scale-invariant power law $P(f) \propto 1/f^{2-\eta}$ governed by sharp physical occlusion boundaries. Seeking an optimal, non-redundant sparse code for natural image patches mathematically forces basis functions to converge to localized, oriented bandpass wavelets.
  \item \textbf{Color Opponency}: On 3-channel RGB imagery, early filters spontaneously split into achromatic (luminance) edge detectors and chromatic color-opponent filters: Red vs. Green ($R - G$) and Blue vs. Yellow ($B - (R+G)/2$). Gradient descent spontaneously rediscovers Hering's 1878 \textbf{opponent-process theory of biological color vision}.
\end{itemize}

As Christopher Bishop cautions~\cite{bishop2024deep}, this convergence does not prove that CNNs function as biophysical models of the mammalian brain; rather, it demonstrates that \textbf{Gabor-like filtering and color opponency are universal mathematical solutions dictated by the statistical physics of natural imagery}. Both biological evolution and artificial backpropagation independently discover the exact same optimal representations.

\subsection{Anatomy of a convolutional layer: Stride, padding, and output geometry}
\label{sec:conv_layer_anatomy}

We now formalize the full multi-channel convolutional layer as deployed in modern deep networks.

\paragraph{Input and Filter Tensor Volumes.}
\begin{itemize}
  \item The input is a 3D tensor volume $\mathbf{X} \in \R^{H_{\text{in}} \times W_{\text{in}} \times C_{\text{in}}}$, where $H_{\text{in}}, W_{\text{in}}$ denote spatial height and width, and $C_{\text{in}}$ denotes channel depth (e.g., $C_{\text{in}} = 3$ for RGB color images).
  \item The layer contains a bank of $C_{\text{out}}$ distinct convolutional filters.
  \item \textbf{Fundamental Depth Rule}: \textbf{Each filter kernel must span the ENTIRE depth $C_{\text{in}}$ of the input volume!} A filter has spatial dimensions $F \times F$, but volumetric dimensions:
    \begin{equation}
    \mathbf{K}^{(k)} \in \R^{F \times F \times C_{\text{in}}}, \quad \forall k \in \{1, \dots, C_{\text{out}}\}.
    \end{equation}
\end{itemize}

Each filter kernel $\mathbf{K}^{(k)}$ computes an element-wise dot product across all $C_{\text{in}}$ channels simultaneously over a local spatial window, adds a scalar bias $b_k \in \R$, and emits a single 2D spatial slice $\mathbf{Y}_{:,:,k} \in \R^{H_{\text{out}} \times W_{\text{out}}}$:
\begin{equation}
\mathbf{Y}_{i, j, k} = \sum_{c=1}^{C_{\text{in}}} \sum_{m=0}^{F-1} \sum_{n=0}^{F-1} \mathbf{X}_{i+m, \, j+n, \, c} \, \mathbf{K}^{(k)}_{m, n, c} + b_k.
\label{eq:multichannel_convolution}
\end{equation}
Stacking the resulting $C_{\text{out}}$ 2D feature maps produces the output volume $\mathbf{Y} \in \R^{H_{\text{out}} \times W_{\text{out}} \times C_{\text{out}}}$.

\begin{figure}[H]
\centering
\begin{tikzpicture}[scale=0.95]
  % 3D Input Volume: 32x32x3
  \draw[fill=deepblue!15, draw=deepblue, thick] (0, 0) -- (2.0, 0) -- (2.0, 2.0) -- (0, 2.0) -- cycle;
  \draw[fill=deepblue!25, draw=deepblue, thick] (2.0, 0) -- (2.6, 0.6) -- (2.6, 2.6) -- (2.0, 2.0) -- cycle;
  \draw[fill=deepblue!10, draw=deepblue, thick] (0, 2.0) -- (0.6, 2.6) -- (2.6, 2.6) -- (2.0, 2.0) -- cycle;
  \node[font=\footnotesize\bfseries\color{politoNavy}, align=center] at (1.0, 1.0) {Input Volume\\$32 \times 32 \times 3$};

  % Local filter spanning full depth
  \draw[fill=examred!30, draw=examred, thick] (0.5, 0.7) -- (1.1, 0.7) -- (1.1, 1.3) -- (0.5, 1.3) -- cycle;
  \draw[fill=examred!40, draw=examred, thick] (1.1, 0.7) -- (1.7, 1.3) -- (1.7, 1.9) -- (1.1, 1.3) -- cycle;
  \draw[fill=examred!20, draw=examred, thick] (0.5, 1.3) -- (1.1, 1.9) -- (1.7, 1.9) -- (1.1, 1.3) -- cycle;
  \node[font=\tiny\color{examred}, align=center] at (1.1, 1.3) {$5 \times 5$\\$\times 3$};

  % Filter bank arrow
  \draw[-{Stealth}, very thick, examred] (2.8, 1.3) -- (4.2, 1.3) node[midway, above, font=\scriptsize, align=center] {$6$ Filters\\$5 \times 5 \times 3$};

  % Output Volume: 28x28x6
  \draw[fill=intuitiongreen!15, draw=intuitiongreen, thick] (4.5, 0.2) -- (6.3, 0.2) -- (6.3, 2.0) -- (4.5, 2.0) -- cycle;
  \draw[fill=intuitiongreen!25, draw=intuitiongreen, thick] (6.3, 0.2) -- (7.5, 1.4) -- (7.5, 3.2) -- (6.3, 2.0) -- cycle;
  \draw[fill=intuitiongreen!10, draw=intuitiongreen, thick] (4.5, 2.0) -- (5.7, 3.2) -- (7.5, 3.2) -- (6.3, 2.0) -- cycle;
  \node[font=\footnotesize\bfseries\color{intuitiongreen}, align=center] at (5.4, 1.1) {Output Volume\\$28 \times 28 \times 6$};
\end{tikzpicture}
\caption{The volumetric geometry of multi-channel convolution. An input volume of $32 \times 32 \times 3$ is convolved by $C_{\text{out}} = 6$ distinct filters. Each individual filter spans the complete depth of the input volume ($5 \times 5 \times 3$), emitting a single 2D feature map ($28 \times 28$). Stacking the 6 resulting activation maps produces the output volume of size $28 \times 28 \times 6$.}
\label{fig:volumetric_convolution}
\end{figure}

\paragraph{Hyperparameters: Stride and Padding.}
The spatial dimensions $H_{\text{out}}, W_{\text{out}}$ are governed by three hyperparameters:
\begin{itemize}
  \item \textbf{Filter Size ($F$)}: The spatial extent of the kernel (typically $F \in \{3, 5, 7\}$).
  \item \textbf{Stride ($S$)}: The step size (in pixels) by which the filter shifts across the input grid. A stride of $S = 1$ shifts pixel-by-pixel; a stride of $S = 2$ skips every second pixel, downsampling spatial resolution by a factor of 2.
  \item \textbf{Zero-Padding ($P$)}: The number of zero-valued pixels appended to the perimeter of the input feature map.
\end{itemize}

\begin{theorem}[Universal Spatial Dimension Formula]
Given an input spatial dimension $N_{\text{in}}$ (height or width), filter size $F$, zero-padding $P$, and stride $S$, the output spatial dimension $N_{\text{out}}$ is given by:
\begin{equation}
N_{\text{out}} = \left\lfloor \frac{N_{\text{in}} - F + 2P}{S} \right\rfloor + 1.
\label{eq:spatial_dimension_formula}
\end{equation}
\end{theorem}

\paragraph{Walkthrough of Slide Examples: How Each Hyperparameter Governs Output Size.}
To understand how each hyperparameter component in Equation~\eqref{eq:spatial_dimension_formula} governs the output, consider the visual decomposition in Figure~\ref{fig:conv_hyperparameter_walkthrough}.

\begin{figure}[H]
\centering
\begin{tikzpicture}[scale=0.88]
  % Case 1: N_in = 7, F = 3, P = 0, S = 1 => N_out = 5
  \begin{scope}[shift={(0,3.6)}]
    \node[font=\footnotesize\bfseries\color{politoNavy}] at (3.5, 1.2) {Case 1: $N_{\text{in}}=7, F=3, P=0, S=1 \implies N_{\text{out}} = \frac{7-3+0}{1} + 1 = 5$};
    \foreach \x in {0,...,6} {
      \draw[fill=deepblue!15, draw=deepblue, thick] (\x, 0) rectangle (\x+0.8, 0.8);
      \node[font=\tiny] at (\x+0.4, 0.4) {$\x$};
    }
    % Filter highlight
    \draw[draw=examred, ultra thick, fill=examred!20, opacity=0.6] (0, 0) rectangle (2.8, 0.8);
    \draw[-{Stealth}, thick, examred] (0.4, -0.3) -- (1.4, -0.3) node[midway, below, font=\tiny] {$S=1$};
  \end{scope}

  % Case 2: N_in = 7, F = 3, P = 0, S = 2 => N_out = 3
  \begin{scope}[shift={(0,1.8)}]
    \node[font=\footnotesize\bfseries\color{politoNavy}] at (3.5, 1.2) {Case 2: $N_{\text{in}}=7, F=3, P=0, S=2 \implies N_{\text{out}} = \frac{7-3+0}{2} + 1 = 3$};
    \foreach \x in {0,...,6} {
      \draw[fill=deepblue!15, draw=deepblue, thick] (\x, 0) rectangle (\x+0.8, 0.8);
      \node[font=\tiny] at (\x+0.4, 0.4) {$\x$};
    }
    % Filter highlight step 1
    \draw[draw=examred, ultra thick, fill=examred!20, opacity=0.6] (0, 0) rectangle (2.8, 0.8);
    \draw[-{Stealth}, thick, examred] (0.4, -0.3) -- (2.4, -0.3) node[midway, below, font=\tiny] {$S=2$ (Skips 1 pixel)};
  \end{scope}

  % Case 3: N_in = 7, F = 3, P = 1, S = 1 => "Same" Padding => N_out = 7
  \begin{scope}[shift={(0,0)}]
    \node[font=\footnotesize\bfseries\color{intuitiongreen}] at (3.5, 1.2) {Case 3: $N_{\text{in}}=7, F=3, P=1, S=1 \implies N_{\text{out}} = \frac{7-3+2(1)}{1} + 1 = 7$ ("Same" Padding)};
    % Padding cells (dashed)
    \draw[fill=gray!20, draw=gray!70, dashed, thick] (-1.0, 0) rectangle (-0.2, 0.8);
    \node[font=\tiny\color{gray!70}] at (-0.6, 0.4) {$0$};
    \foreach \x in {0,...,6} {
      \draw[fill=deepblue!15, draw=deepblue, thick] (\x, 0) rectangle (\x+0.8, 0.8);
      \node[font=\tiny] at (\x+0.4, 0.4) {$\x$};
    }
    \draw[fill=gray!20, draw=gray!70, dashed, thick] (7.0, 0) rectangle (7.8, 0.8);
    \node[font=\tiny\color{gray!70}] at (7.4, 0.4) {$0$};

    % Filter over padded border
    \draw[draw=intuitiongreen, ultra thick, fill=intuitiongreen!20, opacity=0.6] (-1.0, 0) rectangle (1.8, 0.8);
  \end{scope}
\end{tikzpicture}
\caption{Visual walkthrough of convolutional spatial dimension mechanics. Top: A $7 \times 7$ grid with $3 \times 3$ filter and stride 1 produces $5 \times 5$. Middle: Stride 2 shifts the kernel by two pixels at a time, skipping every second location to yield $3 \times 3$. Bottom: Adding zero-padding $P = (F-1)/2 = 1$ appends a boundary of zeros, perfectly preserving the original spatial dimension $7 \times 7$ ("Same" padding).}
\label{fig:conv_hyperparameter_walkthrough}
\end{figure}

The mathematical walkthroughs presented by Prof.~Tommasi demonstrate:
\begin{enumerate}
  \item \textbf{Input $7 \times 7$, Filter $3 \times 3$, Stride $S=1$, Padding $P=0$}:
    \begin{equation}
    N_{\text{out}} = \frac{7 - 3 + 2(0)}{1} + 1 = \frac{4}{1} + 1 = 5 \implies \text{Output is } 5 \times 5.
    \end{equation}
  \item \textbf{Input $7 \times 7$, Filter $3 \times 3$, Stride $S=2$, Padding $P=0$}:
    \begin{equation}
    N_{\text{out}} = \frac{7 - 3 + 2(0)}{2} + 1 = \frac{4}{2} + 1 = 3 \implies \text{Output is } 3 \times 3.
    \end{equation}
  \item \textbf{Input $7 \times 7$, Filter $3 \times 3$, Stride $S=3$, Padding $P=0$}:
    \begin{equation}
    \frac{7 - 3 + 0}{3} + 1 = \frac{4}{3} + 1 = 2.33 \implies \text{Fractional! The filter does NOT fit without clipping!}
    \end{equation}
    In valid convolutional designs, hyperparameters must satisfy $(N_{\text{in}} - F + 2P) \pmod S = 0$.
  \item \textbf{Zero-Padding: "Same" vs. "Valid" Padding}:
    Without padding ($P = 0$, "valid" padding), applying an $F \times F$ convolution inevitably shrinks the spatial resolution by $(F - 1)$ pixels at each layer. In deep networks with dozens of layers, the feature map would collapse to $1 \times 1$ within a few layers, destroying spatial context.
    To preserve spatial resolution ($N_{\text{out}} = N_{\text{in}}$ with stride $S = 1$, known as \textbf{"Same" padding}), we set:
    \begin{equation}
    N_{\text{in}} = \frac{N_{\text{in}} - F + 2P}{1} + 1 \implies 2P = F - 1 \implies P = \frac{F - 1}{2}.
    \label{eq:same_padding_formula}
    \end{equation}
    For standard $3 \times 3$ filters ($F = 3$), setting $P = (3-1)/2 = 1$ perfectly maintains spatial dimensions. For $5 \times 5$ filters ($F = 5$), setting $P = 2$ preserves dimensions.
\end{enumerate}

\paragraph{Parameter and FLOP Counting.}
Let a convolutional layer process input $\mathbf{X} \in \R^{H \times W \times C_{\text{in}}}$ with $C_{\text{out}}$ filters of size $F \times F \times C_{\text{in}}$, stride $S$, and padding $P$:
\begin{itemize}
  \item \textbf{Total Learnable Parameters}: Each filter contains $F \times F \times C_{\text{in}}$ kernel weights plus 1 scalar bias. For $C_{\text{out}}$ filters:
    \begin{equation}
    N_{\text{params}} = \left( F \times F \times C_{\text{in}} + 1 \right) \times C_{\text{out}}.
    \label{eq:conv_param_count}
    \end{equation}
    Notice that $N_{\text{params}}$ is completely independent of the spatial size $H, W$!
  \item \textbf{Total Floating-Point Multiply-Accumulate Operations (FLOPs)}: Computing each scalar output activation requires $(F \times F \times C_{\text{in}})$ multiplications and additions. Emitting the full volume $H_{\text{out}} \times W_{\text{out}} \times C_{\text{out}}$ requires:
    \begin{equation}
    N_{\text{FLOPs}} = 2 \times H_{\text{out}} \times W_{\text{out}} \times \left( F \times F \times C_{\text{in}} \right) \times C_{\text{out}}.
    \label{eq:conv_flops_count}
    \end{equation}
\end{itemize}

\begin{example}[Slide Parameter Walkthrough]
Let input volume be $32 \times 32 \times 3$. We apply $C_{\text{out}} = 10$ filters of size $5 \times 5$, with stride $S = 1$ and padding $P = 2$:
\begin{itemize}[noitemsep]
  \item Spatial Output: $H_{\text{out}} = \frac{32 - 5 + 2(2)}{1} + 1 = 32 \implies \text{Output volume is } 32 \times 32 \times 10$.
  \item Parameters: $(5 \times 5 \times 3 + 1) \times 10 = (75 + 1) \times 10 = 760 \text{ parameters}$.
\end{itemize}
An MLP performing the identical mapping from $3,072$ inputs to $32 \times 32 \times 10 = 10,240$ outputs would require $3,072 \times 10,240 \approx 31.4 \text{ million parameters}$! Convolution achieves a \textbf{$41,000\times$ reduction in parameters} while achieving vastly superior visual generalization.
\end{example}

\subsection{The VGG principle: effective receptive fields and the algebra of stacking small kernels}
\label{sec:vgg_principle_stacking_small_kernels}

Having formalized the output geometry of individual convolutional layers, we analyze a decisive architectural paradigm shift in deep learning: \textit{should deep networks employ large convolutional kernels (such as the $11 \times 11$ and $5 \times 5$ filters in AlexNet~\cite{krizhevsky2012imagenet}), or should they cascade multiple small $3 \times 3$ filters?}~\cite{simonyan2014very,bishop2024deep}

This question was definitively resolved by Karen Simonyan and Andrew Zisserman in the landmark \textbf{VGG Network} (2014)~\cite{simonyan2014very}, establishing what is now universally known as the \textbf{VGG Principle}: \textit{a stack of multiple small $3 \times 3$ convolutional layers strictly dominates a single large convolutional layer in parameter efficiency, computational FLOPs, and non-linear expressive capacity}.

\paragraph{1. Effective Receptive Field ($ERF$) Arithmetic.}
The \textbf{receptive field} $R_l$ of a neuron in layer $l$ is defined as the spatial extent (diameter along one axis) in the input image $\mathbf{X}^{(0)}$ that contributes to the activation of that neuron~\cite{bishop2024deep,prince2023understanding}.

Let $k_l$ denote the kernel size at layer $l$, and let $s_l$ denote the layer stride. With $R_0 = 1$ (the input pixel itself), the cumulative stride (jump) preceding layer $l$ is $J_{l-1} = \prod_{i=1}^{l-1} s_i$ (with $J_0 = 1$). The effective receptive field progression follows the universal linear recurrence:
\begin{equation}
R_l = R_{l-1} + (k_l - 1) \cdot J_{l-1}.
\label{eq:receptive_field_recurrence}
\end{equation}
In standard convolutional blocks prior to pooling, the stride is unit ($s_l = 1, \forall l \implies J_{l-1} = 1$). Equation~\eqref{eq:receptive_field_recurrence} simplifies to:
\begin{equation}
R_l = R_{l-1} + (k_l - 1).
\label{eq:unit_stride_receptive_field}
\end{equation}
For a sequence of $L$ cascaded convolutional layers with identical kernel size $k$:
\begin{equation}
R_L = 1 + \sum_{l=1}^L (k - 1) = 1 + L(k - 1).
\label{eq:cascaded_receptive_field_closed_form}
\end{equation}
Specializing Equation~\eqref{eq:cascaded_receptive_field_closed_form} to standard $3 \times 3$ kernels ($k = 3 \implies k - 1 = 2$):
\begin{itemize}[leftmargin=*]
  \item \textbf{Layer 1 ($L=1$)}: $R_1 = 1 + 1(2) = \mathbf{3}$ ($3 \times 3$ effective receptive field).
  \item \textbf{Layer 2 ($L=2$)}: $R_2 = 3 + 2 = \mathbf{5}$ ($5 \times 5$ effective receptive field).  
    \textit{Analytical Proof}: A hidden unit $h_j^{(2)}$ at layer 2 computes a linear combination of three layer-1 units: $\{h_{j-1}^{(1)}, h_j^{(1)}, h_{j+1}^{(1)}\}$. In turn, the leftmost unit $h_{j-1}^{(1)}$ accesses input pixels $\{x_{j-2}, x_{j-1}, x_j\}$, while the rightmost unit $h_{j+1}^{(1)}$ accesses $\{x_j, x_{j+1}, x_{j+2}\}$. The total spatial span on the input is $[j-2, j+2]$, spanning $(j+2) - (j-2) + 1 = \mathbf{5}$ pixels along each dimension.
  \item \textbf{Layer 3 ($L=3$)}: $R_3 = 5 + 2 = \mathbf{7}$ ($7 \times 7$ effective receptive field).  
    \textit{Analytical Proof}: The three layer-2 units span $[j-3, j+3]$ on the input grid, encompassing $(j+3) - (j-3) + 1 = \mathbf{7}$ pixels along each dimension.
\end{itemize}

\begin{figure}[H]
\centering
\resizebox{0.95\textwidth}{!}{%
\begin{tikzpicture}[every node/.style={transform shape}]
  % Layer 0 (Input)
  \foreach \i in {1,...,7} {
    \node[circle, draw=politoNavy, fill=politoNavy!15, thick, inner sep=2.5pt] (x\i) at (\i*1.1, 0) {\scriptsize $x_{\i}$};
  }
  \node[font=\footnotesize\bfseries\color{politoNavy}, left=0.25cm of x1] {Input $\mathbf{X}$ ($7\times 7$)};

  % Layer 1 (Conv 1: 3x3)
  \foreach \i in {1,...,5} {
    \node[circle, draw=deepblue, fill=skyBlue!35, thick, inner sep=2.5pt] (h1_\i) at (\i*1.1+1.1, 1.4) {\scriptsize $h^{(1)}_{\i}$};
  }
  \node[font=\footnotesize\bfseries\color{deepblue}, left=0.25cm of h1_1] {Conv 1 ($3\times 3$)};

  % Connections Layer 0 -> Layer 1
  \foreach \i in {1,...,5} {
    \pgfmathtruncatemacro{\iA}{\i}
    \pgfmathtruncatemacro{\iB}{\i+1}
    \pgfmathtruncatemacro{\iC}{\i+2}
    \draw[-{Stealth[scale=0.7]}, gray!40] (x\iA) -- (h1_\i);
    \draw[-{Stealth[scale=0.7]}, gray!40] (x\iB) -- (h1_\i);
    \draw[-{Stealth[scale=0.7]}, gray!40] (x\iC) -- (h1_\i);
  }

  % Highlight central unit in Layer 1
  \draw[-{Stealth[scale=0.85]}, deepblue, very thick] (x3) -- (h1_3);
  \draw[-{Stealth[scale=0.85]}, deepblue, very thick] (x4) -- (h1_3);
  \draw[-{Stealth[scale=0.85]}, deepblue, very thick] (x5) -- (h1_3);

  % Layer 2 (Conv 2: 3x3)
  \foreach \i in {1,...,3} {
    \node[circle, draw=intuitiongreen, fill=intuitiongreen!25, thick, inner sep=2.5pt] (h2_\i) at (\i*1.1+2.2, 2.8) {\scriptsize $h^{(2)}_{\i}$};
  }
  \node[font=\footnotesize\bfseries\color{intuitiongreen}, left=0.25cm of h2_1] {Conv 2 ($3\times 3$)};

  % Connections Layer 1 -> Layer 2
  \foreach \i in {1,...,3} {
    \pgfmathtruncatemacro{\iA}{\i}
    \pgfmathtruncatemacro{\iB}{\i+1}
    \pgfmathtruncatemacro{\iC}{\i+2}
    \draw[-{Stealth[scale=0.7]}, gray!40] (h1_\iA) -- (h2_\i);
    \draw[-{Stealth[scale=0.7]}, gray!40] (h1_\iB) -- (h2_\i);
    \draw[-{Stealth[scale=0.7]}, gray!40] (h1_\iC) -- (h2_\i);
  }

  % Highlight central unit in Layer 2
  \draw[-{Stealth[scale=0.85]}, intuitiongreen, very thick] (h1_2) -- (h2_2);
  \draw[-{Stealth[scale=0.85]}, intuitiongreen, very thick] (h1_3) -- (h2_2);
  \draw[-{Stealth[scale=0.85]}, intuitiongreen, very thick] (h1_4) -- (h2_2);

  % Layer 3 (Conv 3: 3x3)
  \node[circle, draw=examred, fill=examred!25, very thick, inner sep=3pt] (h3_1) at (4.4, 4.2) {\scriptsize $h^{(3)}_1$};
  \node[font=\footnotesize\bfseries\color{examred}, left=0.25cm of h3_1] {Conv 3 ($3\times 3$)};

  % Connections Layer 2 -> Layer 3
  \draw[-{Stealth[scale=0.85]}, examred, very thick] (h2_1) -- (h3_1);
  \draw[-{Stealth[scale=0.85]}, examred, very thick] (h2_2) -- (h3_1);
  \draw[-{Stealth[scale=0.85]}, examred, very thick] (h2_3) -- (h3_1);

  % Braces underneath showing Receptive Fields
  \draw[thick, deepblue, decorate, decoration={brace, amplitude=4pt, mirror}] (3.1, -0.4) -- (5.7, -0.4)
    node[midway, below=4pt, font=\scriptsize\bfseries\color{deepblue}] {$\text{RF}_1 = 3\times 3$};
  \draw[thick, intuitiongreen, decorate, decoration={brace, amplitude=4pt, mirror}] (2.0, -1.0) -- (6.8, -1.0)
    node[midway, below=4pt, font=\scriptsize\bfseries\color{intuitiongreen}] {$\text{RF}_2 = 5\times 5$};
  \draw[thick, examred, decorate, decoration={brace, amplitude=5pt, mirror}] (0.9, -1.6) -- (7.9, -1.6)
    node[midway, below=5pt, font=\scriptsize\bfseries\color{examred}] {Effective $\text{RF}_3 = 7\times 7$};

  % Summary card on the right
  \node[draw=politoNavy, fill=deepback, thick, rounded corners=4pt, inner sep=6pt, anchor=west] at (9.1, 2.1) {
    \begin{minipage}{5.8cm}
      \textbf{\textcolor{politoNavy}{The VGG Stacking Principle}}\\[3pt]
      \scriptsize
      \textbf{Parameter Savings}:
      \begin{itemize}[leftmargin=*, noitemsep]
        \item Two $3\times 3$ vs. one $5\times 5$:
        $2 \cdot 9C^2 = 18C^2$ vs. $25C^2$ (\textbf{28\% cut})
        \item Three $3\times 3$ vs. one $7\times 7$:
        $3 \cdot 9C^2 = 27C^2$ vs. $49C^2$ (\textbf{45\% cut})
      \end{itemize}
      \vspace{2pt}
      \textbf{Representational Power}:\\
      Stacking incorporates $3$ non-linear activations (e.g. ReLUs) instead of $1$, forming more discriminative features.
    \end{minipage}
  };
\end{tikzpicture}%
}
\caption{Effective receptive field expansion of stacked $3\times 3$ convolutional kernels. Cascading two $3\times 3$ filters achieves a $5\times 5$ receptive field with $28\%$ parameter reduction, while cascading three $3\times 3$ filters spans a $7\times 7$ receptive field with $45\%$ parameter savings and three successive non-linearities.}
\label{fig:receptive_field_expansion}
\end{figure}

\paragraph{2. Derivation of Parameter Complexity Savings.}
As proved in Figure~\ref{fig:receptive_field_expansion}, cascading small convolutions achieves identical receptive fields to large kernels. We now quantify the exact parameter savings across a block operating on $C$ input and output channels ($C_{\text{in}} = C_{\text{out}} = C$):
\begin{enumerate}[leftmargin=*]
  \item \textbf{Stacking Two $3 \times 3$ vs. One $5 \times 5$}:
    \begin{itemize}[noitemsep]
      \item A single $5 \times 5$ layer has weights:
        \begin{equation}
        N_{\text{single}}^{5\times 5} = 1 \times \left( 5^2 \cdot C^2 \right) = \mathbf{25 C^2}.
        \end{equation}
      \item Two stacked $3 \times 3$ layers have weights:
        \begin{equation}
        N_{\text{stacked}}^{2 \times (3\times 3)} = 2 \times \left( 3^2 \cdot C^2 \right) = 2 \times 9 C^2 = \mathbf{18 C^2}.
        \end{equation}
      \item The parameter reduction is exact:
        \begin{equation}
        \Delta_{\text{params}} = \frac{25 C^2 - 18 C^2}{25 C^2} = \frac{7}{25} = \mathbf{28\% \text{ reduction}}.
        \label{eq:vgg_savings_28}
        \end{equation}
    \end{itemize}
  
  \item \textbf{Stacking Three $3 \times 3$ vs. One $7 \times 7$}:
    \begin{itemize}[noitemsep]
      \item A single $7 \times 7$ layer has weights:
        \begin{equation}
        N_{\text{single}}^{7\times 7} = 1 \times \left( 7^2 \cdot C^2 \right) = \mathbf{49 C^2}.
        \end{equation}
      \item Three stacked $3 \times 3$ layers have weights:
        \begin{equation}
        N_{\text{stacked}}^{3 \times (3\times 3)} = 3 \times \left( 3^2 \cdot C^2 \right) = 3 \times 9 C^2 = \mathbf{27 C^2}.
        \end{equation}
      \item The parameter reduction is exact:
        \begin{equation}
        \Delta_{\text{params}} = \frac{49 C^2 - 27 C^2}{49 C^2} = \frac{22}{49} \approx \mathbf{44.898\% \approx 45\% \text{ reduction}}.
        \label{eq:vgg_savings_45}
        \end{equation}
    \end{itemize}
\end{enumerate}

\paragraph{3. Computational FLOPs Savings.}
For an activation feature map of spatial resolution $H \times W$ with unit stride and padding to preserve dimensions:
\begin{itemize}[leftmargin=*]
  \item A single $7 \times 7$ layer requires $2 \times (49 C^2 H W) = 98 C^2 H W$ FLOPs.
  \item Three stacked $3 \times 3$ layers require $3 \times \big( 2 \times 9 C^2 H W \big) = 54 C^2 H W$ FLOPs.
  \item The computational cost drops by $\frac{98 - 54}{98} = \mathbf{45\%}$!
\end{itemize}

\paragraph{4. Representational Gain: Non-Linear Depth vs. Linear Collapse.}
Beyond parameter and FLOP reductions, cascading smaller filters confers a crucial \textbf{representational expressiveness gain}:
\begin{itemize}[leftmargin=*]
  \item \textbf{Linear Collapse in the Absence of Activations}: If we applied three linear convolutional filters without intermediate activation functions, the associative property of convolution collapses the entire cascade into a single linear filter:
    \begin{equation}
    \mathbf{W}_3 * (\mathbf{W}_2 * (\mathbf{W}_1 * \mathbf{X})) = (\mathbf{W}_3 * \mathbf{W}_2 * \mathbf{W}_1) * \mathbf{X} = \mathbf{W}_{\text{eff}} * \mathbf{X},
    \end{equation}
    where $\mathbf{W}_{\text{eff}} \in \R^{7 \times 7}$. The cascade would simply parameterize a constrained subspace of linear $7 \times 7$ kernels.
  \item \textbf{Discriminative Expressiveness via Cascaded ReLUs}: In contrast, by placing a non-linear activation function $\sigma(\cdot)$ (e.g., ReLU) after each $3 \times 3$ convolution:
    \begin{align}
    \mathbf{H}_1 &= \sigma\left(\mathbf{W}_1 * \mathbf{X} + \mathbf{b}_1\right), \\
    \mathbf{H}_2 &= \sigma\left(\mathbf{W}_2 * \mathbf{H}_1 + \mathbf{b}_2\right), \\
    \mathbf{Y} &= \sigma\left(\mathbf{W}_3 * \mathbf{H}_2 + \mathbf{b}_3\right),
    \end{align}
    the network incorporates \textbf{three non-linear decision stages} instead of one. Three cascaded ReLUs fold and partition the local input space into an exponentially larger number of piecewise-linear decision regions~\cite{montufar2014number,telgarsky2016benefits}. This enables the network to learn substantially more discriminative, hierarchical visual primitives (e.g. oriented edges $\to$ corners $\to$ textured motifs) while simultaneously regularizing the hypothesis space against unstructured spatial overfitting.
\end{itemize}

\subsection{Network In Network and the power of $1 \times 1$ convolutions}
\label{sec:one_by_one_convolutions}

At first glance, a $1 \times 1$ convolutional filter ($F = 1$) might appear redundant: convolving a scalar with a single pixel seems to merely scale pixel values. 

However, as introduced by Min Lin et al. (2013) in \textit{Network In Network}~\cite{lin2013network} and shown in Figure~\ref{fig:one_by_one_conv}, because filters span the entire channel depth $C_{\text{in}}$, a $1 \times 1$ filter is actually a \textbf{cross-channel linear projection}:
\begin{equation}
\mathbf{Y}_{i, j, k} = \sum_{c=1}^{C_{\text{in}}} \mathbf{X}_{i, j, c} \, K_{0, 0, c}^{(k)} + b_k.
\label{eq:one_by_one_math}
\end{equation}

\begin{figure}[H]
\centering
\begin{tikzpicture}[scale=0.92]
  % Input 3D Volume: H x W x C_in
  \draw[fill=deepblue!15, draw=deepblue, thick] (0, 0) rectangle (2.2, 2.2);
  \draw[fill=deepblue!25, draw=deepblue, thick] (2.2, 0) -- (3.2, 0.8) -- (3.2, 3.0) -- (2.2, 2.2) -- cycle;
  \draw[fill=deepblue!10, draw=deepblue, thick] (0, 2.2) -- (1.0, 3.0) -- (3.2, 3.0) -- (2.2, 2.2) -- cycle;
  \node[font=\footnotesize\bfseries\color{politoNavy}, align=center] at (1.1, 1.1) {Input Volume\\$H \times W \times C_{\text{in}}$};

  % Pencil-like 1x1xC_in filter at (i, j)
  \draw[fill=examred!40, draw=examred, very thick] (0.8, 0.8) -- (0.8, 1.0) -- (1.8, 1.8) -- (1.8, 1.6) -- cycle;
  \node[font=\tiny\color{examred}, rotate=38] at (1.4, 1.5) {$1 \times 1 \times C_{\text{in}}$};

  % Projection Arrow
  \draw[-{Stealth}, ultra thick, examred] (3.6, 1.5) -- (5.0, 1.5) node[midway, above, font=\scriptsize, align=center] {$C_{\text{out}}$ Filters\\$1 \times 1 \times C_{\text{in}}$};

  % Output 3D Volume: H x W x C_out
  \draw[fill=intuitiongreen!15, draw=intuitiongreen, thick] (5.4, 0) rectangle (7.6, 2.2);
  \draw[fill=intuitiongreen!25, draw=intuitiongreen, thick] (7.6, 0) -- (8.2, 0.6) -- (8.2, 2.8) -- (7.6, 2.2) -- cycle;
  \draw[fill=intuitiongreen!10, draw=intuitiongreen, thick] (5.4, 2.2) -- (6.0, 2.8) -- (8.2, 2.8) -- (7.6, 2.2) -- cycle;
  \node[font=\footnotesize\bfseries\color{intuitiongreen}, align=center] at (6.5, 1.1) {Output Volume\\$H \times W \times C_{\text{out}}$};

  % Single pixel column output
  \draw[fill=intuitiongreen!40, draw=intuitiongreen, very thick] (6.2, 0.8) -- (6.2, 1.0) -- (6.8, 1.6) -- (6.8, 1.4) -- cycle;
\end{tikzpicture}
\caption{Mechanics of a $1 \times 1$ convolution. A $1 \times 1 \times C_{\text{in}}$ filter spans the entire channel depth at a single spatial pixel coordinate $(i, j)$, computing a localized linear combination across all channels without altering spatial height or width. This enables cross-channel dimensionality reduction (bottlenecks) and non-linear feature pooling.}
\label{fig:one_by_one_conv}
\end{figure}

A $1 \times 1$ convolution operates as a multi-layer perceptron applied identically to every spatial pixel coordinate $(i, j)$:
\begin{enumerate}
  \item \textbf{Channel Dimensionality Reduction / Bottlenecks}: If an input has $C_{\text{in}} = 256$ channels, applying $C_{\text{out}} = 64$ filters of size $1 \times 1 \times 256$ reduces the channel depth from 256 to 64 while leaving spatial dimensions $(H, W)$ completely unchanged. This bottleneck concept is the foundational computational engine of Inception networks and ResNet bottleneck blocks~\cite{he2016deep}, reducing FLOPs by up to $10\times$.
  \item \textbf{Channel Dimensionality Expansion}: Conversely, setting $C_{\text{out}} > C_{\text{in}}$ lifts feature representations into a richer, higher-dimensional channel space.
  \item \textbf{Cross-Channel Non-Linearity}: Placing a ReLU activation after a $1 \times 1$ convolution allows the network to learn non-linear interactions across feature channels at zero spatial cost.
\end{enumerate}

\subsection{Pooling layers and global decision synthesis: dense heads vs. global average pooling}
\label{sec:pooling_and_decision_synthesis}
\label{sec:pooling_layers}

Interspersed between convolutional layers are \textbf{Pooling (Subsampling) layers}. A pooling layer operates independently on each feature channel, replacing local spatial neighborhoods with an aggregated summary statistic~\cite{goodfellow2016deep}.

\begin{figure}[H]
\centering
\begin{tikzpicture}[scale=0.92]
  % 4x4 Input Matrix with 4 Colored Quadrants
  \begin{scope}[shift={(0,0)}]
    \node[font=\footnotesize\bfseries\color{politoNavy}] at (2.0, 4.4) {$4 \times 4$ Feature Map};

    % Top-Left (Green)
    \draw[fill=intuitiongreen!20, draw=intuitiongreen, thick] (0, 2) rectangle (2, 4);
    \node[font=\small\bfseries] at (0.5, 3.5) {12};
    \node[font=\small\bfseries\color{intuitiongreen}] at (1.5, 3.5) {\textbf{20}};
    \node[font=\small\bfseries] at (0.5, 2.5) {8};
    \node[font=\small\bfseries] at (1.5, 2.5) {12};

    % Top-Right (Orange/Red)
    \draw[fill=examred!20, draw=examred, thick] (2, 2) rectangle (4, 4);
    \node[font=\small\bfseries\color{examred}] at (2.5, 3.5) {\textbf{30}};
    \node[font=\small\bfseries] at (3.5, 3.5) {0};
    \node[font=\small\bfseries] at (2.5, 2.5) {2};
    \node[font=\small\bfseries] at (3.5, 2.5) {0};

    % Bottom-Left (Blue)
    \draw[fill=deepblue!20, draw=deepblue, thick] (0, 0) rectangle (2, 2);
    \node[font=\small\bfseries] at (0.5, 1.5) {34};
    \node[font=\small\bfseries] at (1.5, 1.5) {70};
    \node[font=\small\bfseries\color{deepblue}] at (0.5, 0.5) {\textbf{112}};
    \node[font=\small\bfseries] at (1.5, 0.5) {100};

    % Bottom-Right (SkyBlue / Teal)
    \draw[fill=deepdiveTeal!20, draw=deepdiveTeal, thick] (2, 0) rectangle (4, 2);
    \node[font=\small\bfseries\color{deepdiveTeal}] at (2.5, 1.5) {\textbf{37}};
    \node[font=\small\bfseries] at (3.5, 1.5) {4};
    \node[font=\small\bfseries] at (2.5, 0.5) {25};
    \node[font=\small\bfseries] at (3.5, 0.5) {12};

    % Grid lines inside
    \draw[gray!60] (0, 1) -- (4, 1);
    \draw[gray!60] (0, 3) -- (4, 3);
    \draw[gray!60] (1, 0) -- (1, 4);
    \draw[gray!60] (3, 0) -- (3, 4);
  \end{scope}

  % Max Pool Arrow
  \draw[-{Stealth}, ultra thick, politoNavy] (4.6, 2.0) -- (6.2, 2.0) node[midway, above, font=\scriptsize\bfseries, align=center] {Max Pool\\$2 \times 2, S=2$};

  % 2x2 Output Matrix
  \begin{scope}[shift={(6.8,1.0)}]
    \node[font=\footnotesize\bfseries\color{intuitiongreen}] at (1.0, 2.4) {$2 \times 2$ Downsampled Output};

    % Top-Left (Green max = 20)
    \draw[fill=intuitiongreen!35, draw=intuitiongreen, very thick] (0, 1) rectangle (1, 2);
    \node[font=\normalsize\bfseries\color{intuitiongreen!70!black}] at (0.5, 1.5) {20};

    % Top-Right (Red max = 30)
    \draw[fill=examred!35, draw=examred, very thick] (1, 1) rectangle (2, 2);
    \node[font=\normalsize\bfseries\color{examred!70!black}] at (1.5, 1.5) {30};

    % Bottom-Left (Blue max = 112)
    \draw[fill=deepblue!35, draw=deepblue, very thick] (0, 0) rectangle (1, 1);
    \node[font=\normalsize\bfseries\color{deepblue}] at (0.5, 0.5) {112};

    % Bottom-Right (Teal max = 37)
    \draw[fill=deepdiveTeal!35, draw=deepdiveTeal, very thick] (1, 0) rectangle (2, 1);
    \node[font=\normalsize\bfseries\color{deepdiveTeal}] at (1.5, 0.5) {37};
  \end{scope}
\end{tikzpicture}
\caption{Max pooling operation with a $2 \times 2$ filter and stride $S = 2$. The input feature map is partitioned into non-overlapping $2 \times 2$ colored quadrants. The operation routes the maximum value within each quadrant to the downsampled grid, halving height and width ($75\%$ spatial area reduction) with zero learnable parameters.}
\label{fig:max_pooling}
\end{figure}

\paragraph{Max Pooling vs. Average Pooling.}
\begin{itemize}
  \item \textbf{Max Pooling}: Outputs the maximum activation within the pooling window:
    \begin{equation}
    Y_{i, j, c} = \max_{m, n \in \{0, \dots, F-1\}} X_{i \cdot S + m, \, j \cdot S + n, \, c}.
    \label{eq:max_pooling_formula}
    \end{equation}
    Max pooling acts as a logical OR: it tests whether a specific visual feature (e.g., an edge, corner, or texture) is present anywhere within the receptive patch, regardless of its exact coordinate.
  \item \textbf{Average Pooling}: Computes the arithmetic mean of activations:
    \begin{equation}
    Y_{i, j, c} = \frac{1}{F^2} \sum_{m=0}^{F-1} \sum_{n=0}^{F-1} X_{i \cdot S + m, \, j \cdot S + n, \, c}.
    \label{eq:avg_pooling_formula}
    \end{equation}
\end{itemize}

\paragraph{Canonical Settings and Architectural Roles.}
The standard configuration in modern CNNs is a \textbf{$2 \times 2$ filter with stride $S = 2$} ($P = 0$). Because the stride equals the filter size, pooling windows do not overlap.
This operation yields four primary architectural benefits:
\begin{enumerate}
  \item \textbf{Zero Learnable Parameters}: Pooling is a fixed mathematical aggregation rule. It contains zero weights and zero biases ($N_{\text{params}} = 0$).
  \item \textbf{Spatial Downsampling and FLOP Reduction}: A $2 \times 2$ max pool with stride 2 cuts height and width in half ($H/2, W/2$). This reduces the spatial area by a factor of 4 ($75\%$ reduction in activations), slashing the memory footprint and FLOP burden of all subsequent layers.
  \item \textbf{Local Translation Invariance}: If the input image shifts by 1 pixel, the maximum value within a $2 \times 2$ region frequently remains identical. Pooling confers invariance to small spatial shifts and distortions.
  \item \textbf{Receptive Field Expansion}: By halving the spatial grid, a $3 \times 3$ filter applied in the layer \textit{following} a pooling operation covers twice the physical field of view on the original input image.
\end{enumerate}

\paragraph{Global Decision Synthesis: Why Place Fully Connected Heads After Convolution?}
\label{sec:why_fcn_after_cnn}

\begin{figure}[H]
\centering
\begin{tikzpicture}[scale=0.88, every node/.style={transform shape}]
  % Stage 1: Input Image
  \draw[fill=deepblue!15, draw=deepblue, thick] (0, 0) rectangle (1.2, 1.2);
  \node[font=\tiny, align=center] at (0.6, 0.6) {Image\\$32 \times 32$\\$\times 3$};

  % Arrow 1
  \draw[-{Stealth}, thick] (1.3, 0.6) -- (1.8, 0.6) node[midway, above, font=\tiny] {Conv};

  % Stage 2: Conv 1
  \draw[fill=deepblue!25, draw=deepblue, thick] (1.9, -0.1) rectangle (2.9, 1.3);
  \node[font=\tiny, align=center] at (2.4, 0.6) {Conv+ReLU\\$28 \times 28$\\$\times 6$};

  % Arrow 2
  \draw[-{Stealth}, thick] (3.0, 0.6) -- (3.5, 0.6) node[midway, above, font=\tiny] {Pool};

  % Stage 3: Pool 1
  \draw[fill=intuitiongreen!20, draw=intuitiongreen, thick] (3.6, 0.1) rectangle (4.4, 1.1);
  \node[font=\tiny, align=center] at (4.0, 0.6) {MaxPool\\$14 \times 14$\\$\times 6$};

  % Arrow 3
  \draw[-{Stealth}, thick] (4.5, 0.6) -- (5.0, 0.6) node[midway, above, font=\tiny] {Conv};

  % Stage 4: Conv 2
  \draw[fill=deepblue!35, draw=deepblue, thick] (5.1, 0.2) rectangle (5.9, 1.0);
  \node[font=\tiny, align=center] at (5.5, 0.6) {Conv+ReLU\\$10 \times 10$\\$\times 16$};

  % Arrow 4
  \draw[-{Stealth}, thick] (6.0, 0.6) -- (6.5, 0.6) node[midway, above, font=\tiny] {Pool};

  % Stage 5: Pool 2
  \draw[fill=intuitiongreen!30, draw=intuitiongreen, thick] (6.6, 0.3) rectangle (7.2, 0.9);
  \node[font=\tiny, align=center] at (6.9, 0.6) {MaxPool\\$5 \times 5$\\$\times 16$};

  % Arrow 5 (Flatten)
  \draw[-{Stealth}, thick] (7.3, 0.6) -- (7.8, 0.6) node[midway, above, font=\tiny] {Flatten};

  % Stage 6: Flatten Vector (400)
  \draw[fill=gray!20, draw=gray!70, thick] (7.9, -0.6) rectangle (8.2, 1.8);
  \node[font=\tiny, rotate=90] at (8.05, 0.6) {$400 \times 1$};

  % Arrow 6
  \draw[-{Stealth}, thick] (8.3, 0.6) -- (8.7, 0.6);

  % Stage 7: FC1 (120)
  \draw[fill=politoNavy!20, draw=politoNavy, thick] (8.8, -0.4) rectangle (9.1, 1.6);
  \node[font=\tiny, rotate=90] at (8.95, 0.6) {FC 120};

  % Arrow 7
  \draw[-{Stealth}, thick] (9.2, 0.6) -- (9.6, 0.6);

  % Stage 8: FC2 (84)
  \draw[fill=politoNavy!30, draw=politoNavy, thick] (9.7, -0.2) rectangle (10.0, 1.4);
  \node[font=\tiny, rotate=90] at (9.85, 0.6) {FC 84};

  % Arrow 8
  \draw[-{Stealth}, thick] (10.1, 0.6) -- (10.5, 0.6);

  % Stage 9: Output (10 classes)
  \draw[fill=examred!30, draw=examred, thick] (10.6, 0.1) rectangle (10.9, 1.1);
  \node[font=\tiny, rotate=90] at (10.75, 0.6) {Softmax 10};

  % Brackets for Division of Labor
  \draw[thick, deepblue, decorate, decoration={brace, amplitude=5pt}] (0, -0.9) -- (7.2, -0.9) node[midway, below=5pt, font=\scriptsize\bfseries] {Feature Extractor (Local \& Equivariant)};
  \draw[thick, examred, decorate, decoration={brace, amplitude=5pt}] (7.9, -0.9) -- (11.0, -0.9) node[midway, below=5pt, font=\scriptsize\bfseries] {Classifier Head (Global \& Invariant)};
\end{tikzpicture}
\caption{The end-to-end architectural pipeline of a Convolutional Neural Network (LeNet-5 archetype). The convolutional backbone extracts translation-equivariant local visual features across cascading spatial scales, culminating in dense fully connected layers that synthesize global spatial evidence to emit categorical class probabilities.}
\label{fig:end_to_end_cnn_pipeline}
\end{figure}

Having examined convolution, activations, and pooling, we address a foundational architectural puzzle emphasized by Prof.~Tatiana Tommasi: \textit{if Fully Connected Networks fail catastrophically on visual inputs (as proven in Section~\ref{sec:why_mlps_fail_images}), why does every classic deep convolutional architecture (LeNet-5, AlexNet, ZFNet, VGG) conclude by placing multiple Fully Connected Layers at the end of the convolutional stack?}

As illustrated in Figure~\ref{fig:end_to_end_cnn_pipeline}, this architectural design is neither an accident nor an inconsistency. Rather, it represents a mathematically principled \textbf{division of labor between geometric feature extraction and global semantic decision synthesis}.

\paragraph{1. The Complementary Division of Labor: Local Equivariance vs. Global Invariance.}
A complete visual recognition system must reconcile two fundamentally conflicting operational requirements:
\begin{enumerate}[leftmargin=*]
  \item \textbf{The Feature Extraction Requirement (Spatial Equivariance)}: In early and intermediate visual processing, the network must detect local features (edges, corners, textures, motifs, object parts) while \textbf{preserving their exact spatial coordinates}. The convolutional backbone achieves this through shared kernels sliding across the 2D grid, computing an activation tensor:
  \begin{equation}
  \mathbf{F} \in \R^{H' \times W' \times C'},
  \end{equation}
  where channel $c$ at coordinate $(i, j)$ answers the localized question: \textit{``Does visual feature $c$ appear in the receptive field centered at $(i, j)$?''} This representation is translation-equivariant: shifting the object shifts the activation map identically.
  
  \item \textbf{The Classification Requirement (Global Invariance and Relational Logic)}: In contrast, categorical object classification (e.g., determining whether an image depicts a ``Cat'' or a ``Truck'') is fundamentally a \textbf{global, translation-invariant, non-spatial decision problem}:
  \begin{equation}
  \hat{\mathbf{y}} = P(y \in \{1, \dots, K\} \mid \mathbf{X}) \in \Delta^{K-1}.
  \end{equation}
  The final output is a single categorical probability distribution over $K$ classes that must be invariant to where the object was found on the canvas. 
\end{enumerate}

\paragraph{2. The Relational Synthesis Problem: Part-to-Whole Composition.}
Why can the convolutional layers not simply make the final classification decision directly?
Because recognizing a holistic visual concept requires reasoning over the \textbf{global spatial co-occurrence and relative geometric arrangement of multiple independent parts}:
\begin{itemize}[leftmargin=*]
  \item Identifying a ``Face'' requires verifying that two eye features appear above a nose feature, which in turn appears above a mouth feature.
  \item Identifying a ``Car'' requires verifying that a front wheel and a rear wheel appear beneath a metallic chassis and windshield. If an ear appears floating next to an exhaust pipe, the semantic combination is invalid.
\end{itemize}
Because convolutional filters are spatially constrained (typically $3 \times 3$ or $5 \times 5$), a convolutional unit can only combine information within its local spatial receptive field. Even deep in the network, convolutional kernels compute spatially localized combinations with strictly shared weights; they cannot establish arbitrary, unconstrained all-to-all cross-spatial correlations between a feature in the upper-left quadrant and a feature in the lower-right quadrant without adding an impractical number of extra layers.

\paragraph{3. Why the Fully Connected Layer Solves Global Synthesis.}
The Fully Connected layers placed at the end of the convolutional backbone resolve this relational bottleneck through \textbf{dense all-to-all connectivity}. By unrolling the final 3D feature tensor into a 1D representation:
\begin{equation}
\mathbf{v} = \operatorname{vec}(\mathbf{F}) \in \R^{d_{\text{feat}}}, \quad \text{where } d_{\text{feat}} = H' \cdot W' \cdot C',
\end{equation}
the first fully connected layer computes:
\begin{equation}
\mathbf{a}_{\text{fc}}^{(1)} = \sigma\left( \mathbf{W}_{\text{fc}}^{(1)} \mathbf{v} + \mathbf{b}_{\text{fc}}^{(1)} \right).
\end{equation}
In scalar notation, each hidden unit $m$ in the dense head computes:
\begin{equation}
z_m = \sum_{c=1}^{C'} \sum_{i=1}^{H'} \sum_{j=1}^{W'} W_{m, (i, j, c)} F_{i, j, c} + b_m.
\label{eq:dense_all_to_all_synthesis}
\end{equation}
Equation~\eqref{eq:dense_all_to_all_synthesis} reveals the unique power of the FCN head: \textbf{every neuron receives inputs from every spatial location $(i, j)$ and every semantic channel $c$ simultaneously}. The dense layer can freely learn unconstrained global relational templates:
\begin{equation}
\big[ F_{1, 2, 42} \text{ (ear)} \gg 0 \;\land\; F_{3, 4, 88} \text{ (snout)} \gg 0 \big] \implies z_{\text{cat}} \gg 0.
\end{equation}
The fully connected layers act as the \textbf{global decision synthesizer}, mapping the spatially distributed evidence discovered by the convolutional backbone into final categorical class logits.

\paragraph{4. Why the Parameter Explosion Is Tamed at the End.}
A critical student question often arises: \textit{if dense layers cause catastrophic parameter explosion at the input, why does using an FCN at the end not suffer the same fate?}

The answer lies in the \textbf{exponential spatial compression} performed by the convolutional backbone:
\begin{itemize}[leftmargin=*]
  \item \textbf{At the Input (Raw Image)}: The raw input space is massive ($224 \times 224 \times 3 = 150,528$ dimensions). Flattening directly at the input requires over $616$ million weights for a single layer of $4,096$ units, operating on raw uninformative pixel intensities.
  \item \textbf{At the End (Latent Feature Tensor)}: Through cascading convolutions and pooling/striding operations, spatial dimensions are downsampled exponentially while channel capacity expands:
  \begin{equation}
  \begin{aligned}
  (224 \times 224 \times 3) &\xrightarrow{\text{stride/pool}} (112 \times 112 \times 64) \to (56 \times 56 \times 128) \\
  &\to (28 \times 28 \times 256) \to (14 \times 14 \times 512) \to (7 \times 7 \times 512).
  \end{aligned}
  \end{equation}
  The spatial area has been compressed by a factor of $(224 / 7)^2 = 32^2 = \mathbf{1,024\times}$!
  \item \textbf{Semantic Density}: When flattened at this late stage, the feature vector has length $d_{\text{feat}} = 7 \times 7 \times 512 = 25,088$. While still sizable, connecting to $4,096$ neurons requires $25,088 \times 4,096 \approx 102.7 \times 10^6$ parameters. Crucially, each of these $25,088$ inputs is not a fragile pixel intensity, but a rich, highly abstracted, non-linear semantic descriptor (e.g., an entire wheel, a dog snout, or a bird wing). 
\end{itemize}

\paragraph{5. The Decoupling Principle.}
Placing the FCN at the end decouples the computer vision problem into two specialized stages:
\begin{enumerate}[noitemsep]
  \item \textbf{Stage 1 (CNN Backbone)}: Solves the continuous geometry and physics of visual perception (spatial locality, translation equivariance, illumination invariance, edge filtering, and multi-scale hierarchy).
  \item \textbf{Stage 2 (FCN Head)}: Solves the discrete semantic categorization problem, acting as an unconstrained non-linear classifier that maps high-level feature evidence into class probabilities via Softmax:
  \begin{equation}
  \mathbf{z}_{\text{out}} = \mathbf{W}_{\text{out}} \mathbf{a}_{\text{fc}} + \mathbf{b}_{\text{out}} \in \R^K \implies \hat{y}_k = \frac{e^{z_k}}{\sum_{j=1}^K e^{z_j}}.
  \end{equation}
\end{enumerate}

\paragraph{6. The Modern Evolution: Dense FC Heads vs. Global Average Pooling (GAP).}
While dense classification heads provide immense expressive power for global spatial-relational reasoning, they introduce a significant drawback in deep networks: in architectures such as AlexNet and VGG-16, the fully connected layers account for over $85\%$--$90\%$ of the total parameters ($58.6\text{M}$ of $61\text{M}$ weights in AlexNet; $102.7\text{M}$ of $138\text{M}$ in VGG-16!). This concentration of parameters creates severe overfitting risks and requires aggressive regularization such as Dropout ($p=0.5$)~\cite{srivastava2014dropout}.

To eliminate this parameter bottleneck, modern architectures (pioneered by Network In Network~\cite{lin2013network} and established by GoogLeNet and ResNet~\cite{he2016deep}) replace the dense flattening head with \textbf{Global Average Pooling (GAP)}:
\begin{equation}
\bar{z}_c \triangleq \frac{1}{H' \times W'} \sum_{i=1}^{H'} \sum_{j=1}^{W'} F_{i, j, c}, \quad \forall c \in \{1, \dots, K\}.
\label{eq:gap_formula}
\end{equation}
By setting the final convolutional layer to output exactly $K$ feature channels (one channel per target category) and averaging each spatial map across its entire area, GAP produces a $K$-dimensional vector directly fed into the Softmax function.

\paragraph{Exam-Critical Comparison: Dense FCN Head vs. Global Average Pooling.}
In oral examinations, professors frequently ask students to compare the architectural trade-offs between Dense FC Heads and Global Average Pooling:
\begin{itemize}[leftmargin=*]
  \item \textbf{Spatial Position Sensitivity}: Dense FCN heads retain spatial position sensitivity (preserving the relative spatial coordinate within the $7 \times 7$ grid where a feature occurred). This property is indispensable in structured spatial tasks such as \textbf{object detection} (e.g., the bounding-box regression heads in Fast/Faster R-CNN RoI heads) and \textbf{keypoint estimation}. In contrast, GAP completely integrates out spatial coordinates, enforcing strict spatial invariance.
  \item \textbf{Parameter Efficiency and Overfitting}: Dense heads require tens of millions of parameters, making them prone to overfitting on small datasets. GAP introduces zero trainable parameters, acting as a structural regularizer that prevents overfitting and eliminates the need for Dropout.
  \item \textbf{Arbitrary Input Resolution}: Dense heads require a fixed input resolution because the weight matrix $\mathbf{W}_{\text{fc}}$ depends rigidly on the flattened dimension $H' W' C'$. GAP, by computing a spatial average, can seamlessly ingest input images of arbitrary resolutions at test time without architectural modifications.
\end{itemize}

\paragraph{Causal Bridge to Downstream Computer Vision Tasks.}
Having mastered the mechanics of convolutional feature extraction, effective receptive fields, and global decision synthesis, a pivotal architectural question emerges: \textit{how do these structural priors adapt when the downstream objective extends beyond simple whole-image categorization?} In real-world visual perception, an intelligent agent must not only categorize what is present, but delineate where objects reside via continuous bounding geometries or parse every individual coordinate with dense pixel-level categorical labels. In Section~\ref{sec:downstream_cv_tasks_imagenet}, we formalize the three canonical downstream computer vision tasks---Classification, Detection, and Segmentation---and trace the empirical revolution ignited by deep convolutional architectures across the historic ImageNet benchmark.

---

\section[Downstream computer vision tasks and ImageNet]{Downstream computer vision tasks and the ImageNet revolution}
\label{sec:downstream_cv_tasks_imagenet}

Convolutional neural networks achieved pre-eminence not merely by matching human perceptual capabilities, but by fundamentally transforming how machine learning formulates visual tasks. Rather than requiring distinct, handcrafted algorithmic pipelines for every problem, modern computer vision organizes perception into three foundational downstream tasks sharing common convolutional backbones but diverging in their spatial geometries, symmetry constraints, and output heads~\cite{prince2023understanding,bishop2024deep}.

\subsection{The taxonomy of visual perception: classification, detection, and segmentation}
\label{sec:cv_tasks_taxonomy}

Visual perception operates along a spectrum of spatial granularity, transitioning from coarse whole-image semantic abstraction to dense coordinate-level parsing.

\paragraph{1. Image Classification.}
In whole-image categorical classification, the input image $\mathbf{X} \in \R^{H \times W \times C}$ is mapped to a single probability distribution across $K$ mutually exclusive categories:
\begin{equation}
\hat{\mathbf{y}} = P(y \in \{1,\dots,K\} \mid \mathbf{X}) \in \Delta^{K-1},
\label{eq:image_classification_mapping}
\end{equation}
where $\Delta^{K-1} \triangleq \{\mathbf{p} \in \R^K : p_k \ge 0, \, \sum_{k=1}^K p_k = 1\}$ denotes the probability simplex. Given a high-level latent feature representation $\mathbf{h} \in \R^d$ produced by the convolutional backbone and global pooling, the classification head computes logits $\mathbf{z} \in \R^K$ and Softmax probabilities:
\begin{equation}
\mathbf{z} = \mathbf{W}_{\text{head}}\mathbf{h} + \mathbf{b}_{\text{head}}, \quad \hat{y}_k = \frac{\exp(z_k)}{\sum_{j=1}^K \exp(z_j)}, \quad \forall k \in \{1,\dots,K\}.
\label{eq:classification_softmax_logits}
\end{equation}
Training minimizes the Categorical Cross-Entropy (Negative Log-Likelihood) loss:
\begin{equation}
\mathcal{L}_{\text{CE}}(\mathbf{y}, \hat{\mathbf{y}}) \triangleq -\sum_{k=1}^K y_k \ln \hat{y}_k = -\ln \hat{y}_{c^*},
\label{eq:classification_cross_entropy}
\end{equation}
where $c^* \in \{1,\dots,K\}$ denotes the ground-truth categorical label. Categorical classification enforces strict \textbf{translation invariance}: shifting the image coordinates must leave the predicted class probability vector completely unaltered:
\begin{equation}
f\left( \mathcal{T}_{\boldsymbol{\delta}}\mathbf{X} \right) = f(\mathbf{X}), \quad \forall \boldsymbol{\delta} \in \mathbb{Z}^2.
\end{equation}
While intermediate convolutional layers preserve spatial geometry through translation equivariance, the terminal global pooling stage or dense flattening head intentionally destroys spatial coordinates to enforce invariant decision boundaries.

\paragraph{2. Comparative Task Taxonomy.}
Table~\ref{tab:cv_tasks_comparison} establishes the formal mathematical comparison among the three canonical vision tasks.

\begin{table}[H]
\centering
\footnotesize
\setlength{\tabcolsep}{5pt}
\caption{Taxonomy of canonical computer vision tasks comparing target objectives, output mathematical domains, symmetry requirements, head architectures, and primary loss formulations.}
\label{tab:cv_tasks_comparison}
\vspace{4pt}
\begin{tabularx}{\textwidth}{>{\raggedright\arraybackslash}p{2.3cm} >{\raggedright\arraybackslash}p{2.6cm} >{\raggedright\arraybackslash}p{2.2cm} >{\raggedright\arraybackslash}p{2.4cm} >{\raggedright\arraybackslash}p{2.5cm} >{\raggedright\arraybackslash}X}
\toprule
\textbf{Task} & \textbf{Objective} & \textbf{Output Domain} & \textbf{Symmetry} & \textbf{Head Architecture} & \textbf{Primary Loss} \\
\midrule
\textbf{Image Classification} & Whole-image categorical recognition & Simplex $\Delta^{K-1}$ ($1 \times K$) & Translation Invariance ($f(\mathcal{T}\mathbf{X}) = f(\mathbf{X})$) & Dense FC or Global Average Pooling (GAP) & Categorical Cross-Entropy ($\mathcal{L}_{\text{CE}}$) \\
\addlinespace
\textbf{Object Detection} & Simultaneous classification \& spatial localization of $M$ objects & $\{1,\dots,K\} \times \R^4$ per instance & Equivariant backbone, coordinate regression & RoI Pooling / RoI Align with sibling FC branches & Multi-task: Cross-Entropy $+$ Smooth $L_1$ \\
\addlinespace
\textbf{Semantic Segmentation} & Dense pixel-level categorical classification & Grid $\{1,\dots,K\}^{H \times W}$ ($H \times W \times K$) & Translation Equivariance ($Y(\mathcal{T}\mathbf{X}) = \mathcal{T}[Y(\mathbf{X})]$) & Fully convolutional decoder (U-Net) with skip connections & Mean pixel-wise Cross-Entropy ($\mathcal{L}_{\text{seg}}$) \\
\bottomrule
\end{tabularx}
\end{table}

\subsection{Object detection: continuous bounding boxes, IoU, and non-maximum suppression}
\label{sec:object_detection_mechanics}

Object detection extends classification by requiring an algorithm to simultaneously answer \textit{``What is present?''} and \textit{``Where is it located?''} for an arbitrary, unknown number of instances $M \ge 0$ scattered across the image~\cite{prince2023understanding}.

\paragraph{1. Continuous 4D Bounding Box Parameterization.}
Each visual object instance is tightly enclosed by a 2D rectangular bounding box. In continuous spatial coordinates, a bounding box is parameterized by a 4D vector $\mathbf{b} \in \R^4$:
\begin{equation}
\mathbf{b} = [x_c, \, y_c, \, w, \, h]^\top,
\label{eq:bounding_box_center_param}
\end{equation}
where $(x_c, y_c) \in [0, W] \times [0, H]$ denotes the centroid of the bounding box, and $w > 0, h > 0$ denote its spatial width and height. Alternatively, in corner coordinate notation:
\begin{equation}
\mathbf{b}_{\text{corner}} = [x_{\min}, \, y_{\min}, \, x_{\max}, \, y_{\max}]^\top = \left[x_c - \frac{w}{2}, \, y_c - \frac{h}{2}, \, x_c + \frac{w}{2}, \, y_c + \frac{h}{2}\right]^\top.
\label{eq:bounding_box_corner_param}
\end{equation}

\paragraph{2. The Multi-Task Detection Objective.}
Unlike classification networks which optimize a single loss, modern object detectors optimize a composite \textbf{multi-task objective} balancing discrete categorical recognition with continuous geometric regression~\cite{girshick2015fast}:
\begin{equation}
\mathcal{L}(p, u, \mathbf{t}^u, \mathbf{v}) = \mathcal{L}_{\text{cls}}(p, u) + \lambda \, [u \ge 1] \, \mathcal{L}_{\text{loc}}(\mathbf{t}^u, \mathbf{v}).
\label{eq:detection_multitask_loss}
\end{equation}
In Equation~\eqref{eq:detection_multitask_loss}:
\begin{itemize}[leftmargin=*]
  \item $\mathcal{L}_{\text{cls}}(p, u) = -\ln p_u$ is the log loss for true category $u \in \{0, 1, \dots, K\}$, where class $u=0$ designates background.
  \item The Iverson bracket $[u \ge 1]$ evaluates to $1$ when the region corresponds to a positive foreground object and $0$ for background proposals, ensuring localization gradients are never computed on empty canvas.
  \item $\mathbf{t}^u = (t_x^u, t_y^u, t_w^u, t_h^u)$ is the predicted scale-invariant parameterized coordinate offset vector, regressed toward normalized ground-truth target $\mathbf{v} = (v_x, v_y, v_w, v_h)$.
  \item $\mathcal{L}_{\text{loc}}$ is evaluated using the robust \textbf{Smooth $L_1$ loss}:
    \begin{equation}
    \mathcal{L}_{\text{loc}}(\mathbf{t}^u, \mathbf{v}) = \sum_{i \in \{x, y, w, h\}} \operatorname{smooth}_{L_1}(t_i^u - v_i), \quad \operatorname{smooth}_{L_1}(z) = \begin{cases} 0.5 z^2 & \text{if } |z| < 1, \\ |z| - 0.5 & \text{otherwise.} \end{cases}
    \label{eq:smooth_l1_loss}
    \end{equation}
    The Smooth $L_1$ loss behaves quadratically ($L_2$) near zero error, ensuring smooth gradient descent convergence, but transitions to linear ($L_1$) behavior for large regression errors, preventing exploding gradients caused by wild outlier proposals.
\end{itemize}

\paragraph{3. Quantitative Localization Metric: Intersection over Union (IoU).}
To quantify how accurately a predicted bounding box $B_{\text{pred}}$ matches ground-truth box $B_{\text{gt}}$, computer vision uses the \textbf{Jaccard Index}, universally termed \textbf{Intersection over Union (IoU)}:
\begin{equation}
\text{IoU}(B_{\text{pred}}, B_{\text{gt}}) \triangleq \frac{\operatorname{Area}(B_{\text{pred}} \cap B_{\text{gt}})}{\operatorname{Area}(B_{\text{pred}} \cup B_{\text{gt}})} = \frac{\operatorname{Area}(B_{\text{pred}} \cap B_{\text{gt}})}{\operatorname{Area}(B_{\text{pred}}) + \operatorname{Area}(B_{\text{gt}}) - \operatorname{Area}(B_{\text{pred}} \cap B_{\text{gt}})}.
\label{eq:iou_definition}
\end{equation}
Analytically, the intersection area of two axis-aligned bounding boxes is evaluated by coordinate clipping:
\begin{align}
\operatorname{Area}(B_{\text{pred}} \cap B_{\text{gt}}) &= \max\left(0, \min(x_{\max}^p, x_{\max}^g) - \max(x_{\min}^p, x_{\min}^g)\right) \nonumber \\
&\quad \times \max\left(0, \min(y_{\max}^p, y_{\max}^g) - \max(y_{\min}^p, y_{\min}^g)\right).
\label{eq:iou_overlap_area}
\end{align}

\begin{figure}[H]
\centering
\resizebox{0.96\textwidth}{!}{%
\begin{tikzpicture}
  % -------------------------------------------------------------
  % CANVAS FRAME (Simulated Image)
  % -------------------------------------------------------------
  \draw[draw=gray!50, fill=gray!5, thick, rounded corners=3pt] (0, 0) rectangle (9.8, 6.4);
  \node[font=\scriptsize\bfseries\color{gray!70}, anchor=north west] at (0.3, 6.15) {Image Canvas $\mathbf{X} \in \mathbb{R}^{H \times W \times 3}$};

  % -------------------------------------------------------------
  % 1. BACKGROUND FILLS (Drawn first so fills never obscure text or outlines)
  % -------------------------------------------------------------
  % Ground Truth background: [0.6, 4.8] x [0.8, 3.8]
  \draw[draw=none, fill=deepblue!10] (0.6, 0.8) rectangle (4.8, 3.8);
  % Predicted Box background: [3.4, 7.6] x [1.8, 4.8]
  \draw[draw=none, fill=examred!10] (3.4, 1.8) rectangle (7.6, 4.8);
  % Overlap (Intersection) background: [3.4, 4.8] x [1.8, 3.8]
  \draw[draw=none, fill=deepdiveTeal!30] (3.4, 1.8) rectangle (4.8, 3.8);

  % -------------------------------------------------------------
  % 2. OUTLINES
  % -------------------------------------------------------------
  \draw[draw=politoNavy, very thick] (0.6, 0.8) rectangle (4.8, 3.8);
  \draw[draw=examred, very thick, dashed] (3.4, 1.8) rectangle (7.6, 4.8);
  \draw[draw=deepdiveTeal, thick] (3.4, 1.8) rectangle (4.8, 3.8);

  % -------------------------------------------------------------
  % 3. LABELS AND CENTROIDS (Completely disjoint spatial regions)
  % -------------------------------------------------------------
  % Ground Truth Label (Stacked two lines in non-overlapping left region: x in [0.8, 2.4], y in [3.1, 3.65])
  \node[font=\footnotesize\bfseries\color{politoNavy}, align=left, anchor=north west] at (0.8, 3.65) {
    Ground Truth\\[1pt]$B_{\text{gt}}$
  };
  
  % Ground Truth Centroid (Exact geometric center: xc=2.7, yc=2.3)
  \node[circle, fill=politoNavy, inner sep=1.8pt] (gt_c) at (2.7, 2.3) {};
  \node[font=\scriptsize\bfseries\color{politoNavy}, anchor=north east] at (2.55, 2.2) {$(x_c^{\text{gt}}, y_c^{\text{gt}})$};

  % Predicted Box Label (Stacked two lines in non-overlapping top-right region: x in [5.8, 7.4], y in [4.1, 4.65])
  \node[font=\footnotesize\bfseries\color{examred}, align=right, anchor=north east] at (7.4, 4.65) {
    Predicted\\[1pt]$B_{\text{pred}}$
  };
  
  % Predicted Centroid (Exact geometric center: xc=5.5, yc=3.3)
  % Label placed below-right so it has ample separation from top-right title
  \node[circle, fill=examred, inner sep=1.8pt] (pred_c) at (5.5, 3.3) {};
  \node[font=\scriptsize\bfseries\color{examred}, anchor=north west] at (5.65, 3.15) {$(x_c^{\text{pred}}, y_c^{\text{pred}})$};

  % Overlap Area (Intersection) Label (Center of intersection: x=4.1, y=2.8)
  \node[font=\scriptsize\bfseries\color{deepdiveTeal!90!black}, align=center] at (4.1, 2.8) {Intersection\\[2pt]$B_{\text{pred}} \cap B_{\text{gt}}$};

  % -------------------------------------------------------------
  % 4. DIMENSION BRACES FOR B_pred
  % -------------------------------------------------------------
  % Width brace at y=5.05 (centered at x=5.5; canvas text is at x=0.3, y=6.15 -> completely disjoint)
  \draw[thick, examred, decorate, decoration={brace, amplitude=4pt}] (3.4, 5.05) -- (7.6, 5.05)
    node[midway, above=3pt, font=\scriptsize\bfseries\color{examred}] {Width $w_{\text{pred}}$};
    
  % Height brace at x=7.85 (midway at y=3.3; label ends at x=9.5, canvas border is at x=9.8)
  \draw[thick, examred, decorate, decoration={brace, amplitude=4pt}] (7.85, 4.8) -- (7.85, 1.8)
    node[midway, right=4pt, font=\scriptsize\bfseries\color{examred}] {Height $h_{\text{pred}}$};

  % -------------------------------------------------------------
  % 5. MATHEMATICAL IoU PANEL (Generously shifted to x=10.8 to guarantee zero collision)
  % -------------------------------------------------------------
  \begin{scope}[shift={(10.8, 0.4)}]
    % Formula Header (IoU =)
    \node[font=\normalsize\bfseries, anchor=east] at (0.6, 3.0) {$\text{IoU} = $};

    % Division bar
    \draw[very thick, black] (0.8, 3.0) -- (5.4, 3.0);

    % Numerator: Intersection Subpanel
    \draw[draw=deepdiveTeal, fill=deepdiveTeal!20, thick, rounded corners=2pt] (1.0, 3.4) rectangle (2.8, 4.8);
    \node[font=\tiny\bfseries\color{deepdiveTeal!90!black}, align=center] at (1.9, 4.1) {$\operatorname{Area}(B_{\text{pred}} \cap B_{\text{gt}})$\\(Intersection)};
    \node[anchor=west, font=\scriptsize] at (3.0, 4.1) {$[w_{\cap} \times h_{\cap}]$ overlap area};

    % Denominator: Union Subpanel
    \draw[draw=politoNavy, fill=politoNavy!15, thick, rounded corners=2pt] (1.0, 1.2) rectangle (2.8, 2.6);
    \node[font=\tiny\bfseries\color{politoNavy}, align=center] at (1.9, 1.9) {$\operatorname{Area}(B_{\text{pred}} \cup B_{\text{gt}})$\\(Total Union)};
    \node[anchor=west, font=\scriptsize] at (3.0, 1.9) {$\operatorname{Area}_{\text{pred}} + \operatorname{Area}_{\text{gt}} - \operatorname{Area}_{\cap}$};

    % Evaluation Criterion Box
    \node[draw=examred, fill=examback, thick, rounded corners=3pt, inner sep=6pt, anchor=north west] at (0.0, 0.5) {
      \begin{minipage}{5.6cm}
        \scriptsize\bfseries\textcolor{examred}{Standard Evaluation Gate:}\\[3pt]
        \normalfont\scriptsize
        $\bullet$ $\text{IoU} \ge 0.5 \implies$ \textbf{True Positive (TP)} (Correct detection).\\
        $\bullet$ $\text{IoU} < 0.5 \implies$ \textbf{False Positive (FP)} (Localization error).\\
        $\bullet$ Strict COCO AP uses average over thresholds $\text{IoU} \in [0.50 : 0.05 : 0.95]$.
      \end{minipage}
    };
  \end{scope}
\end{tikzpicture}%
}
\caption{Continuous 4D bounding box parameterization and Intersection over Union (IoU) localization evaluation metric. The ratio between the area of overlap (intersection) and the total combined area (union) quantifies spatial agreement; detections with $\text{IoU} \ge 0.5$ are deemed true positives in the PASCAL VOC benchmark.}
\label{fig:iou_bounding_box_schematic}
\end{figure}

As shown in Figure~\ref{fig:iou_bounding_box_schematic}, in benchmark evaluations (e.g., PASCAL VOC), a candidate detection is scored as a \textbf{True Positive (TP)} if and only if $\text{IoU}(B_{\text{pred}}, B_{\text{gt}}) \ge 0.5$ and the categorical label matches ground truth. If $\text{IoU} < 0.5$, the prediction is penalized as a \textbf{False Positive (FP)}. In the modern MS-COCO benchmark, performance is evaluated across a spectrum of ten thresholds $\text{IoU} \in [0.50 : 0.05 : 0.95]$, rewarding sub-pixel localization accuracy.

\paragraph{4. Non-Maximum Suppression (NMS) Algorithm.}
Because dense detection architectures evaluate thousands of potential spatial anchors or proposals across adjacent coordinates, the network routinely fires dozens of highly confident bounding boxes around the exact same target. To eliminate this spatial redundancy, modern detectors execute \textbf{Non-Maximum Suppression (NMS)}, a greedy deduplication algorithm:
\begin{enumerate}[leftmargin=*]
  \item \textbf{Input}: A set of candidate bounding boxes $\mathcal{B} = \{B_1, \dots, B_N\}$, corresponding category confidence scores $\mathcal{S} = \{s_1, \dots, s_N\}$, score threshold $s_{\min}$, and IoU suppression threshold $N_{\text{th}} \in [0.3, 0.5]$.
  \item \textbf{Pruning}: Discard all candidate boxes whose classification score satisfies $s_i < s_{\min}$. Initialize the set of confirmed final detections as empty: $\mathcal{D} \leftarrow \emptyset$.
  \item \textbf{Iterative Suppression}: While $\mathcal{B} \neq \emptyset$:
    \begin{enumerate}[noitemsep]
      \item Select the bounding box with the highest confidence score:
        \begin{equation}
        M \leftarrow \argmax_{B_i \in \mathcal{B}} s_i.
        \end{equation}
      \item Move $M$ into the set of confirmed detections: $\mathcal{D} \leftarrow \mathcal{D} \cup \{M\}$, and remove $M$ from candidates: $\mathcal{B} \leftarrow \mathcal{B} \setminus \{M\}$.
      \item For every remaining candidate box $B_j \in \mathcal{B}$, compute spatial overlap with $M$. If:
        \begin{equation}
        \text{IoU}(M, B_j) \ge N_{\text{th}},
        \end{equation}
        discard $B_j$ from $\mathcal{B}$ as a duplicate detection.
    \end{enumerate}
  \item \textbf{Output}: Return the pruned, deduplicated detection set $\mathcal{D}$.
\end{enumerate}

\paragraph{5. Computational Evolution: From Sliding Windows to Shared Convolutions.}
The computational challenge of object detection motivated one of deep learning's greatest architectural evolutions:
\begin{itemize}[leftmargin=*]
  \item \textbf{Naive Sliding Windows}: Evaluated every possible spatial coordinate $(x, y)$, scale, and aspect ratio independently through an entire convolutional network. Evaluating millions of crops independently incurred an astronomical computational cost of $O(H^2 W^2)$ forward passes per image.
  \item \textbf{R-CNN (Girshick et al., 2014)}~\cite{girshick2014rich}: Reduced the search space by utilizing an external Selective Search heuristic to propose $\sim 2,000$ candidate regions. However, R-CNN cropped and warped each candidate patch independently through AlexNet, requiring $\approx 47$ seconds per image on GPU and preventing real-time deployment.
  \item \textbf{Fast R-CNN (Girshick, 2015)}~\cite{girshick2015fast}: Introduced \textbf{Shared Convolutional Feature Maps}. Instead of passing 2,000 image patches through the CNN, Fast R-CNN feeds the entire input image through the backbone \textit{once} to extract a high-resolution shared feature map $\mathbf{F}$.
  \item \textbf{Region of Interest (RoI) Pooling}: Given arbitrary proposed rectangle coordinates $(r, c, h_r, w_r)$ projected onto the shared feature map, RoI Pooling divides the sub-region into a fixed $7 \times 7$ grid of sub-windows and performs max pooling within each grid cell:
    \begin{equation}
    (h_r \times w_r \times C) \xrightarrow{\text{RoI Pooling}} (7 \times 7 \times C).
    \label{eq:roi_pooling_mapping}
    \end{equation}
    Because every proposal extracts its fixed-size feature vector directly from the shared feature map, Fast R-CNN replaced 2,000 forward passes with 1 forward pass, achieving a \textbf{$146\times$ acceleration} during inference.
  \item \textbf{Faster R-CNN (Ren et al., 2015)}~\cite{ren2015faster}: Replaced the CPU-bound Selective Search heuristic entirely with an internal, fully convolutional \textbf{Region Proposal Network (RPN)} that slides $3 \times 3$ convolutional filters over the shared feature map, predicting multi-scale anchor offsets and objectness scores in unified end-to-end GPU memory.
\end{itemize}

\subsection{Semantic image segmentation: dense pixel categorization and encoder-decoder topologies}
\label{sec:semantic_segmentation_mechanics}

While object detection encloses targets within coarse bounding boxes, \textbf{Semantic Image Segmentation} requires dense, pixel-level visual understanding: assigning a categorical class label to every individual coordinate $(i, j)$ of the image grid~\cite{prince2023understanding,bishop2024deep}.

\paragraph{1. Dense Pixel-Level Mathematical Mapping.}
In semantic segmentation, the network maps an image $\mathbf{X} \in \R^{H \times W \times C}$ to a dense categorical label canvas:
\begin{equation}
\hat{y}_{i, j} \in \{1, \dots, K\}, \quad \forall (i, j) \in \{1, \dots, H\} \times \{1, \dots, W\}.
\label{eq:semantic_segmentation_mapping}
\end{equation}
The architecture outputs a 3D tensor of unnormalized pixel logits $\mathbf{Z} \in \R^{H \times W \times K}$. At every spatial coordinate $(i, j)$, a spatial Softmax produces a categorical distribution over classes:
\begin{equation}
\hat{P}(y_{i,j} = k \mid \mathbf{X}) = \frac{\exp(Z_{i, j, k})}{\sum_{c=1}^K \exp(Z_{i, j, c})}, \quad \forall k \in \{1, \dots, K\}.
\label{eq:pixelwise_softmax}
\end{equation}
The network is optimized end-to-end by minimizing the mean pixel-wise Categorical Cross-Entropy loss:
\begin{equation}
\mathcal{L}_{\text{seg}}(\mathbf{Y}, \hat{\mathbf{Y}}) = -\frac{1}{H \cdot W} \sum_{i=1}^H \sum_{j=1}^W \sum_{k=1}^K Y_{i, j, k} \ln \hat{P}(y_{i,j} = k \mid \mathbf{X}).
\label{eq:pixelwise_cross_entropy}
\end{equation}
Unlike whole-image classification (which demands translation invariance), semantic segmentation strictly demands \textbf{translation equivariance}: shifting an input object by $\boldsymbol{\delta}$ must shift the corresponding segmentation mask by the identical displacement vector $\boldsymbol{\delta}$.

\paragraph{2. The Architectural Dilemma: Semantic Context vs. Spatial Resolution.}
Designing deep segmentation networks introduces a fundamental representational trade-off:
\begin{itemize}[leftmargin=*]
  \item \textbf{Context Aggregation (The Need to Downsample)}: Individual raw pixels contain no semantic meaning (a single brown pixel could belong to a wooden chair, a dog, or a tree trunk). Correctly recognizing what an object is requires large effective receptive fields spanning wide contextual areas. As derived in Section~\ref{sec:conv_layer_anatomy}, pooling and strided convolutions downsample spatial dimensions ($H \to H/32$), allowing deep neurons to encode global scene semantics (``\textit{What}'').
  \item \textbf{Boundary Obliteration (The Cost of Downsampling)}: However, downsampling the spatial grid by a factor of 32 obliterates high-frequency spatial boundaries and geometric edges (``\textit{Where}''). Directly upsampling a coarse $7 \times 7$ feature map produces blurry, blob-like segmentation masks devoid of fine contours.
\end{itemize}

\paragraph{3. The Encoder-Decoder Architectural Topology.}
To reconcile semantic context with sharp spatial localization, modern deep segmentation architectures adopt the \textbf{Encoder-Decoder (Hourglass)} paradigm:
\begin{enumerate}[leftmargin=*]
  \item \textbf{Contractive Path (Encoder)}: A standard convolutional backbone (e.g., VGG or ResNet) that systematically alternates convolutions with strided pooling. Spatial height and width are progressively halved while channel capacity doubles, culminating in a low-resolution semantic bottleneck tensor:
    \begin{equation}
    \mathbf{Z}_{\text{bottleneck}} \in \R^{\frac{H}{32} \times \frac{W}{32} \times C_{\text{deep}}}.
    \end{equation}
  \item \textbf{Expansive Path (Decoder)}: A sequence of spatial upsampling layers and convolutions that progressively doubles spatial resolution while reducing channel depth ($\frac{H}{32} \to \frac{H}{16} \to \dots \to H$), reconstructing the full-resolution prediction canvas of dimensions $H \times W \times K$.
\end{enumerate}

\paragraph{4. Mechanisms of Spatial Upsampling.}
Restoring spatial resolution in the decoder requires principled upsampling operators:
\begin{enumerate}[leftmargin=*]
  \item \textbf{Bilinear and Nearest-Neighbor Interpolation}: Fixed, parameter-free geometric interpolation. While computationally efficient, these operators contain zero learnable weights and cannot adaptively adjust to complex boundary statistics.
  \item \textbf{Bed of Nails and Max Unpooling (Noh et al., 2015)}~\cite{noh2015learning,zeiler2014visualizing}: During the contractive path, each max pooling operation records the exact switch coordinates (argmax pixel locations within each $2 \times 2$ window). During decoding, max unpooling places each activation back into its exact original switch coordinate, filling all other positions with zeros. This preserves high-frequency edge locations at zero parameter cost.
  \item \textbf{Transposed Convolution (Fractionally-Strided Convolution)}:
    In standard 1D forward convolution with stride $S$, the linear transformation is expressed as matrix multiplication $\mathbf{y} = \mathbf{C} \mathbf{x}$, where $\mathbf{C} \in \R^{M \times N}$ ($M < N$) is a banded Toeplitz matrix downsampling the input.
    The \textbf{transposed convolution} applies the transpose operator:
    \begin{equation}
    \hat{\mathbf{x}} = \mathbf{C}^\top \mathbf{y}, \quad \text{where } \mathbf{C}^\top \in \R^{N \times M}.
    \end{equation}
    Spatially, a transposed convolution with stride $S > 1$ inserts $S - 1$ zeros between adjacent rows and columns of the input feature map, subsequently applying a standard convolution with learnable kernel weights $\mathbf{K} \in \R^{C_{\text{out}} \times C_{\text{in}} \times k \times k}$. The output spatial dimension is governed by:
    \begin{equation}
    H_{\text{out}} = (H_{\text{in}} - 1) \cdot S - 2P + k + P_{\text{out}},
    \label{eq:transposed_conv_output_dim}
    \end{equation}
    where $P$ is input padding and $P_{\text{out}}$ is output padding. Because kernel weights $\mathbf{K}$ are optimized via backpropagation, transposed convolutions learn data-driven, task-adaptive upsampling filters.
\end{enumerate}

\paragraph{5. Skip Connections and the U-Net Archetype.}
Even with learnable transposed convolutions, deep encoders lose irrecoverable fine-grained spatial information through repeated pooling. In their landmark work, Long, Shelhamer, and Darrell (2015) introduced \textbf{Fully Convolutional Networks (FCN)}~\cite{long2015fully}, which added shallow feature maps to deep upsampled maps via element-wise addition (FCN-32s, FCN-16s, FCN-8s).

This concept was perfected by Olaf Ronneberger et al. (2015) in the \textbf{U-Net} architecture~\cite{ronneberger2015u}, illustrated in Figure~\ref{fig:unet_architecture_schematic}. U-Net introduced \textbf{long skip connections with channel concatenation}:
\begin{equation}
\mathbf{X}_{\text{dec}}^{(l)} = \operatorname{ConvBlock}\left( \left[ \operatorname{Upsample}\left(\mathbf{X}_{\text{dec}}^{(l-1)}\right), \; \mathbf{X}_{\text{enc}}^{(l)} \right] \right).
\label{eq:unet_skip_concatenation}
\end{equation}
By concatenating pristine, uncompressed feature maps $\mathbf{X}_{\text{enc}}^{(l)}$ from early encoder layers directly into corresponding decoder stages, U-Net supplies exact boundary localization coordinates directly to the decoder, overcoming the spatial bottleneck while retaining high-level semantic context.

\begin{figure}[H]
\centering
\resizebox{0.95\textwidth}{!}{%
\begin{tikzpicture}
  % Styles
  \tikzset{
    encblock/.style={rectangle, draw=politoNavy, fill=politoNavy!20, thick, minimum width=1.3cm, align=center, font=\scriptsize},
    decblock/.style={rectangle, draw=deepblue, fill=skyBlue!30, thick, minimum width=1.3cm, align=center, font=\scriptsize},
    botblock/.style={rectangle, draw=deepdiveTeal, fill=deepdiveTeal!25, very thick, text width=4.3cm, align=center, font=\scriptsize},
    poolarrow/.style={-{Stealth[scale=0.85]}, thick, examred},
    upsparrow/.style={-{Stealth[scale=0.85]}, thick, deepdiveTeal},
    skiparrow/.style={-{Stealth[scale=0.85]}, very thick, dashed, intuitiongreen}
  }

  % CONTRACTIVE PATH (LEFT)
  % Level 1 Encoder
  \node[encblock, minimum height=3.6cm] (e1) at (0, 3.6) {Conv $3\times 3$\\[2pt]$H \times W$\\[2pt]$C=64$};
  \node[font=\tiny\bfseries\color{politoNavy}, above=4pt of e1] {Encoder Level 1};

  % Pool 1
  \node[circle, fill=examred!20, draw=examred, inner sep=2pt] (p1) at (0, 1.4) {\tiny $\downarrow 2$};
  \draw[poolarrow] (e1.south) -- (p1.north);

  % Level 2 Encoder
  \node[encblock, minimum height=2.6cm] (e2) at (1.9, 1.4) {Conv $3\times 3$\\[2pt]$(H/2) \times (W/2)$\\[2pt]$C=128$};
  \draw[poolarrow] (p1.east) -- ++(0.35, 0) |- (e2.west);

  % Pool 2
  \node[circle, fill=examred!20, draw=examred, inner sep=2pt] (p2) at (1.9, -0.4) {\tiny $\downarrow 2$};
  \draw[poolarrow] (e2.south) -- (p2.north);

  % Level 3 Encoder
  \node[encblock, minimum height=1.8cm] (e3) at (3.8, -0.4) {Conv $3\times 3$\\[2pt]$(H/4) \times (W/4)$\\[2pt]$C=256$};
  \draw[poolarrow] (p2.east) -- ++(0.35, 0) |- (e3.west);

  % Pool 3 (shifted cleanly below e3, above bottleneck connection)
  \node[circle, fill=examred!20, draw=examred, inner sep=2pt] (p3) at (3.8, -2.1) {\tiny $\downarrow 2$};
  \draw[poolarrow] (e3.south) -- (p3.north);

  % BOTTLENECK (BOTTOM CENTER)
  \node[botblock, minimum height=1.3cm] (bot) at (7.0, -3.2) {Bottleneck Conv Block\\[2pt]$(H/8) \times (W/8) \times 512$\\[2pt]\textbf{Deepest Context (``What'')}};
  \draw[poolarrow] (p3.south) |- (bot.west);

  % EXPANSIVE PATH (RIGHT)
  % Up 3 (shifted cleanly below d3, above bottleneck connection)
  \node[circle, fill=deepdiveTeal!20, draw=deepdiveTeal, inner sep=2pt] (u3) at (10.2, -2.1) {\tiny $\uparrow 2$};
  \draw[upsparrow] (bot.east) -| (u3.south);

  % Level 3 Decoder
  \node[decblock, minimum height=1.8cm] (d3) at (10.2, -0.4) {UpConv + Cat\\[2pt]$(H/4) \times (W/4)$\\[2pt]$C=256$};
  \draw[upsparrow] (u3.north) -- (d3.south);

  % Up 2
  \node[circle, fill=deepdiveTeal!20, draw=deepdiveTeal, inner sep=2pt] (u2) at (12.1, -0.4) {\tiny $\uparrow 2$};
  \draw[upsparrow] (d3.east) -- (u2.west);

  % Level 2 Decoder
  \node[decblock, minimum height=2.6cm] (d2) at (12.1, 1.4) {UpConv + Cat\\[2pt]$(H/2) \times (W/2)$\\[2pt]$C=128$};
  \draw[upsparrow] (u2.north) -- (d2.south);

  % Up 1
  \node[circle, fill=deepdiveTeal!20, draw=deepdiveTeal, inner sep=2pt] (u1) at (14.0, 1.4) {\tiny $\uparrow 2$};
  \draw[upsparrow] (d2.east) -- (u1.west);

  % Level 1 Decoder
  \node[decblock, minimum height=3.6cm] (d1) at (14.0, 3.6) {UpConv + Cat\\[2pt]$H \times W$\\[2pt]$C=64$};
  \draw[upsparrow] (u1.north) -- (d1.south);
  \node[font=\tiny\bfseries\color{deepblue}, above=4pt of d1] {Decoder Level 1};

  % Output Head (1x1 conv)
  \node[rectangle, draw=examred, fill=examred!20, thick, minimum height=3.6cm, minimum width=0.9cm, align=center, font=\scriptsize] (out) at (16.2, 3.6) {$1\times 1$\\[2pt]Conv\\[2pt]Head};
  \draw[-{Stealth[scale=0.9]}, very thick, examred] (d1.east) -- (out.west);
  \node[font=\tiny\bfseries\color{examred}, align=center, above=4pt of out] {Output Map\\[1pt]$H \times W \times K$};

  % SKIP CONNECTIONS (Horizontal bridges)
  \draw[skiparrow] (e1.east) -- node[above, font=\scriptsize\bfseries\color{intuitiongreen}] {Skip Connection (Level 1: High-Res Details)} (d1.west);
  \draw[skiparrow] (e2.east) -- node[above, font=\scriptsize\bfseries\color{intuitiongreen}] {Skip Connection (Level 2)} (d2.west);
  \draw[skiparrow] (e3.east) -- node[above, font=\scriptsize\bfseries\color{intuitiongreen}] {Skip Connection (Level 3)} (d3.west);

  % Architectural Legend / Role Callout Box (Positioned below the bottleneck with ample clearance)
  \node[draw=politoNavy!50, fill=politoNavy!5, rounded corners=4pt, line width=0.8pt, font=\scriptsize, align=left, anchor=north, text width=16.2cm, inner sep=6pt] at (7.9, -4.3) {
    \textbf{The Semantic-Resolution Division of Labor in U-Net}:\\
    $\bullet$ \textcolor{politoNavy}{\textbf{Contractive Path (Encoder)}}: Halves spatial grid, doubles channels $\to$ captures global scene semantics (``\textit{What}'').\\
    $\bullet$ \textcolor{deepblue}{\textbf{Expansive Path (Decoder)}}: Doubles spatial grid via transposed convolutions $\to$ restores input dimensions.\\
    $\bullet$ \textcolor{intuitiongreen}{\textbf{Skip Connections}}: Transfer pristine uncompressed spatial edges $\to$ recovers boundary localization (``\textit{Where}'').
  };
\end{tikzpicture}%
}
\caption{The Encoder-Decoder architecture with long skip connections for semantic image segmentation (U-Net archetype). The contractive encoder compresses spatial resolution to capture high-level semantic context, while the expansive decoder upsamples feature maps back to full resolution. Skip connections route fine-grained spatial boundary details across the bottleneck, resolving the trade-off between global semantics and pixel-level localization.}
\label{fig:unet_architecture_schematic}
\end{figure}

\subsection{The empirical paradigm shift: ImageNet benchmark evolution (2011--2015)}
\label{sec:imagenet_benchmark_revolution}

The theoretical principles governing convolutional networks---translation equivariance, receptive field arithmetic, and multi-task loss formulation---found their decisive historical validation through the \textbf{ImageNet Large Scale Visual Recognition Challenge (ILSVRC)}~\cite{deng2009imagenet,russakovsky2015imagenet}.

\paragraph{1. The ImageNet Benchmark Setting.}
Constructed by Fei-Fei Li and colleagues in 2009~\cite{deng2009imagenet}, ImageNet organized over 14 million images mapped to the WordNet lexical hierarchy. The annual ILSVRC competition established the rigorous benchmark:
\begin{itemize}[noitemsep]
  \item \textbf{Data Scale}: $1.28$ million training images, $50,000$ validation images, and $100,000$ test images.
  \item \textbf{Category Granularity}: Exactly $K = 1,000$ fine-grained object categories (including 120 distinct breeds of dogs).
  \item \textbf{Evaluation Metrics}:
    \begin{itemize}
      \item \textbf{Top-1 Error}: The fraction of test images where the model's single highest-probability class prediction is incorrect ($\hat{y}_{\text{pred}} \neq y_{\text{true}}$).
      \item \textbf{Top-5 Error}: The fraction of test images where the true category is not among the model's top 5 most confident class predictions ($\text{True Label} \notin \{c_{(1)}, c_{(2)}, c_{(3)}, c_{(4)}, c_{(5)}\}$). Top-5 error served as the official ranking metric to account for inherent semantic ambiguity in fine-grained labels.
    \end{itemize}
\end{itemize}

\paragraph{2. The Pre-2012 Handcrafted Plateau.}
Prior to 2012, computer vision was dominated by classical multi-stage pipelines:
\begin{enumerate}[leftmargin=*]
  \item Dense handcrafted local feature extraction: Scale-Invariant Feature Transform (SIFT, Lowe 2004)~\cite{lowe2004distinctive} and Histogram of Oriented Gradients (HOG, Dalal \& Triggs 2005)~\cite{dalal2005histograms}.
  \item Vector quantization and aggregation: Bag-of-Visual-Words (BoVW), Spatial Pyramid Matching, and Fisher Vectors (Perronnin et al., 2010).
  \item Supervised linear classification: One-vs-Rest kernel Support Vector Machines (SVMs).
\end{enumerate}
Despite decades of meticulous mathematical design by vision scientists, these handcrafted pipelines stalled at an empirical plateau: the winning 2011 system achieved a \textbf{top-5 error of $26.2\%$}. Many researchers believed that computer vision had encountered an asymptotic ceiling governed by the limits of visual representations.

\paragraph{3. The 2012 AlexNet Watershed.}
In October 2012, Alex Krizhevsky, Ilya Sutskever, and Geoffrey Hinton entered \textbf{AlexNet}~\cite{krizhevsky2012imagenet} into ILSVRC 2012, achieving a top-5 error of \textbf{16.4\%} (and $15.3\%$ with model ensembling).
This marked an unprecedented \textbf{9.8 percentage point drop} in error, cutting relative classification error by nearly \textbf{40\%} in a single year. AlexNet demonstrated that deep neural networks trained with backpropagation could vastly outperform decades of human feature engineering.

AlexNet succeeded by combining four technological pillars:
\begin{enumerate}[noitemsep]
  \item \textbf{Architectural Depth}: 8 learned layers (5 convolutional layers followed by 3 dense layers, encompassing 61 million parameters).
  \item \textbf{Non-Saturating Activations}: Replacing sigmoidal functions with Rectified Linear Units ($\text{ReLU}(z) = \max(0, z)$), accelerating gradient propagation and speeding up training convergence by $6\times$.
  \item \textbf{Hardware Acceleration}: Parallelizing custom CUDA convolutional kernels across two NVIDIA GeForce GTX 580 GPUs (3 GB VRAM each).
  \item \textbf{Overfitting Regularization}: Applying Dropout ($p=0.5$) in the dense layers and massive data augmentation (random spatial crops, horizontal reflections, and PCA color jittering).
\end{enumerate}

\paragraph{4. The ILSVRC Architectural Progression (2011--2015).}
Table~\ref{tab:imagenet_historical_evolution} chronicles the landmark architectural progression on the ImageNet benchmark, culminating in deep residual learning surpassing human visual classification performance.

\begin{table}[H]
\centering
\footnotesize
\setlength{\tabcolsep}{4pt}
\caption{Historical evolution of top-5 classification error on the ImageNet (ILSVRC) 1,000-class benchmark, documenting the transition from pre-2012 handcrafted feature pipelines to ultra-deep residual networks.}
\label{tab:imagenet_historical_evolution}
\vspace{4pt}
\begin{tabularx}{\textwidth}{>{\raggedright\arraybackslash}p{2.6cm} c c c c >{\raggedright\arraybackslash}X}
\toprule
\textbf{Model / Paradigm} & \textbf{Year} & \textbf{Depth} & \textbf{Params} & \textbf{Top-5 Err.} & \textbf{Key Innovations \& Architectural Milestones} \\
\midrule
\textbf{Handcrafted Baseline} & 2011 & Shallow & $\sim 10^7$ & $26.2\%$ & Dense SIFT, Fisher Vectors, Bag-of-Visual-Words (BoVW), Spatial Pyramid Matching, linear/kernel SVMs. Severe empirical ceiling. \\
\addlinespace
\textbf{AlexNet} & 2012 & 8 & $61\,\text{M}$ & $16.4\%$ & Deep CNN watershed: cut relative error by $\sim 40\%$. ReLU activations, Dropout ($p=0.5$), 2$\times$ GTX 580 GPUs, overlapping max pooling~\cite{krizhevsky2012imagenet}. \\
\addlinespace
\textbf{ZFNet} & 2013 & 8 & $60\,\text{M}$ & $11.7\%$ & Deconvolutional feature visualization; reduced Layer 1 filter size from $11\times 11$ ($S=4$) to $7\times 7$ ($S=2$) to preserve high-frequency features~\cite{zeiler2014visualizing}. \\
\addlinespace
\textbf{VGG-16} & 2014 & 16 & $138\,\text{M}$ & $7.3\%$ & Homogeneous $3\times 3$ kernel stacking (VGG principle); proves depth and simplicity trump large filters; $45\%$ parameter savings vs. $7\times 7$~\cite{simonyan2014very}. \\
\addlinespace
\textbf{GoogLeNet} & 2014 & 22 & $6.8\,\text{M}$ & $6.7\%$ & Multi-scale Inception modules ($1\times 1, 3\times 3, 5\times 5$); $1\times 1$ conv dimensionality reduction; Global Average Pooling; $20\times$ fewer params than VGG~\cite{szegedy2015going}. \\
\addlinespace
\textbf{ResNet-152} & 2015 & 152 & $60.2\,\text{M}$ & $\mathbf{3.57\%}$ & Residual skip connections $\mathbf{y} = \mathcal{F}(\mathbf{x}) + \mathbf{x}$; identity shortcuts solve the gradient degradation problem; enabled 152-layer optimization~\cite{he2016deep}. \\
\midrule
\textbf{Human Baseline} & 2015 & Cortex & $\sim 10^{14}$ syn. & $\approx \mathbf{5.1\%}$ & Andrej Karpathy human benchmark study (1,500 test images). ResNet-152 became the first model in history to surpass human performance. \\
\bottomrule
\end{tabularx}
\end{table}

\paragraph{5. Conceptual Takeaways of the Deep Revolution.}
The historical trajectory in Table~\ref{tab:imagenet_historical_evolution} yields three profound theoretical conclusions:
\begin{enumerate}[leftmargin=*]
  \item \textbf{End-to-End Representation Learning}: The superiority of learned representations over engineered descriptors demonstrated that optimizing feature extractors directly on the downstream objective is vastly more powerful than separating feature design from classifier optimization.
  \item \textbf{Depth and Inductive Bias as Regularizers}: As evidenced by VGG, GoogLeNet, and ResNet, network performance scales directly with architectural depth ($8 \to 16 \to 22 \to 152$ layers). Crucially, depth is only trainable when paired with principled architectural inductive biases: small kernel stacking, multi-scale bottleneck projections, and residual identity mappings.
  \item \textbf{Generalization Beyond ImageNet}: Representations learned on ImageNet serve as universal visual representations. A convolutional backbone pre-trained on ImageNet can be fine-tuned or transferred to object detection (Faster R-CNN), semantic segmentation (FCN/U-Net), medical imaging, and autonomous driving, establishing pre-training and transfer learning as the foundational paradigm of modern machine learning.
\end{enumerate}
